\documentclass[11pt]{article}

\usepackage[
    paperwidth=176mm,
    paperheight=250mm,
    top=18mm,
    bottom=22mm,
    left=18mm,
    right=18mm,
    footskip=10mm,
]{geometry}

\usepackage{graphicx}
\usepackage{xcolor}
\usepackage{fancyhdr}
\usepackage{parskip}
\usepackage{ragged2e}
\usepackage{fontspec}
\usepackage{lastpage}

% Use Arial — present on every Windows install, looked up via fontconfig.
% fontspec will fall back to the system default if Arial is unavailable.
\IfFontExistsTF{Arial}{%
    \setmainfont{Arial}%
}{%
    % Latin Modern is bundled with every TeX distribution; safe fallback.
    \setmainfont{Latin Modern Roman}%
}

\pagestyle{fancy}
\fancyhf{}
\renewcommand{\headrulewidth}{0pt}
\renewcommand{\footrulewidth}{0.4pt}

\newcommand{\theposttitle}{}
\newcommand{\thepostdate}{}

\fancyfoot[L]{\color{gray!70}\footnotesize\itshape\theposttitle}
\fancyfoot[C]{\color{gray!70}\footnotesize\thepage\,/\,\pageref{LastPage}}
\fancyfoot[R]{\color{gray!70}\footnotesize\thepostdate}

\setlength{\parindent}{0pt}

\begin{document}
\renewcommand{\theposttitle}{I.D.E.A. Newsletter 6: Keystone Habits \& Growth Zones}
\renewcommand{\thepostdate}{2022-07-01}
\vspace*{\fill}
\begin{center}
{\Huge\bfseries INTERESTING}\\[6mm]
\end{center}
\begin{center}
\includegraphics[width=\textwidth, height=0.55\textheight, keepaspectratio]{C:/GitHub/VisualIdeas/data/images/2d8adbbfe6610df97852fafcdb869054.png}
\end{center}
\vspace{2mm}
\begin{center}
\begin{minipage}{0.92\textwidth}
\small
Recently on a visit to my kid's doctor appointment, I rediscovered this toy I loved. My favorite way of playing with these was grabbing the bottom-most block and swooping it around from one end of the toy to the other. It was always so satisfying feeling it bring the rest of the blocks along for the ride.This made me think about habit stacking from the book, "The Power of Habit". This was a first draft of the idea and I hope to refine it in a future post. 😊
\end{minipage}
\end{center}
\vspace*{\fill}
\newpage
\renewcommand{\theposttitle}{I.D.E.A. Newsletter 6: Keystone Habits \& Growth Zones}
\renewcommand{\thepostdate}{2022-07-01}
\vspace*{\fill}
\begin{center}
{\Huge\bfseries ENCHANTING}\\[6mm]
\end{center}
\begin{center}
\includegraphics[width=\textwidth, height=0.55\textheight, keepaspectratio]{C:/GitHub/VisualIdeas/data/images/a6a573c20f5f02d2dc7a469286ac854d.png}
\end{center}
\vspace{2mm}
\begin{center}
\begin{minipage}{0.92\textwidth}
\small
This is more or less my creative process, sprinkled in with experience and moments of serendipity. A lot of ideas come to me when I'm driving or in the shower since these are my quiet times.
\end{minipage}
\end{center}
\vspace*{\fill}
\newpage
\renewcommand{\theposttitle}{I.D.E.A. Newsletter 6: Keystone Habits \& Growth Zones}
\renewcommand{\thepostdate}{2022-07-01}
\vspace*{\fill}
\begin{center}
{\Huge\bfseries ANALOGY}\\[6mm]
\end{center}
\begin{center}
\includegraphics[width=\textwidth, height=0.55\textheight, keepaspectratio]{C:/GitHub/VisualIdeas/data/images/0d32e79f100487855f985885e5b4fed8.png}
\end{center}
\vspace{2mm}
\begin{center}
\begin{minipage}{0.92\textwidth}
\small
I have a LEGO ship in a bottle sitting in my office. There's something comforting knowing the ship is safe and protected from the outside world so it got me thinking how it's pretty much equal to a comfort zone.
\end{minipage}
\end{center}
\vspace*{\fill}
\newpage
\renewcommand{\theposttitle}{The Beginner's Mind and Finding Clarity}
\renewcommand{\thepostdate}{2022-07-08}
\vspace*{\fill}
\begin{center}
{\Huge\bfseries INTERESTING}\\[6mm]
\end{center}
\begin{center}
\includegraphics[width=\textwidth, height=0.55\textheight, keepaspectratio]{C:/GitHub/VisualIdeas/data/images/88fed55bf7c04e9b80d33e30c706f162.png}
\end{center}
\vspace{2mm}
\begin{center}
\begin{minipage}{0.92\textwidth}
\small
The secret sauce to solving difficult problems is combining the efficiency of the expert with the openness of a beginners mind.There’s tremendous value in specialization and expertise.But there's also a hidden trap.To become an expert, you have to zoom in and learn deeply about a singular domain. The more you do something, the more efficient you get. But as you specialize and your efficiency improves, your attention gets narrower and you lose the ability to think laterally.Said another way, when you become an expert you sacrifice creativity for efficiency. The visual above celebrates the beginner's mind, or \textit{shoshin }(初心) as its known in Japanese Zen Buddhism. In the context of Zen, a beginner’s mind is like when you eat something for the very first time. Imagine that first bite. Its taste. Its texture. Its aroma...fully immersed in each moment of that experience. Having a beginner's mind is about: • snapping yourself out of auto-pilot mode• sustaining curiosity• and checking your assumptions about what's possible.

\textbf{DIVE DEEPER: }My favorite real world example of this is the the case of the scientist who discovered H. Pylori.
\end{minipage}
\end{center}
\vspace*{\fill}
\newpage
\renewcommand{\theposttitle}{The Beginner's Mind and Finding Clarity}
\renewcommand{\thepostdate}{2022-07-08}
\vspace*{\fill}
\begin{center}
{\Huge\bfseries ENCHANTING}\\[6mm]
\end{center}
\begin{center}
\includegraphics[width=\textwidth, height=0.55\textheight, keepaspectratio]{C:/GitHub/VisualIdeas/data/images/06da478c5e162eb9e66926399b395f6d.png}
\end{center}
\vspace{2mm}
\begin{center}
\begin{minipage}{0.92\textwidth}
\small
Had so much fun drawing this, especially figuring out how to create the illusion of shaking up the globe.

When it’s manageable, anxiety and tension can help motivate you towards dreams and aspirations.But if you feel a paralyzing chaos welling up daily, I’d invite you to start a journaling practice to calm the mind and settle the soul.Journaling can be more than an inventory of the events of the day. It can be a method to reflect on what's going on internally and become more consciously aware of what you're feeling.Any kind of reflection is helpful, but I’m a proponent of journaling by hand because it helps you slow down and take note of what you're experiencing.
\end{minipage}
\end{center}
\vspace*{\fill}
\newpage
\renewcommand{\theposttitle}{The Beginner's Mind and Finding Clarity}
\renewcommand{\thepostdate}{2022-07-08}
\vspace*{\fill}
\begin{center}
{\Huge\bfseries ANALOGY}\\[6mm]
\end{center}
\begin{center}
\includegraphics[width=\textwidth, height=0.55\textheight, keepaspectratio]{C:/GitHub/VisualIdeas/data/images/cabfc77005ddd09f821a3eeab3b3b602.png}
\end{center}
\vspace{2mm}
\begin{center}
\begin{minipage}{0.92\textwidth}
\small
As a teacher, an analogy is the most effective way to communicate foreign concepts. I came up with this one as I backed out of the driveway thinking about how to explain why my students should write to clarify their thinking about what they wanted to achieve for the year. I'm flattered to have been DM'd several times over the last couple months asking about my process of coming up with visuals. It may sound weird, but writing is an essential component of visualization. Without the clarity found through exploring an idea, a visual has no soul. That said, I still feel a deep anxiety with every visual before posting, wondering "Is This Anything?"[affiliate link] much like Jerry Seinfeld thinks before a joke.
\end{minipage}
\end{center}
\vspace*{\fill}
\newpage
\renewcommand{\theposttitle}{The Cycle of Growth and Stoic Philosophy}
\renewcommand{\thepostdate}{2022-07-15}
\vspace*{\fill}
\begin{center}
{\Huge\bfseries INTERESTING}\\[6mm]
\end{center}
\begin{center}
\includegraphics[width=\textwidth, height=0.55\textheight, keepaspectratio]{C:/GitHub/VisualIdeas/data/images/f635d99b631a1ba5ff075e5cce550266.png}
\end{center}
\vspace{2mm}
\begin{center}
\begin{minipage}{0.92\textwidth}
\small
This image made it to the Interesting slot because it's the first image I posted where I received a ton of feedback on the analogy itself (which I love actually because it forces me to improve my thinking and iterate the drawing).

For me, the key idea is when you're working on something alone, the low-hanging fruit is the easiest to grab.

But if you want a sweeter reward, getting there in a team is easier, even if the fruit (🥁) of that labor may need to be shared.

Honestly though...working together with a community toward a goal is more fun and better for the spirit.

I imagine the characters in this picture cracking up as they try not to topple over one another. 😂

\textbf{Artist's Note:} Iteration for Clarity• Include all three of the characters grabbing low-hanging fruit with three golden apples at the top. • When they're grabbing the fruit on the right, "Bluey" will have a basket full of 2 golden apples as he holds the third.
\end{minipage}
\end{center}
\vspace*{\fill}
\newpage
\renewcommand{\theposttitle}{The Cycle of Growth and Stoic Philosophy}
\renewcommand{\thepostdate}{2022-07-15}
\vspace*{\fill}
\begin{center}
{\Huge\bfseries ENCHANTING}\\[6mm]
\end{center}
\begin{center}
\includegraphics[width=\textwidth, height=0.55\textheight, keepaspectratio]{C:/GitHub/VisualIdeas/data/images/9e1499116ca07159c80fe7ac38143f71.png}
\end{center}
\vspace{2mm}
\begin{center}
\begin{minipage}{0.92\textwidth}
\small
99\% of people who embrace a growth mindset fall into the trap thinking there's a singular moment when they'll be...well...fully grown (AKA good at the thing).

But this is a mistake.

\textbf{Growth is a cycle,} not a destination.

You're playing an infinite game. Enjoy the swim.

\textbf{Artist's Note: }I saw this weird looking circle in one of my kid's children's books. It got me thinking about this birds eye shot of an ocean with this image of a guy swimming around and around, lap by lap, hour by hour, with no goal but to keep swimming. If I drew this in my 20s, I'd probably draw the guy swimming the whole cycle through. But being in my 40s now, I see the value of rest and how it contributes to a more solid lasting foundation of growth.
\end{minipage}
\end{center}
\vspace*{\fill}
\newpage
\renewcommand{\theposttitle}{The Cycle of Growth and Stoic Philosophy}
\renewcommand{\thepostdate}{2022-07-15}
\vspace*{\fill}
\begin{center}
{\Huge\bfseries ANALOGY}\\[6mm]
\end{center}
\begin{center}
\includegraphics[width=\textwidth, height=0.55\textheight, keepaspectratio]{C:/GitHub/VisualIdeas/data/images/addeba8703ded30b901e0d18cf244283.png}
\end{center}
\vspace{2mm}
\begin{center}
\begin{minipage}{0.92\textwidth}
\small
Inspired by one of my favorite ideas of Stoicism illustrated most eloquently in this line:

“The impediment to action advances action. What stands in the way becomes the way.”

How you perceive the obstacles in your life defines your path forward.

The obstacle isn't a bug. It's the feature.

\textbf{Artist's Note:} This idea came to me when I was looking at my staircase. When I looked straight at them, the tiling made it look as if one stair blended into the next! It got me thinking...this is actually quite dangerous...😅
\end{minipage}
\end{center}
\vspace*{\fill}
\newpage
\renewcommand{\theposttitle}{I.D.E.A. Newsletter 9: Storytelling and Comfort Zones}
\renewcommand{\thepostdate}{2022-07-22}
\vspace*{\fill}
\begin{center}
{\Huge\bfseries INTERESTING}\\[6mm]
\end{center}
\begin{center}
\includegraphics[width=\textwidth, height=0.55\textheight, keepaspectratio]{C:/GitHub/VisualIdeas/data/images/4d0110068e4ee382f4b22ff2562e5d05.jpeg}
\end{center}
\vspace{2mm}
\begin{center}
\begin{minipage}{0.92\textwidth}
\small
The highlight of my week was when my friend @EvansNifty posted this visual demonstrating my philosophy on storytelling.

If you want your content to stand out, involve the audience as active participants of your work.Telling a good story is as if you're drawing a circle for your audience.

But instead of forming the circle completely by yourself, you want to leave it slightly incomplete.

That small gap at the end is for the audience to complete themselves.

Don't give the audience the answer (4). Engage them with the question (2+2).
\end{minipage}
\end{center}
\vspace*{\fill}
\newpage
\renewcommand{\theposttitle}{I.D.E.A. Newsletter 9: Storytelling and Comfort Zones}
\renewcommand{\thepostdate}{2022-07-22}
\vspace*{\fill}
\begin{center}
{\Huge\bfseries ENCHANTING}\\[6mm]
\end{center}
\begin{center}
\includegraphics[width=\textwidth, height=0.24\textheight, keepaspectratio]{C:/GitHub/VisualIdeas/data/images/998ae232699231de89393e811985d724.png}
\end{center}
\vspace{5mm}
\begin{center}

\includegraphics[width=\textwidth, height=0.24\textheight, keepaspectratio]{C:/GitHub/VisualIdeas/data/images/1707e8138ed61b9220cdd548236b3de6.png}
\end{center}
\vspace{2mm}
\begin{center}
\begin{minipage}{0.92\textwidth}
\small
The very last lines of "The Practice" by Seth Godin reads:

"The magic is that there is no magic. Start where you're at. Don't Stop."

The top of the Maslow's Hierarchy of Needs pyramid should be changed from "Self-Actualization" to "Continuous Progress."

There's a reason why we begin most movies with our hero in an Ordinary World.

This is the comfort zone.

Soon after our hero's introduced, there's a Call to Action, followed swiftly by the hero's Refusal of the Call.

Why?

Because few people intentionally step into the unknown and into discomfort.

They often have to be pushed into it.

The tragic truth of growth is that it's something we recognize only in hindsight and not as it's happening.

But growth requires us to step into the fray and trust we'll get through to the other side stronger for it.

A flower only grows when it reaches for the sun.
\end{minipage}
\end{center}
\vspace*{\fill}
\newpage
\renewcommand{\theposttitle}{I.D.E.A. Newsletter 9: Storytelling and Comfort Zones}
\renewcommand{\thepostdate}{2022-07-22}
\vspace*{\fill}
\begin{center}
{\Huge\bfseries ANALOGY}\\[6mm]
\end{center}
\begin{center}
\includegraphics[width=\textwidth, height=0.55\textheight, keepaspectratio]{C:/GitHub/VisualIdeas/data/images/0cfade5a7562e079497b077a3a497a16.png}
\end{center}
\vspace{2mm}
\begin{center}
\begin{minipage}{0.92\textwidth}
\small
This metaphor tries to tap into multiple different areas and meanings of the words "consuming" and "creating." On face value, consuming is taking from the world, whereas creating is giving back.But also the earth is a finite and fragile resource. We need to drink from it to survive. If we consume without giving back, it will run dry.

Creating a piece of writing, a piece of art, or even a business to serve customers and solve their problems is a generous act.

Be generous. Create something and share it with the world. Give more than you take.
\end{minipage}
\end{center}
\vspace*{\fill}
\newpage
\renewcommand{\theposttitle}{Read More, Build Momentum}
\renewcommand{\thepostdate}{2022-08-05}
\vspace*{\fill}
\begin{center}
{\Huge\bfseries INTERESTING}\\[6mm]
\end{center}
\begin{center}
\includegraphics[width=\textwidth, height=0.55\textheight, keepaspectratio]{C:/GitHub/VisualIdeas/data/images/6f1c16af873e6dd57a95c890a8c66298.png}
\end{center}
\vspace{2mm}
\begin{center}
\begin{minipage}{0.92\textwidth}
\small
A simple quote from Naval Ravikant, but one of my favorites. • This inspired a lot of interesting comments, especially on IG.• Most mentioned answers: gym, reading, and teaching!

\textbf{Artist's Note: }When you create visuals, everything and anything can be a source of inspiration. I came up with this idea playing with a PJ (🥁) Mask toy at the house. 😂 The light blue cage is part of the Headquarter set!
\end{minipage}
\end{center}
\vspace*{\fill}
\newpage
\renewcommand{\theposttitle}{Read More, Build Momentum}
\renewcommand{\thepostdate}{2022-08-05}
\vspace*{\fill}
\begin{center}
{\Huge\bfseries DESIGN}\\[6mm]
\end{center}
\begin{center}
\includegraphics[width=\textwidth, height=0.55\textheight, keepaspectratio]{C:/GitHub/VisualIdeas/data/images/80fbbbcf0892c42fdf44f20d973d94f1.png}
\end{center}
\vspace{2mm}
\begin{center}
\begin{minipage}{0.92\textwidth}
\small
\textbf{Shoot me a reply with design/drawing questions and I'll add it to the list.}
\end{minipage}
\end{center}
\vspace*{\fill}
\newpage
\renewcommand{\theposttitle}{Read More, Build Momentum}
\renewcommand{\thepostdate}{2022-08-05}
\vspace*{\fill}
\begin{center}
{\Huge\bfseries ENCHANTING}\\[6mm]
\end{center}
\begin{center}
\includegraphics[width=\textwidth, height=0.24\textheight, keepaspectratio]{C:/GitHub/VisualIdeas/data/images/242257cb6018c6e11fd1192ae5c9660e.png}
\end{center}
\vspace{2mm}
\begin{center}
\begin{minipage}{0.92\textwidth}
\small
I can't go to the library without checking out at least half the library so I wanted to celebrate the power of reading with this visual.It's true that some skills are impossible to learn by reading alone...like riding a bike or swimming!

But some skills are slow to learn by action alone.

You don't know what you don't know so read more!

\textbf{Artist's Note: }Design is all about problem solving.

If you look at the original post on Twitter, it reads:"READ MORE BOOKS TO GROW" (if read top =\textgreater{} bottom) but some people read it as:"GROW TO BOOKS MORE READ" (bottom =\textgreater{} top). 😅

Why?

The problem is with the visual design 👇🏻
\end{minipage}
\end{center}
\vspace{5mm}
\begin{center}
\includegraphics[width=\textwidth, height=0.24\textheight, keepaspectratio]{C:/GitHub/VisualIdeas/data/images/deb1d48630b8238064718346cecd9922.png}
\end{center}
\vspace{2mm}
\begin{center}
\begin{minipage}{0.92\textwidth}
\small
Solution?

Iterate. Come up with 5 words that make sense read forward AND backward:"Step By Step, More Reading," which can also be read:"Reading More, Step by Step."

Not perfect, but better! 💪
\end{minipage}
\end{center}
\vspace*{\fill}
\newpage
\renewcommand{\theposttitle}{Read More, Build Momentum}
\renewcommand{\thepostdate}{2022-08-05}
\vspace*{\fill}
\begin{center}
{\Huge\bfseries ANALOGY}\\[6mm]
\end{center}
\begin{center}
\includegraphics[width=\textwidth, height=0.55\textheight, keepaspectratio]{C:/GitHub/VisualIdeas/data/images/193c51162fc5aa3dd7600eb8cc388221.png}
\end{center}
\vspace{2mm}
\begin{center}
\begin{minipage}{0.92\textwidth}
\small
My depiction of the "flow state."\textbf{Artist's Note:} I was playing with my kids, showing them the slinky-stair trick when this idea came to me. This is the 2nd iteration. You can see the first on my IG.

If you want to make your own visuals, give yourself permission to be imperfect. Know you can make it better with v.2, v.3, and so on till you get the idea right.
\end{minipage}
\end{center}
\vspace*{\fill}
\newpage
\renewcommand{\theposttitle}{Online Mindsets and Pruning Your Commitments}
\renewcommand{\thepostdate}{2022-08-23}
\vspace*{\fill}
\begin{center}
{\Huge\bfseries INTERESTING}\\[6mm]
\end{center}
\begin{center}
\begin{minipage}{0.92\textwidth}
\small
\textbf{3 Social Media Mindsets, but only one makes an impact.}
\end{minipage}
\end{center}
\vspace{4mm}
\begin{center}
\includegraphics[width=\textwidth, height=0.17\textheight, keepaspectratio]{C:/GitHub/VisualIdeas/data/images/ad5dbab945684a923b5ad81b2d689250.png}
\end{center}
\vspace{5mm}
\begin{center}
\includegraphics[width=\textwidth, height=0.17\textheight, keepaspectratio]{C:/GitHub/VisualIdeas/data/images/43057a6c2ce0a230a211e05e2c23d192.png}
\end{center}
\vspace{5mm}
\begin{center}

\includegraphics[width=\textwidth, height=0.17\textheight, keepaspectratio]{C:/GitHub/VisualIdeas/data/images/1422087ca9e17848a53a3fe76546316c.png}
\end{center}
\vspace*{\fill}
\newpage
\renewcommand{\theposttitle}{Online Mindsets and Pruning Your Commitments}
\renewcommand{\thepostdate}{2022-08-23}
\vspace*{\fill}
\begin{center}
{\Huge\bfseries DESIGN}\\[6mm]
\end{center}
\begin{center}
\includegraphics[width=\textwidth, height=0.55\textheight, keepaspectratio]{C:/GitHub/VisualIdeas/data/images/322cae304e7e7af54967c40e2adb0afc.png}
\end{center}
\vspace{2mm}
\begin{center}
\begin{minipage}{0.92\textwidth}
\small
A simple but easy to apply lesson in drawing the same object in perspective.
\end{minipage}
\end{center}
\vspace*{\fill}
\newpage
\renewcommand{\theposttitle}{Online Mindsets and Pruning Your Commitments}
\renewcommand{\thepostdate}{2022-08-23}
\vspace*{\fill}
\begin{center}
{\Huge\bfseries ENCHANTING}\\[6mm]
\end{center}
\begin{center}
\includegraphics[width=\textwidth, height=0.55\textheight, keepaspectratio]{C:/GitHub/VisualIdeas/data/images/a025e4684749b1dad60f202433b69c0b.png}
\end{center}
\vspace{2mm}
\begin{center}
\begin{minipage}{0.92\textwidth}
\small
An ode to the Sistine Chapel "Creation of Adam."
\end{minipage}
\end{center}
\vspace*{\fill}
\newpage
\renewcommand{\theposttitle}{Online Mindsets and Pruning Your Commitments}
\renewcommand{\thepostdate}{2022-08-23}
\vspace*{\fill}
\begin{center}
{\Huge\bfseries ANALOGY}\\[6mm]
\end{center}
\begin{center}
\includegraphics[width=\textwidth, height=0.55\textheight, keepaspectratio]{C:/GitHub/VisualIdeas/data/images/293b49ad4f52a21f0ad68e76c8ec2df1.png}
\end{center}
\vspace{2mm}
\begin{center}
\begin{minipage}{0.92\textwidth}
\small
The word priority initially meant 1 thing above all others. Having multiple priorities is the same as having no priority.

I might be mixing my metaphors here, but if you plant 10 seeds and ration a cup of water over all 10 seeds, none of them sprout.But if you plant a single seed and use the entire cup to nourish the soil surrounding the seed, it will bear the most delicious of fruits.

When you cut off the excess, you cultivate the space to grow.
\end{minipage}
\end{center}
\vspace*{\fill}
\newpage
\renewcommand{\theposttitle}{Risks, Throbbers, and the Ego!}
\renewcommand{\thepostdate}{2022-09-06}
\vspace*{\fill}
\begin{center}
{\Huge\bfseries INTERESTING}\\[6mm]
\end{center}
\begin{center}
\includegraphics[width=\textwidth, height=0.55\textheight, keepaspectratio]{C:/GitHub/VisualIdeas/data/images/fec72b03a8a2590e95c4dcb486fd5a09.png}
\end{center}
\vspace{2mm}
\begin{center}
\begin{minipage}{0.92\textwidth}
\small
I made this inspired by Justin Welsh's Saturday Solopreneur newsletter, which was a great reminder to take small bets.

Said another way, take the big thing you're trying to achieve and do a 1\% version of it. (idea from Derek Sivers's book, "Anything You Want.")

If there's anything in your life you've always wanted to do, excavate and execute the 1\% version of it. If there's a big gamble you've been considering, is there a way to make a smaller gamble to test your assumptions?
\end{minipage}
\end{center}
\vspace*{\fill}
\newpage
\renewcommand{\theposttitle}{Risks, Throbbers, and the Ego!}
\renewcommand{\thepostdate}{2022-09-06}
\vspace*{\fill}
\begin{center}
{\Huge\bfseries DESIGN}\\[6mm]
\end{center}
\begin{center}
\includegraphics[width=\textwidth, height=0.55\textheight, keepaspectratio]{C:/GitHub/VisualIdeas/data/images/fd1bd3dda9cd1c91bf58e98e38d843f3.png}
\end{center}
\vspace{2mm}
\begin{center}
\begin{minipage}{0.92\textwidth}
\small
These loading icons are called 'throbbers' (cue high school student snickering). Apple decided to shift from a wristwatch frozen at 9 o'clock (where you're reminded the effect of how you're losing time) to a spinning pinwheel (aka beachball, which was a bit more playful). Of course, we've adapted to the interface and now call it "the spinning wheel of death"...but hey, still a brilliant move in UX design..
\end{minipage}
\end{center}
\vspace*{\fill}
\newpage
\renewcommand{\theposttitle}{Risks, Throbbers, and the Ego!}
\renewcommand{\thepostdate}{2022-09-06}
\vspace*{\fill}
\begin{center}
{\Huge\bfseries ENCHANTING}\\[6mm]
\end{center}
\begin{center}
\includegraphics[width=\textwidth, height=0.55\textheight, keepaspectratio]{C:/GitHub/VisualIdeas/data/images/6292630fe00d49aa945edcac0f3c32e1.png}
\end{center}
\vspace{2mm}
\begin{center}
\begin{minipage}{0.92\textwidth}
\small
A makeover of one of my first visuals on Twitter.

Goes to show the value of putting your work out there and iterating on it. 😃
\end{minipage}
\end{center}
\vspace*{\fill}
\newpage
\renewcommand{\theposttitle}{Risks, Throbbers, and the Ego!}
\renewcommand{\thepostdate}{2022-09-06}
\vspace*{\fill}
\begin{center}
{\Huge\bfseries ANALOGY}\\[6mm]
\end{center}
\begin{center}
\includegraphics[width=\textwidth, height=0.55\textheight, keepaspectratio]{C:/GitHub/VisualIdeas/data/images/cbfc3276460d220fa5b4970e996b729f.png}
\end{center}
\vspace{2mm}
\begin{center}
\begin{minipage}{0.92\textwidth}
\small
An important reminder for all of us: Don't let it all get to your head. 😅
\end{minipage}
\end{center}
\vspace*{\fill}
\newpage
\renewcommand{\theposttitle}{Managing vs Leading \& Life Events vs Storytelling}
\renewcommand{\thepostdate}{2022-09-20}
\vspace*{\fill}
\begin{center}
{\Huge\bfseries INTERESTING}\\[6mm]
\end{center}
\begin{center}
\includegraphics[width=\textwidth, height=0.55\textheight, keepaspectratio]{C:/GitHub/VisualIdeas/data/images/7425d378105c3426c6db0651bea74f45.png}
\end{center}
\vspace{2mm}
\begin{center}
\begin{minipage}{0.92\textwidth}
\small
There's a bunch of advice on time management but little on energy management.

To some willpower is grit or the ability to persevere.

To others it's the ability to resist.

For me, it's the discipline to stick to my commitments.

I found if I front-load my day by following through on commitments, I have less anxiety the rest of the day.

I used to be the opposite, but recently I've recognized I'm depleted after dinner and got very little left in the tank.

How do you manage your energy?
\end{minipage}
\end{center}
\vspace*{\fill}
\newpage
\renewcommand{\theposttitle}{Managing vs Leading \& Life Events vs Storytelling}
\renewcommand{\thepostdate}{2022-09-20}
\vspace*{\fill}
\begin{center}
{\Huge\bfseries DESIGN}\\[6mm]
\end{center}
\begin{center}
\includegraphics[width=\textwidth, height=0.55\textheight, keepaspectratio]{C:/GitHub/VisualIdeas/data/images/7f0787a6fc83b89c0a041de4f5af15f6.png}
\end{center}
\vspace{2mm}
\begin{center}
\begin{minipage}{0.92\textwidth}
\small
A simple but easy to apply lesson in drawing the same object in perspective. These simple guidelines alone are helpful, but when they're combined with other guidelines can become powerful.
\end{minipage}
\end{center}
\vspace*{\fill}
\newpage
\renewcommand{\theposttitle}{Managing vs Leading \& Life Events vs Storytelling}
\renewcommand{\thepostdate}{2022-09-20}
\vspace*{\fill}
\begin{center}
{\Huge\bfseries ENCHANTING}\\[6mm]
\end{center}
\begin{center}
\includegraphics[width=\textwidth, height=0.55\textheight, keepaspectratio]{C:/GitHub/VisualIdeas/data/images/32e95463bf3b20ade6acaa8f67a9150d.png}
\end{center}
\vspace{2mm}
\begin{center}
\begin{minipage}{0.92\textwidth}
\small
We may think we don't have stories worth sharing, but that's just because we haven't connected the dots yet.

“Make it your mission to find, see, remember, and identify stories, and you will begin to see your life in a new and more compelling light.”

There're lot of books on storytelling, but few have as many practical insights into the artform as Storyworthy.

It's made me reconsider what I should look for in an engaging story, which has proven invaluable in my teaching practice.
\end{minipage}
\end{center}
\vspace*{\fill}
\newpage
\renewcommand{\theposttitle}{Managing vs Leading \& Life Events vs Storytelling}
\renewcommand{\thepostdate}{2022-09-20}
\vspace*{\fill}
\begin{center}
{\Huge\bfseries ANALOGY}\\[6mm]
\end{center}
\begin{center}
\includegraphics[width=\textwidth, height=0.55\textheight, keepaspectratio]{C:/GitHub/VisualIdeas/data/images/a24338966ccda6ac3a175008013a9b72.png}
\end{center}
\vspace{2mm}
\begin{center}
\begin{minipage}{0.92\textwidth}
\small
I am fascinated with how words and labels shift the way we experience the world.

Visuals can help clarify and distill subtle differences in words. Often management and leadership are words used interchangeably but there is a difference.

Both are skills that can be learned.

Both are essential in an organization.

But you should know what you're trying to do and take care to recognize when you should lead and when you need to manage. This was one of my first visuals I posted on Twitter and you can see in my process video, how messy the discovery process of the design and colors were. 😅
\end{minipage}
\end{center}
\vspace*{\fill}
\newpage
\renewcommand{\theposttitle}{Cycle of Growth \& Intentionality}
\renewcommand{\thepostdate}{2022-10-04}
\vspace*{\fill}
\begin{center}
{\Huge\bfseries INTERESTING}\\[6mm]
\end{center}
\begin{center}
\includegraphics[width=\textwidth, height=0.55\textheight, keepaspectratio]{C:/GitHub/VisualIdeas/data/images/d481ab6de052f99f43785095ec29023f.png}
\end{center}
\vspace{2mm}
\begin{center}
\begin{minipage}{0.92\textwidth}
\small
"Intelligence is traditionally viewed as the ability to think and learn.

Yet in a turbulent world, there's another set of cognitive skills that might matter more: the ability to rethink and unlearn."

-Adam Grant, Think Again: The Power of Knowing What You Don't Know
\end{minipage}
\end{center}
\vspace*{\fill}
\newpage
\renewcommand{\theposttitle}{Cycle of Growth \& Intentionality}
\renewcommand{\thepostdate}{2022-10-04}
\vspace*{\fill}
\begin{center}
{\Huge\bfseries DESIGN}\\[6mm]
\end{center}
\begin{center}
\includegraphics[width=\textwidth, height=0.55\textheight, keepaspectratio]{C:/GitHub/VisualIdeas/data/images/7e3df53c4baf97284191ba9cddd86cb9.png}
\end{center}
\vspace{2mm}
\begin{center}
\begin{minipage}{0.92\textwidth}
\small
The key is to keep the diagonal lines (the red lines) parallel to each other and you'll get perfect perspective when drawing anything that has a pathway and gaps: • railroads, • steps • sidewalks, • bridges, etc. Erase the diagonals at the end and you're good to go. Notice the walls I made are parallel lines going straight up.
\end{minipage}
\end{center}
\vspace*{\fill}
\newpage
\renewcommand{\theposttitle}{Cycle of Growth \& Intentionality}
\renewcommand{\thepostdate}{2022-10-04}
\vspace*{\fill}
\begin{center}
{\Huge\bfseries ENCHANTING}\\[6mm]
\end{center}
\begin{center}
\includegraphics[width=\textwidth, height=0.55\textheight, keepaspectratio]{C:/GitHub/VisualIdeas/data/images/7db558b85005b5d9dbe14e44c3cd1092.png}
\end{center}
\vspace{2mm}
\begin{center}
\begin{minipage}{0.92\textwidth}
\small
This may only be enchanting to me, but I just updated this visual which was one of the first visuals I posted online. 🥰
\end{minipage}
\end{center}
\vspace*{\fill}
\newpage
\renewcommand{\theposttitle}{Cycle of Growth \& Intentionality}
\renewcommand{\thepostdate}{2022-10-04}
\vspace*{\fill}
\begin{center}
{\Huge\bfseries ANALOGY}\\[6mm]
\end{center}
\begin{center}
\includegraphics[width=\textwidth, height=0.55\textheight, keepaspectratio]{C:/GitHub/VisualIdeas/data/images/f33317c67d15a5c34ee0d496bfaa08f9.png}
\end{center}
\vspace{2mm}
\begin{center}
\begin{minipage}{0.92\textwidth}
\small
When I first drew this, it was a simple apple, but it felt that it just didn't carry the message completely.

Honestly, this is how I feel if I'm not careful. It's really easy to get caught up with getting through the day instead of living the day and this image was a personal meditation of that experience.
\end{minipage}
\end{center}
\vspace*{\fill}
\newpage
\renewcommand{\theposttitle}{Hard Decisions \& Deep Work}
\renewcommand{\thepostdate}{2022-10-18}
\vspace*{\fill}
\begin{center}
{\Huge\bfseries INTERESTING}\\[6mm]
\end{center}
\begin{center}
\includegraphics[width=\textwidth, height=0.55\textheight, keepaspectratio]{C:/GitHub/VisualIdeas/data/images/75475293c62d4cedd3859dbe432f9c57.png}
\end{center}
\vspace{2mm}
\begin{center}
\begin{minipage}{0.92\textwidth}
\small
I made this visualization as an experiment to see if it would make sense even though it breaks how most graphs are read.

The verdict? It works! (most everyone understood albeit maybe not agreeing with the message)

Of course there were plenty of reminders about how to use a y-axis 😅
\end{minipage}
\end{center}
\vspace*{\fill}
\newpage
\renewcommand{\theposttitle}{Hard Decisions \& Deep Work}
\renewcommand{\thepostdate}{2022-10-18}
\vspace*{\fill}
\begin{center}
{\Huge\bfseries DESIGN}\\[6mm]
\end{center}
\begin{center}
\includegraphics[width=\textwidth, height=0.55\textheight, keepaspectratio]{C:/GitHub/VisualIdeas/data/images/d979e77774c22d0675ffa93e78c70dd7.png}
\end{center}
\vspace{2mm}
\begin{center}
\begin{minipage}{0.92\textwidth}
\small
If there's a single tip to focus on to make dramatic improvements to your compositions, it's think about how you're using *diagonal lines* in your image.

In this case, a slight adjustment changes the type of perspective we're using.

Vanishing Point is most common, but Parallel Perspective is used a lot these days in isometric designs.
\end{minipage}
\end{center}
\vspace*{\fill}
\newpage
\renewcommand{\theposttitle}{Hard Decisions \& Deep Work}
\renewcommand{\thepostdate}{2022-10-18}
\vspace*{\fill}
\begin{center}
{\Huge\bfseries ENCHANTING}\\[6mm]
\end{center}
\begin{center}
\includegraphics[width=\textwidth, height=0.55\textheight, keepaspectratio]{C:/GitHub/VisualIdeas/data/images/43c39f5666c52917460e0700427b3e4d.png}
\end{center}
\vspace{2mm}
\begin{center}
\begin{minipage}{0.92\textwidth}
\small
Deep Insights = Focus x Time
\end{minipage}
\end{center}
\vspace*{\fill}
\newpage
\renewcommand{\theposttitle}{Hard Decisions \& Deep Work}
\renewcommand{\thepostdate}{2022-10-18}
\vspace*{\fill}
\begin{center}
{\Huge\bfseries ANALOGY}\\[6mm]
\end{center}
\begin{center}
\includegraphics[width=\textwidth, height=0.55\textheight, keepaspectratio]{C:/GitHub/VisualIdeas/data/images/5483ff303f7541a8748ae99b5f260b67.png}
\end{center}
\vspace{2mm}
\begin{center}
\begin{minipage}{0.92\textwidth}
\small
I don't publish a lot of what I write these days but the practice of crystalizing my ideas on paper is vital to gain clarity in thinking. I write and draw in the same book combining 'morning pages' with 'morning drawings.' It's a cleansing ritual I'd highly recommend.
\end{minipage}
\end{center}
\vspace*{\fill}
\newpage
\renewcommand{\theposttitle}{Creative Resistance and Finding Purpose}
\renewcommand{\thepostdate}{2022-11-01}
\vspace*{\fill}
\begin{center}
{\Huge\bfseries INTERESTING}\\[6mm]
\end{center}
\begin{center}
\includegraphics[width=\textwidth, height=0.55\textheight, keepaspectratio]{C:/GitHub/VisualIdeas/data/images/11f5061c3d3677c03624e5319326bf9d.png}
\end{center}
\vspace{2mm}
\begin{center}
\begin{minipage}{0.92\textwidth}
\small
Where we place our focus affects the range and depth of our expertise.

As generalists develop range by spreading their focus over many skills, research suggests they are better creative, out-of-the-box thinking, because they see dots from other domains unseen by the specialist.

As specialists develop filling singular skillsets with depth, they are more efficient and can probably better innovate on their given area of expertise.

There're situations where generalists shine and situations where specialists shine.

Where do you stand on the generalist vs specialist argument?
\end{minipage}
\end{center}
\vspace*{\fill}
\newpage
\renewcommand{\theposttitle}{Creative Resistance and Finding Purpose}
\renewcommand{\thepostdate}{2022-11-01}
\vspace*{\fill}
\begin{center}
{\Huge\bfseries DESIGN}\\[6mm]
\end{center}
\begin{center}
\includegraphics[width=\textwidth, height=0.55\textheight, keepaspectratio]{C:/GitHub/VisualIdeas/data/images/f280f729fbbc1e3d4186e96ec570bd06.png}
\end{center}
\vspace{2mm}
\begin{center}
\begin{minipage}{0.92\textwidth}
\small
I love using shadows in my work. They increase the overall turnaround time of finishing each illustration, but they're not just for looks. They're a practical tool to add to your tool belt for conveying space and where objects are in relativity to one another.
\end{minipage}
\end{center}
\vspace*{\fill}
\newpage
\renewcommand{\theposttitle}{Creative Resistance and Finding Purpose}
\renewcommand{\thepostdate}{2022-11-01}
\vspace*{\fill}
\begin{center}
{\Huge\bfseries ENCHANTING}\\[6mm]
\end{center}
\begin{center}
\includegraphics[width=\textwidth, height=0.55\textheight, keepaspectratio]{C:/GitHub/VisualIdeas/data/images/26dac188294a878b3dc2c2100ab52279.png}
\end{center}
\vspace{2mm}
\begin{center}
\begin{minipage}{0.92\textwidth}
\small
“Resistance will tell you anything to keep you from doing your work.”

Resistance is sneaky because it's always present.

It's like the proverbial fish in David Foster Wallace's now famous commencement speech, "This is Water."

If you've never heard it, here's an excerpt:

"There are these two young fish swimming along and they happen to meet an older fish swimming the other way, who nods at them and says:

'Morning, boys. How’s the water?'

And the two young fish swim on for a bit, and then eventually one of them looks over at the other and goes:

'What the hell is water?'"

-David Foster Wallace

Resistance is our water. It's the default. It slows progress.

But it's completely possible to push through to the other side.
\end{minipage}
\end{center}
\vspace*{\fill}
\newpage
\renewcommand{\theposttitle}{Creative Resistance and Finding Purpose}
\renewcommand{\thepostdate}{2022-11-01}
\vspace*{\fill}
\begin{center}
{\Huge\bfseries ANALOGY}\\[6mm]
\end{center}
\begin{center}
\includegraphics[width=\textwidth, height=0.55\textheight, keepaspectratio]{C:/GitHub/VisualIdeas/data/images/1bf0058cadea3b6c64856d11ad3f3aae.png}
\end{center}
\vspace{2mm}
\begin{center}
\begin{minipage}{0.92\textwidth}
\small
"Passion brings immediate joy.

Values provide lasting meaning."

There's something about the fuel a purpose-driven life provides that helps forge our hearts to be more resilient.

If you can somehow melt the two together, then you'll be able to sustain your mission even when the immediate burst of joy has burnt out.
\end{minipage}
\end{center}
\vspace*{\fill}
\newpage
\renewcommand{\theposttitle}{How Americans Spend Time and Shiny Object Syndrome}
\renewcommand{\thepostdate}{2022-12-06}
\vspace*{\fill}
\begin{center}
{\Huge\bfseries INTERESTING}\\[6mm]
\end{center}
\begin{center}
\begin{minipage}{0.92\textwidth}
\small
Sahil Bloom posted graphs earlier this month citing data from \textit{American Time Use Survey} about how the Western world spends their time.

I felt compelled to simplify these graphs, which I admit are sobering (even depressing) but they're also eye-opening.

Please be sure to check out the reminders under each image.

I'd even encourage you to consider your own takeaway because it wasn't till I thought of the reminders to myself that I felt hopeful.
\end{minipage}
\end{center}
\vspace{4mm}
\begin{center}
\includegraphics[width=\textwidth, height=0.13\textheight, keepaspectratio]{C:/GitHub/VisualIdeas/data/images/5c73af8b66bcdf36942940c7dc3a8e5a.png}
\end{center}
\vspace{2mm}
\begin{center}
\begin{minipage}{0.92\textwidth}
\small
\textbf{☝️Reminder to be present, even and especially during challenging times.}
\end{minipage}
\end{center}
\vspace{5mm}
\begin{center}
\includegraphics[width=\textwidth, height=0.13\textheight, keepaspectratio]{C:/GitHub/VisualIdeas/data/images/243d23bcb2256e17d319854b51fdbe49.png}
\end{center}
\vspace{2mm}
\begin{center}
\begin{minipage}{0.92\textwidth}
\small
\textbf{☝️}\textbf{Reminder to appreciate the limited time you have with your parents.}
\end{minipage}
\end{center}
\vspace{5mm}
\begin{center}
\includegraphics[width=\textwidth, height=0.13\textheight, keepaspectratio]{C:/GitHub/VisualIdeas/data/images/9dca997d325d61cfc09df4cce1ea9570.png}
\end{center}
\vspace{2mm}
\begin{center}
\begin{minipage}{0.92\textwidth}
\small
\textbf{☝️}\textbf{Reminder to choose the friends you share your time with wisely.}
\end{minipage}
\end{center}
\vspace{5mm}
\begin{center}
\includegraphics[width=\textwidth, height=0.13\textheight, keepaspectratio]{C:/GitHub/VisualIdeas/data/images/198f4fcafc69815799f9d88ef5b3da66.png}
\end{center}
\vspace{2mm}
\begin{center}
\begin{minipage}{0.92\textwidth}
\small
\textbf{☝️}\textbf{Reminder to embody the kind of person you'd want to be around.}
\end{minipage}
\end{center}
\vspace*{\fill}
\newpage
\renewcommand{\theposttitle}{How Americans Spend Time and Shiny Object Syndrome}
\renewcommand{\thepostdate}{2022-12-06}
\vspace*{\fill}
\begin{center}
{\Huge\bfseries DESIGN}\\[6mm]
\end{center}
\begin{center}

\includegraphics[width=\textwidth, height=0.55\textheight, keepaspectratio]{C:/GitHub/VisualIdeas/data/images/0a6d20fffeef62c5583b0d13f046b88d.png}
\end{center}
\vspace{2mm}
\begin{center}
\begin{minipage}{0.92\textwidth}
\small
I don't have a design tip this issue, but I have a design \textit{tool} to share.

If you've been thinking of playing with AI art, Midjourney is quite compelling. It takes a bit of a learning curve, but it's scary good.
\end{minipage}
\end{center}
\vspace*{\fill}
\newpage
\renewcommand{\theposttitle}{How Americans Spend Time and Shiny Object Syndrome}
\renewcommand{\thepostdate}{2022-12-06}
\vspace*{\fill}
\begin{center}
{\Huge\bfseries ENCHANTING}\\[6mm]
\end{center}
\begin{center}
\includegraphics[width=\textwidth, height=0.55\textheight, keepaspectratio]{C:/GitHub/VisualIdeas/data/images/1d538845a41e5626e20c4afcd48b5052.png}
\end{center}
\vspace{2mm}
\begin{center}
\begin{minipage}{0.92\textwidth}
\small
Write between the lines. (this visual took me ages to get right)
\end{minipage}
\end{center}
\vspace*{\fill}
\newpage
\renewcommand{\theposttitle}{How Americans Spend Time and Shiny Object Syndrome}
\renewcommand{\thepostdate}{2022-12-06}
\vspace*{\fill}
\begin{center}
{\Huge\bfseries ANALOGY}\\[6mm]
\end{center}
\begin{center}
\includegraphics[width=\textwidth, height=0.55\textheight, keepaspectratio]{C:/GitHub/VisualIdeas/data/images/a580fc45ca634d7de2253562ba53e9ac.png}
\end{center}
\vspace{2mm}
\begin{center}
\begin{minipage}{0.92\textwidth}
\small
\textbf{Shiny object syndrome} is when you're constantly jumping from project to project because you think the next new thing is going to bring more value than the current thing.

It's essentially thinking "the grass is greener on the other side" and then hopping from one lawn to the next in search of the greenest grass. The irony is that the reason the grass is greener somewhere else is because you never stay long enough where you are to take care of it.

I have such a problem with shiny object syndrome. I do it with books, with hobbies, even with cleaning the house (much to my wife's irritation).

Recently I'm learning how to leverage my focus for deeper value. And any time I've been able to stay focused and dive deep on one thing, I find not only actual results, but deeper fulfillment because I \textit{feel} the progress.
\end{minipage}
\end{center}
\vspace*{\fill}
\newpage
\renewcommand{\theposttitle}{Collecting Ideas vs. Creating Outcomes}
\renewcommand{\thepostdate}{2023-02-08}
\vspace*{\fill}
\begin{center}
{\Huge\bfseries DESIGN}\\[6mm]
\end{center}
\begin{center}
\includegraphics[width=\textwidth, height=0.55\textheight, keepaspectratio]{C:/GitHub/VisualIdeas/data/images/774fbbca0f14196baff2dbd82a88f0d4.png}
\end{center}
\vspace{2mm}
\begin{center}
\begin{minipage}{0.92\textwidth}
\small
\textbf{The Unsuspecting Balloon Metaphor} is a classic. The key to doing this one well is to reveal something the audience is not expecting that's a threat to what they want. Bonus points for specificity.Classic Example: \textbf{Positive Desired Outcome:} Ability to Focus\textbf{Non-obvious Thread:} A Single Notification
\end{minipage}
\end{center}
\vspace*{\fill}
\newpage
\renewcommand{\theposttitle}{Collecting Ideas vs. Creating Outcomes}
\renewcommand{\thepostdate}{2023-02-08}
\vspace*{\fill}
\begin{center}
{\Huge\bfseries ENCHANTING}\\[6mm]
\end{center}
\begin{center}
\includegraphics[width=\textwidth, height=0.55\textheight, keepaspectratio]{C:/GitHub/VisualIdeas/data/images/25f11ffd5e92a861cf854e0e59b28a27.png}
\end{center}
\vspace{2mm}
\begin{center}
\begin{minipage}{0.92\textwidth}
\small
\textbf{Collecting Ideas vs. Creating Outcomes}There's a concept called the collector's fallacy:

Believing that by collecting information, we somehow absorb it.

It gives us a false sense of progress.

There's nothing wrong with reading for pleasure. But if you're reading to learn how to do something, there's no substitute for action. I hesitated putting my work online for years because I didn't know 'exactly' what I wanted to do.

Finally started posting in March '22 and things started happening. Real progress is made only through behavior and action.
\end{minipage}
\end{center}
\vspace*{\fill}
\newpage
\renewcommand{\theposttitle}{Collecting Ideas vs. Creating Outcomes}
\renewcommand{\thepostdate}{2023-02-08}
\vspace*{\fill}
\begin{center}
{\Huge\bfseries ANALOGY}\\[6mm]
\end{center}
\begin{center}
\includegraphics[width=\textwidth, height=0.55\textheight, keepaspectratio]{C:/GitHub/VisualIdeas/data/images/5332b6a92259c9472debe33fe252bb45.png}
\end{center}
\vspace{2mm}
\begin{center}
\begin{minipage}{0.92\textwidth}
\small
\textbf{"Reading is essential for those who seek to rise above the ordinary."}-Jim Rohn
\end{minipage}
\end{center}
\vspace*{\fill}
\newpage
\renewcommand{\theposttitle}{Fluid vs Crystalized Intelligence}
\renewcommand{\thepostdate}{2023-03-01}
\vspace*{\fill}
\begin{center}
{\Huge\bfseries INTERESTING}\\[6mm]
\end{center}
\begin{center}
\includegraphics[width=\textwidth, height=0.55\textheight, keepaspectratio]{C:/GitHub/VisualIdeas/data/images/5d6dff56e22281550173223d586e5cbc.png}
\end{center}
\vspace{2mm}
\begin{center}
\begin{minipage}{0.92\textwidth}
\small
I've been thinking about age a lot recently so I was intrigued to come across the above concept from Catell's Intelligence Theory. Apparently, many physicists, scientists, and even entrepreneurs hit their innovation peak early in their careers. And they waste the latter end of their lives miserable because they continue to chase success to rival that of their youth. Turns out, there's a scientific reason for this inability.

There're two types of intelligence (and they're affected by age).

\textbf{Fluid Intelligence:} Ability to solve problems by coming up with new innovative ways to attack the problem. Peaks somewhere between ages 20-25 and then slowly declines afterward.

\textbf{Crystalized Intelligence:} Ability to solve problems from accumulated knowledge (make connections). This grows stronger with age.

The problem comes when people neglect making the jump from fluid intelligence to leverage crystallized intelligence during midlife. Important to note, these intelligences are not related to IQ.

From the book: "From Strength to Strength" by Arthur C. Brooks. (amazon link)
\end{minipage}
\end{center}
\vspace*{\fill}
\newpage
\renewcommand{\theposttitle}{Fluid vs Crystalized Intelligence}
\renewcommand{\thepostdate}{2023-03-01}
\vspace*{\fill}
\begin{center}
{\Huge\bfseries DESIGN}\\[6mm]
\end{center}
\begin{center}
\includegraphics[width=\textwidth, height=0.55\textheight, keepaspectratio]{C:/GitHub/VisualIdeas/data/images/958a1b03244bee19854c1f7ae2ef4c5f.png}
\end{center}
\vspace{2mm}
\begin{center}
\begin{minipage}{0.92\textwidth}
\small
For this week's template, we're tapping into UX which is a valuable inspiration for visual thinking.

Example: "It's time to work out." Choice 1: "LET'S GO" Choice 2: "Snooze"

In the above example, the question is optional (it can even be a declarative sentence or a warning signal). It's just used to frame the context of the choice.

Note where the mouse is. You can hover it over a choice to nudge your audience to one of the choices. Alternatively, you can hover the mouse in the middle to let your audience decide.
\end{minipage}
\end{center}
\vspace*{\fill}
\newpage
\renewcommand{\theposttitle}{Fluid vs Crystalized Intelligence}
\renewcommand{\thepostdate}{2023-03-01}
\vspace*{\fill}
\begin{center}
{\Huge\bfseries ENCHANTING}\\[6mm]
\end{center}
\begin{center}
\includegraphics[width=\textwidth, height=0.55\textheight, keepaspectratio]{C:/GitHub/VisualIdeas/data/images/c2821c8436c6f7d1db13a7770b619547.png}
\end{center}
\vspace{2mm}
\begin{center}
\begin{minipage}{0.92\textwidth}
\small
\textbf{Success is a team support.}
\end{minipage}
\end{center}
\vspace*{\fill}
\newpage
\renewcommand{\theposttitle}{Fluid vs Crystalized Intelligence}
\renewcommand{\thepostdate}{2023-03-01}
\vspace*{\fill}
\begin{center}
{\Huge\bfseries ANALOGY}\\[6mm]
\end{center}
\begin{center}
\includegraphics[width=\textwidth, height=0.55\textheight, keepaspectratio]{C:/GitHub/VisualIdeas/data/images/b861421f643569c6fc3a2adc7a2578eb.png}
\end{center}
\vspace{2mm}
\begin{center}
\begin{minipage}{0.92\textwidth}
\small
\textbf{Use a creative lens to see the master pieces.}
\end{minipage}
\end{center}
\vspace*{\fill}
\newpage
\renewcommand{\theposttitle}{Finding Inspiration and Hofstadter's Law}
\renewcommand{\thepostdate}{2023-03-14}
\vspace*{\fill}
\begin{center}
{\Huge\bfseries INTERESTING}\\[6mm]
\end{center}
\begin{center}
\includegraphics[width=\textwidth, height=0.55\textheight, keepaspectratio]{C:/GitHub/VisualIdeas/data/images/83267f37543d91d3b4b7158f09afd947.png}
\end{center}
\vspace{2mm}
\begin{center}
\begin{minipage}{0.92\textwidth}
\small
\textbf{Hofstadter's Law }is helpful to consider when planning for something. It states that we always underestimate how long it'll take to complete a task, even when we're conscious of Hofstadter's Law.

This law has single-mindedly targeted me all my life. 😂
\end{minipage}
\end{center}
\vspace*{\fill}
\newpage
\renewcommand{\theposttitle}{Finding Inspiration and Hofstadter's Law}
\renewcommand{\thepostdate}{2023-03-14}
\vspace*{\fill}
\begin{center}
{\Huge\bfseries DESIGN}\\[6mm]
\end{center}
\begin{center}

\includegraphics[width=\textwidth, height=0.24\textheight, keepaspectratio]{C:/GitHub/VisualIdeas/data/images/ca53f3c8c8d50ed092a2adc69c3e3fb5.png}
\end{center}
\vspace{2mm}
\begin{center}
\begin{minipage}{0.92\textwidth}
\small
What's important about this template for visual thinking is the small yolk has to be considered a part of the bigger reality too.

Here's one way I used it recently:
\end{minipage}
\end{center}
\vspace{5mm}
\begin{center}

\includegraphics[width=\textwidth, height=0.24\textheight, keepaspectratio]{C:/GitHub/VisualIdeas/data/images/212ecbfbc04eb997ddcc37c946f37c1c.png}
\end{center}
\vspace*{\fill}
\newpage
\renewcommand{\theposttitle}{Finding Inspiration and Hofstadter's Law}
\renewcommand{\thepostdate}{2023-03-14}
\vspace*{\fill}
\begin{center}
{\Huge\bfseries ENCHANTING}\\[6mm]
\end{center}
\begin{center}
\includegraphics[width=\textwidth, height=0.55\textheight, keepaspectratio]{C:/GitHub/VisualIdeas/data/images/dc43a3771c86d16b2096e4d244d0b157.png}
\end{center}
\vspace{2mm}
\begin{center}
\begin{minipage}{0.92\textwidth}
\small
\textbf{The faintest ink is more powerful than the strongest memory. }-Chinese Proverb

Literally inspired in the shower with this one 😂
\end{minipage}
\end{center}
\vspace*{\fill}
\newpage
\renewcommand{\theposttitle}{Finding Inspiration and Hofstadter's Law}
\renewcommand{\thepostdate}{2023-03-14}
\vspace*{\fill}
\begin{center}
{\Huge\bfseries ANALOGY}\\[6mm]
\end{center}
\begin{center}
\includegraphics[width=\textwidth, height=0.55\textheight, keepaspectratio]{C:/GitHub/VisualIdeas/data/images/4a3ade94062600bd699b36e6d90993fb.png}
\end{center}
\vspace{2mm}
\begin{center}
\begin{minipage}{0.92\textwidth}
\small
\textbf{The water you drink affects the ideas you think.}

As with all my visuals, there're elements of personal experience embedded in the image. For me:

Netflix - consumed me instead of me consuming it. (I felt like my watch list was a to-do list!)

TikTok - there's educational content but it's really difficult to be intentional in this format for me.

News - I have personally not been inspired to create anything after watching the news. It usually paralyzes me with despair.

We fuel or pollute our creative well with the content we consume. Curious, what fuels (or pollutes) you to be inspired to create something?
\end{minipage}
\end{center}
\vspace*{\fill}
\newpage
\renewcommand{\theposttitle}{Mental Models and Setting a Priority}
\renewcommand{\thepostdate}{2023-03-28}
\vspace*{\fill}
\begin{center}
{\Huge\bfseries INTERESTING}\\[6mm]
\end{center}
\begin{center}
\includegraphics[width=\textwidth, height=0.55\textheight, keepaspectratio]{C:/GitHub/VisualIdeas/data/images/b6c2b0898b7571285b59d2fadc0fdc55.png}
\end{center}
\vspace{2mm}
\begin{center}
\begin{minipage}{0.92\textwidth}
\small
I've been diving more into visualizing mental models as they apply to visual thinking (also applies in writing).

Here's how to use \textbf{Occam's Razor. }

• Choose the simplest explanation• Remove distracting details• Keep only the essential

Why use it?• Increases clarity• Amplifies your message• Simplifies your approach
\end{minipage}
\end{center}
\vspace*{\fill}
\newpage
\renewcommand{\theposttitle}{Mental Models and Setting a Priority}
\renewcommand{\thepostdate}{2023-03-28}
\vspace*{\fill}
\begin{center}
{\Huge\bfseries DESIGN}\\[6mm]
\end{center}
\begin{center}
\includegraphics[width=\textwidth, height=0.24\textheight, keepaspectratio]{C:/GitHub/VisualIdeas/data/images/eec9e59002d6f25abefb63eb77ba0210.png}
\end{center}
\vspace{2mm}
\begin{center}
\begin{minipage}{0.92\textwidth}
\small
This is one framework I use myself pretty regularly. It's a template communicating a \textbf{comparison in quantity}, similar to a pie chart.

\textbf{Important note:} As with any visual, the most important thing is to communicate the difference and not the exact numbers. In my work, I focus on communicating the spirit of the idea. I am not so much creating a data visualization.

Here's an example of it in action:
\end{minipage}
\end{center}
\vspace{5mm}
\begin{center}
\includegraphics[width=\textwidth, height=0.24\textheight, keepaspectratio]{C:/GitHub/VisualIdeas/data/images/b5692b27f866ad06fdb32f55e0887199.png}
\end{center}
\vspace*{\fill}
\newpage
\renewcommand{\theposttitle}{Mental Models and Setting a Priority}
\renewcommand{\thepostdate}{2023-03-28}
\vspace*{\fill}
\begin{center}
{\Huge\bfseries ENCHANTING}\\[6mm]
\end{center}
\begin{center}
\includegraphics[width=\textwidth, height=0.55\textheight, keepaspectratio]{C:/GitHub/VisualIdeas/data/images/67c165b435353791ba61ef8b5ba1f471.png}
\end{center}
\vspace{2mm}
\begin{center}
\begin{minipage}{0.92\textwidth}
\small
Being a parent is not easy. In the same breath you can be fed up with frustration and filled up with love. But one of my most treasured and surreal memories of fatherhood was teaching my son to ride a bicycle. I get emotional just thinking about it.
\end{minipage}
\end{center}
\vspace*{\fill}
\newpage
\renewcommand{\theposttitle}{Mental Models and Setting a Priority}
\renewcommand{\thepostdate}{2023-03-28}
\vspace*{\fill}
\begin{center}
{\Huge\bfseries ANALOGY}\\[6mm]
\end{center}
\begin{center}
\includegraphics[width=\textwidth, height=0.55\textheight, keepaspectratio]{C:/GitHub/VisualIdeas/data/images/8ef554cb5a580ef8243dcb557d9ad953.png}
\end{center}
\vspace{2mm}
\begin{center}
\begin{minipage}{0.92\textwidth}
\small
\textbf{If you don't prioritize your life someone else will.}- Gregory McKeown, author of EssentialismI've been thinking a lot about the word "priority" lately. In McKeown's book, he talks about how the word priority initially meant "one thing." The idea of multiple priorities is a modern dilemma.

I haven't landed on what my singular priority is yet, but it's something I'm working through. Coming to grips with not being able to do everything is tough, but I believe it might help me feel a bit less overwhelmed.

What about you? What are you prioritizing in this season of life? Please feel free to reply in this email (I read every one).
\end{minipage}
\end{center}
\vspace*{\fill}
\newpage
\renewcommand{\theposttitle}{Telephoto Thinking vs Wide-Angle Thinking}
\renewcommand{\thepostdate}{2023-04-11}
\vspace*{\fill}
\begin{center}
{\Huge\bfseries INTERESTING}\\[6mm]
\end{center}
\begin{center}
\includegraphics[width=\textwidth, height=0.55\textheight, keepaspectratio]{C:/GitHub/VisualIdeas/data/images/badb981bf35963e6f26bb34abc3d301e.png}
\end{center}
\vspace{2mm}
\begin{center}
\begin{minipage}{0.92\textwidth}
\small
\textbf{Which are you--a telephoto or wide-angle lens thinker?}

Most people want to be both (variable lens) but if forced to choose, I'd say most starting online content creators lean toward generalists.

The thing about telephoto lenses is that they're extremely difficult to focus unless you place them on a stable footing. This feels in line with how most creatives feel when they're told to niche down and specialize if they want to develop an identity to be recognized online.
\end{minipage}
\end{center}
\vspace*{\fill}
\newpage
\renewcommand{\theposttitle}{Telephoto Thinking vs Wide-Angle Thinking}
\renewcommand{\thepostdate}{2023-04-11}
\vspace*{\fill}
\begin{center}
{\Huge\bfseries DESIGN}\\[6mm]
\end{center}
\begin{center}
\includegraphics[width=\textwidth, height=0.55\textheight, keepaspectratio]{C:/GitHub/VisualIdeas/data/images/4d4d7a731fc4414a07498017838ab9f7.png}
\end{center}
\vspace{2mm}
\begin{center}
\begin{minipage}{0.92\textwidth}
\small
\textbf{"It gets easier" Design. }

The goal of this design is to change the perception of what seems impossible at first. Essentially create progressively smaller versions of a symbol representing what you want to convey getting easier over time.
\end{minipage}
\end{center}
\vspace*{\fill}
\newpage
\renewcommand{\theposttitle}{Telephoto Thinking vs Wide-Angle Thinking}
\renewcommand{\thepostdate}{2023-04-11}
\vspace*{\fill}
\begin{center}
{\Huge\bfseries ENCHANTING}\\[6mm]
\end{center}
\begin{center}
\includegraphics[width=\textwidth, height=0.55\textheight, keepaspectratio]{C:/GitHub/VisualIdeas/data/images/dd9b15df45bbf28633fca79e435e9e8b.png}
\end{center}
\vspace{2mm}
\begin{center}
\begin{minipage}{0.92\textwidth}
\small
\textbf{"You can do anything, but not everything."}-Greg McKeown

A reminder I got tired of having to tell myself so I illustrated the message hoping it will stick! 😁
\end{minipage}
\end{center}
\vspace*{\fill}
\newpage
\renewcommand{\theposttitle}{Telephoto Thinking vs Wide-Angle Thinking}
\renewcommand{\thepostdate}{2023-04-11}
\vspace*{\fill}
\begin{center}
{\Huge\bfseries ANALOGY}\\[6mm]
\end{center}
\begin{center}
\includegraphics[width=\textwidth, height=0.55\textheight, keepaspectratio]{C:/GitHub/VisualIdeas/data/images/e8f70e6cc9303004d74aba687bef4717.png}
\end{center}
\vspace{2mm}
\begin{center}
\begin{minipage}{0.92\textwidth}
\small
Unfortunately, getting older means also feeling the inevitable loss of the people we love most. And each time feels like the first.

In these moments, I find myself feeling tremendous guilt and shame. Guilt if I try to live a life without them and shame I didn't do more when they were alive.

There's a beautiful TED talk I was reminded of recently off a grief post from Michael J. Boorman (@wisdommadeeasy) that we don't ever really move on from grief. We move forward with it.

A quote I think about that helps me move forward with grief:

"What is grief if not love persevering?"
\end{minipage}
\end{center}
\vspace*{\fill}
\newpage
\renewcommand{\theposttitle}{Hard Decisions and the Clarity of Humility}
\renewcommand{\thepostdate}{2023-05-02}
\vspace*{\fill}
\begin{center}
{\Huge\bfseries INTERESTING}\\[6mm]
\end{center}
\begin{center}
\includegraphics[width=\textwidth, height=0.55\textheight, keepaspectratio]{C:/GitHub/VisualIdeas/data/images/8da9bf15f767b0c1cfc0af1726d5dbca.png}
\end{center}
\vspace{2mm}
\begin{center}
\begin{minipage}{0.92\textwidth}
\small
Beware the path of least resistance. \\
Careful walking the fine line between thoughtful reflection and willful procrastination as the urgent often overshadows the important.\\
Read the full post here.

I had an interesting comment to this visual on LinkedIn. Specifically, that on the opposite side of “Pain” should be the word “Pleasure.” I respectfully disagreed. The absence of pain is not necessarily pleasure. 

What do you think?
\end{minipage}
\end{center}
\vspace*{\fill}
\newpage
\renewcommand{\theposttitle}{Hard Decisions and the Clarity of Humility}
\renewcommand{\thepostdate}{2023-05-02}
\vspace*{\fill}
\begin{center}
{\Huge\bfseries DESIGN}\\[6mm]
\end{center}
\begin{center}
\includegraphics[width=\textwidth, height=0.55\textheight, keepaspectratio]{C:/GitHub/VisualIdeas/data/images/2750da5366f1b38fefb6b1e8bcdff7c1.png}
\end{center}
\vspace{2mm}
\begin{center}
\begin{minipage}{0.92\textwidth}
\small
This visual uses a \textbf{“This to That” Framework}. All that means is comparing \textbf{this} (fear in imagination) to \textbf{that} (fear in reality). It’s a simple but powerful device to help you frame something in context.

Couple this with one of my favorite creative levers to pull: \textit{exaggeration}.

Exaggeration can help you interpret an idea specifically to your own point of view.

Be careful though. Exaggeration is a double-edged sword in visual thinking. An exaggerated analogy may have more difficulty getting your message across because a more literal audience might get stuck in saying, “It’s not really like that.”
\end{minipage}
\end{center}
\vspace*{\fill}
\newpage
\renewcommand{\theposttitle}{Hard Decisions and the Clarity of Humility}
\renewcommand{\thepostdate}{2023-05-02}
\vspace*{\fill}
\begin{center}
{\Huge\bfseries ENCHANTING}\\[6mm]
\end{center}
\begin{center}
\includegraphics[width=\textwidth, height=0.55\textheight, keepaspectratio]{C:/GitHub/VisualIdeas/data/images/a5a9179e83c9e0d66ecd508f0305d102.png}
\end{center}
\vspace{2mm}
\begin{center}
\begin{minipage}{0.92\textwidth}
\small
Name and recognize failures, but focus more on ingraining the lessons they teach you.

Let temporary setbacks wash away.
\end{minipage}
\end{center}
\vspace*{\fill}
\newpage
\renewcommand{\theposttitle}{Hard Decisions and the Clarity of Humility}
\renewcommand{\thepostdate}{2023-05-02}
\vspace*{\fill}
\begin{center}
{\Huge\bfseries ANALOGY}\\[6mm]
\end{center}
\begin{center}
\includegraphics[width=\textwidth, height=0.17\textheight, keepaspectratio]{C:/GitHub/VisualIdeas/data/images/ebab181dd6723a792b4ddd632fec62e4.png}
\end{center}
\vspace{2mm}
\begin{center}
\begin{minipage}{0.92\textwidth}
\small
\textbf{An empty pot makes the loudest noise.}\\
-Swahili Proverb
\end{minipage}
\end{center}
\vspace{5mm}
\begin{center}

\includegraphics[width=\textwidth, height=0.17\textheight, keepaspectratio]{C:/GitHub/VisualIdeas/data/images/d8cc89735aa0543f63af3a516be0dc7f.png}
\end{center}
\vspace{2mm}
\begin{center}
\begin{minipage}{0.92\textwidth}
\small
Say hi 👋🏻 on Twitter or LinkedIn

If you'd like to join the 3,244 subscribers getting I.D.E.A.s straight into their inbox, subscribe below:

Thanks for reading! Reply any time with feedback or questions.

I made this newsletter with \textbf{beehiiv}\textbf{, }the best place to start and scale a newsletter in 2023.

When you’re ready, I’m working on building a course to teach visual thinking: Join the Waiting List
\end{minipage}
\end{center}
\vspace{5mm}
\begin{center}

\includegraphics[width=\textwidth, height=0.17\textheight, keepaspectratio]{C:/GitHub/VisualIdeas/data/images/f2833dff016936c751c34b46c7d76fa1.png}
\end{center}
\vspace*{\fill}
\newpage
\renewcommand{\theposttitle}{Authenticity}
\renewcommand{\thepostdate}{2023-05-30}
\vspace*{\fill}
\begin{center}
{\Huge\bfseries INTERESTING}\\[6mm]
\end{center}
\begin{center}
\includegraphics[width=\textwidth, height=0.24\textheight, keepaspectratio]{C:/GitHub/VisualIdeas/data/images/5e2bb680e892eba7ff1c35face045ee5.png}
\end{center}
\vspace{2mm}
\begin{center}
\begin{minipage}{0.92\textwidth}
\small
How will we look at these tools 20 years from now?

Every corner of the internet is buzzing with “AI this” or “AI that.” 😵‍💫

I’m still working out how I feel about it as a teacher. but love it or hate it, it’s here to stay. I wonder if people will one day look at AI-assisted writing in English class the we way we look at calculators in Math.

On the other hand:
\end{minipage}
\end{center}
\vspace{5mm}
\begin{center}
\includegraphics[width=\textwidth, height=0.24\textheight, keepaspectratio]{C:/GitHub/VisualIdeas/data/images/cd97b3718cffcf53f403a27e6930f506.png}
\end{center}
\vspace{2mm}
\begin{center}
\begin{minipage}{0.92\textwidth}
\small
This is no way to build strength.
\end{minipage}
\end{center}
\vspace*{\fill}
\newpage
\renewcommand{\theposttitle}{Authenticity}
\renewcommand{\thepostdate}{2023-05-30}
\vspace*{\fill}
\begin{center}
{\Huge\bfseries DESIGN}\\[6mm]
\end{center}
\begin{center}
\includegraphics[width=\textwidth, height=0.55\textheight, keepaspectratio]{C:/GitHub/VisualIdeas/data/images/32bf1e9ae115ca7a30fce009d9a2526c.png}
\end{center}
\vspace{2mm}
\begin{center}
\begin{minipage}{0.92\textwidth}
\small
The books I most often recommend to start learning visual thinking:

Visual Thinking by Williemien Brand

Visual Doing by Williemien Brand

Non-Designer's Design Book by Robin Williams

Illustration School: Let's Draw a Story by Sachiko Umoto
\end{minipage}
\end{center}
\vspace*{\fill}
\newpage
\renewcommand{\theposttitle}{Authenticity}
\renewcommand{\thepostdate}{2023-05-30}
\vspace*{\fill}
\begin{center}
{\Huge\bfseries ENCHANTING}\\[6mm]
\end{center}
\begin{center}
\includegraphics[width=\textwidth, height=0.55\textheight, keepaspectratio]{C:/GitHub/VisualIdeas/data/images/0aaae0c6aae10f566a16d1f14ff65181.png}
\end{center}
\vspace{2mm}
\begin{center}
\begin{minipage}{0.92\textwidth}
\small
\textbf{We complete each other.}
\end{minipage}
\end{center}
\vspace*{\fill}
\newpage
\renewcommand{\theposttitle}{Authenticity}
\renewcommand{\thepostdate}{2023-05-30}
\vspace*{\fill}
\begin{center}
{\Huge\bfseries ANALOGY}\\[6mm]
\end{center}
\begin{center}
\includegraphics[width=\textwidth, height=0.55\textheight, keepaspectratio]{C:/GitHub/VisualIdeas/data/images/f4bccdc24537d9b0b6cf175e8bc9325e.png}
\end{center}
\vspace{2mm}
\begin{center}
\begin{minipage}{0.92\textwidth}
\small
\textbf{It is the cracked pot that lets the light in.}\\
-Leonard Cohen

Not a direct quote, but a paraphrased one.

This image drew up some challenging comments on Instagram. A few people remarked how they were betrayed by revealing too much to the people they trusted.

This reminded me how I’ve been burned too, more times than I can count. In fact, the one thing I remember from my dad growing up is how many times his friends swindled him out of his hard-earned money because he revealed too many scars, which created real fears in me around trust.

But I’ve also been lucky to meet people who restored some of that trust. \\
• An art teacher who encouraged me out of a sea of 4,000 other low-income, public school students to keep creating. \\
• An elderly Japanese man in Hiroshima who saw me lost in the train station and spent 2 hours to help me get to where I was going. \\
• A friend I just met in Thailand inviting me over to share in Christmas dinner and presents so I wouldn’t spend the holidays alone.

Some scars heal and some scars fester. All I know is I saw what the festering did to my dad and that’s something I’d like to avoid.

I’m curious, how do you navigate revealing who you are on social media?
\end{minipage}
\end{center}
\vspace*{\fill}
\newpage
\renewcommand{\theposttitle}{How to get the morning gunk out of your head}
\renewcommand{\thepostdate}{2023-06-20}
\vspace*{\fill}
\begin{center}
{\Huge\bfseries INTERESTING}\\[6mm]
\end{center}
\begin{center}
\includegraphics[width=\textwidth, height=0.55\textheight, keepaspectratio]{C:/GitHub/VisualIdeas/data/images/336a3df0f51023ce33f1440a740f4ca8.png}
\end{center}
\vspace{2mm}
\begin{center}
\begin{minipage}{0.92\textwidth}
\small
Clear your mind for the day.

Inspired by the concept of Morning Pages from the book The Artist’s Way by Julia Cameron. But honestly, any-time journaling is better than nothing. 

I have one on hand whenever and wherever I go to 'process' and get the gunk out. 🤓
\end{minipage}
\end{center}
\vspace*{\fill}
\newpage
\renewcommand{\theposttitle}{How to get the morning gunk out of your head}
\renewcommand{\thepostdate}{2023-06-20}
\vspace*{\fill}
\begin{center}
{\Huge\bfseries DESIGN}\\[6mm]
\end{center}
\begin{center}
\begin{minipage}{0.92\textwidth}
\small
I recently iterated on my “Imposter Syndrome to Growth Mindset” visual. 

What design problem pops out to you from this first iteration? ↴
\end{minipage}
\end{center}
\vspace{4mm}
\begin{center}
\includegraphics[width=\textwidth, height=0.17\textheight, keepaspectratio]{C:/GitHub/VisualIdeas/data/images/971a88c423fd26d133da4fcd598dd46b.png}
\end{center}
\vspace{2mm}
\begin{center}
\begin{minipage}{0.92\textwidth}
\small
That’s right! The point of intersection between the graphs! 

Where the graphs overlap needs to make insightful sense to both “growth mindsets” and “imposter syndrome.” Unfortunately, the point of intersection wasn’t the point of the message!

So back to the drawing board (🥁)
\end{minipage}
\end{center}
\vspace{5mm}
\begin{center}
\includegraphics[width=\textwidth, height=0.17\textheight, keepaspectratio]{C:/GitHub/VisualIdeas/data/images/7ba2aa7ff7884883ba0b8d41fff1d658.png}
\end{center}
\vspace{2mm}
\begin{center}
\begin{minipage}{0.92\textwidth}
\small
My first attempt

• Reads top-down and left-to-right, so problem = growth mindset is seen first

• Red graph dominates the image

• Thought bubbles make less sense

• Helps me identify the central problem = Reveal each element so it can be read in the correct order

Below is the final version solving the central design problem. Curious, how might you have solved this differently?
\end{minipage}
\end{center}
\vspace{5mm}
\begin{center}
\includegraphics[width=\textwidth, height=0.17\textheight, keepaspectratio]{C:/GitHub/VisualIdeas/data/images/d7b41c6174cc33b86f8ba5af7b155aac.png}
\end{center}
\vspace{2mm}
\begin{center}
\begin{minipage}{0.92\textwidth}
\small
After
\end{minipage}
\end{center}
\vspace*{\fill}
\newpage
\renewcommand{\theposttitle}{How to get the morning gunk out of your head}
\renewcommand{\thepostdate}{2023-06-20}
\vspace*{\fill}
\begin{center}
{\Huge\bfseries ENCHANTING}\\[6mm]
\end{center}
\begin{center}
\includegraphics[width=\textwidth, height=0.55\textheight, keepaspectratio]{C:/GitHub/VisualIdeas/data/images/e52d38aa550159dbdfffb51f467d3181.png}
\end{center}
\vspace{2mm}
\begin{center}
\begin{minipage}{0.92\textwidth}
\small
What’s a book that completely changed how you think about how you think?

In case you’re interested in recommendations this summer, here are my three favorites of the last year:

Four Thousand Weeks: Time Management for Mortals by Oliver Burkeman

From Strength to Strength: Finding Success, Happiness, and Deep Purpose in the Second Half of Life by Arthur C. Brooks

The Anthropocene Reviewed by John Green
\end{minipage}
\end{center}
\vspace*{\fill}
\newpage
\renewcommand{\theposttitle}{How to get the morning gunk out of your head}
\renewcommand{\thepostdate}{2023-06-20}
\vspace*{\fill}
\begin{center}
{\Huge\bfseries ANALOGY}\\[6mm]
\end{center}
\begin{center}
\includegraphics[width=\textwidth, height=0.55\textheight, keepaspectratio]{C:/GitHub/VisualIdeas/data/images/7d5775ddd2dafa8d5f7486fcd8153f4d.png}
\end{center}
\vspace{2mm}
\begin{center}
\begin{minipage}{0.92\textwidth}
\small
Inspired by the Stoic version of the serenity prayer:

\textit{“The chief task in life is simply this: to identify and separate matters so that I can say clearly to myself which are externals not under my control, and which have to do with the choices I actually control. }

\textit{Where then do I look for good and evil? }

\textit{Not to uncontrollable externals, but within myself to the choices that are my own…” }\\
— Epictetus, Discourses, 2.5.4–5

Login or Subscribe to participate
\end{minipage}
\end{center}
\vspace*{\fill}
\newpage
\renewcommand{\theposttitle}{Growth Minded Celebrations and Disagreements}
\renewcommand{\thepostdate}{2023-07-04}
\vspace*{\fill}
\begin{center}
{\Huge\bfseries INTERESTING}\\[6mm]
\end{center}
\begin{center}
\includegraphics[width=\textwidth, height=0.55\textheight, keepaspectratio]{C:/GitHub/VisualIdeas/data/images/2b453d0ec14d657a3b483629a0bb5049.png}
\end{center}
\vspace{2mm}
\begin{center}
\begin{minipage}{0.92\textwidth}
\small
The correct way to infuse life

This has been an ongoing theme this year for me: trying to remember that life is more than work.
\end{minipage}
\end{center}
\vspace*{\fill}
\newpage
\renewcommand{\theposttitle}{Growth Minded Celebrations and Disagreements}
\renewcommand{\thepostdate}{2023-07-04}
\vspace*{\fill}
\begin{center}
{\Huge\bfseries DESIGN}\\[6mm]
\end{center}
\begin{center}
\begin{minipage}{0.92\textwidth}
\small
One of the most common issues creators face when visual thinking is speed of execution. The main reason for this is a lack of a system to help avoid falling in the trap of “pixel pushing.” 

I do two simple things that help me speed up my workflow and make it less likely I waste my design time on moving elements around:

1) Brainstorm on LARGE sheets paper (I use 11×17) or a journal (my notebook of choice if interested) \textbf{with a pen} because:\\
• No way to erase forces efficiency and an attitude of imperfection\\
• Having a large canvas allows for multiple iterations on the same page

2) Design on Canva to get a sense of where things should go. \\
• Canva can be a great free tool to test how things will look with colors and design\\
• Can speed up getting sizing and proportions of different elements without worrying about drawing 

Here’s a before and after of a cheeky version of the productivity jar meme I made this weekend.
\end{minipage}
\end{center}
\vspace{4mm}
\begin{center}
\includegraphics[width=\textwidth, height=0.24\textheight, keepaspectratio]{C:/GitHub/VisualIdeas/data/images/3f32fd57ed485e7fbe63e01388968b2e.png}
\end{center}
\vspace{2mm}
\begin{center}
\begin{minipage}{0.92\textwidth}
\small
Layout done in Canva
\end{minipage}
\end{center}
\vspace{5mm}
\begin{center}
\includegraphics[width=\textwidth, height=0.24\textheight, keepaspectratio]{C:/GitHub/VisualIdeas/data/images/413b01b7afb538ea5fdc14b5121b8ce1.png}
\end{center}
\vspace{2mm}
\begin{center}
\begin{minipage}{0.92\textwidth}
\small
Final version done in Procreate
\end{minipage}
\end{center}
\vspace*{\fill}
\newpage
\renewcommand{\theposttitle}{Growth Minded Celebrations and Disagreements}
\renewcommand{\thepostdate}{2023-07-04}
\vspace*{\fill}
\begin{center}
{\Huge\bfseries ENCHANTING}\\[6mm]
\end{center}
\begin{center}

\includegraphics[width=\textwidth, height=0.55\textheight, keepaspectratio]{C:/GitHub/VisualIdeas/data/images/95ce04bfd88b68f181fae2b8008884c4.png}
\end{center}
\vspace{2mm}
\begin{center}
\begin{minipage}{0.92\textwidth}
\small
Inspired by an excerpt from Sahil Bloom’s Curiosity Chronicle Newsletter:

“Your success is not determined by your achievement of the extraordinary, but by your ability to show up, day-in, day-out, and lay one new brick in your life's wall.” \\
-Sahil Bloom
\end{minipage}
\end{center}
\vspace*{\fill}
\newpage
\renewcommand{\theposttitle}{Growth Minded Celebrations and Disagreements}
\renewcommand{\thepostdate}{2023-07-04}
\vspace*{\fill}
\begin{center}
{\Huge\bfseries ANALOGY}\\[6mm]
\end{center}
\begin{center}
\includegraphics[width=\textwidth, height=0.55\textheight, keepaspectratio]{C:/GitHub/VisualIdeas/data/images/ec2d11376a50b5b95ed8fafa2a1248ea.png}
\end{center}
\vspace{2mm}
\begin{center}
\begin{minipage}{0.92\textwidth}
\small
Deadbolt Mindset vs Blue Sky Mindset

When I disagree with someone, I often assume I already know what they’re gonna say before they say it. This approach has rarely if ever helped me resolve an issue. 

So over the last several months, I’ve made a concerted effort to prompt myself to stay curious in arguments (helps a lot in relationship squabbles).

I simply do this through asking questions. What’s critical (and didn’t initially come natural to me), is avoiding leading or rhetorical questions. They’re only helpful if they’re earnest questions.

Login or Subscribe to participate
\end{minipage}
\end{center}
\vspace*{\fill}
\newpage
\renewcommand{\theposttitle}{Notemaking and the Stability of your Identity}
\renewcommand{\thepostdate}{2023-07-18}
\vspace*{\fill}
\begin{center}
{\Huge\bfseries INTERESTING}\\[6mm]
\end{center}
\begin{center}
\includegraphics[width=\textwidth, height=0.55\textheight, keepaspectratio]{C:/GitHub/VisualIdeas/data/images/32158f14c766b982f668efb4c6d0b086.png}
\end{center}
\vspace*{\fill}
\newpage
\renewcommand{\theposttitle}{Notemaking and the Stability of your Identity}
\renewcommand{\thepostdate}{2023-07-18}
\vspace*{\fill}
\begin{center}
{\Huge\bfseries DESIGN}\\[6mm]
\end{center}
\begin{center}
\begin{minipage}{0.92\textwidth}
\small
Behind-the-Scenes designing of \textit{“Run Your Own Race”} 👇
\end{minipage}
\end{center}
\vspace{4mm}
\begin{center}
\includegraphics[width=\textwidth, height=0.13\textheight, keepaspectratio]{C:/GitHub/VisualIdeas/data/images/7242a0028ed7042356b60fb43e859a7f.png}
\end{center}
\vspace{2mm}
\begin{center}
\begin{minipage}{0.92\textwidth}
\small
\textbf{Problem 1: }Text and messaging feel like they can be more coherent.

Let’s use \textbf{color} to connect the labels \textbf{“10 Years Ago”} and \textbf{“Today” }more closely with the \textbf{START sign}.
\end{minipage}
\end{center}
\vspace{5mm}
\begin{center}
\includegraphics[width=\textwidth, height=0.13\textheight, keepaspectratio]{C:/GitHub/VisualIdeas/data/images/f23378b86c6a4fa4934a6e67ad8214c4.jpg}
\end{center}
\vspace{5mm}
\begin{center}
\includegraphics[width=\textwidth, height=0.13\textheight, keepaspectratio]{C:/GitHub/VisualIdeas/data/images/669f3d6a74b4a898d2552096a8059b32.jpg}
\end{center}
\vspace{5mm}
\begin{center}
\includegraphics[width=\textwidth, height=0.13\textheight, keepaspectratio]{C:/GitHub/VisualIdeas/data/images/0e1a12675da4e543c48357e412059298.png}
\end{center}
\vspace{2mm}
\begin{center}
\begin{minipage}{0.92\textwidth}
\small
Not bad! It even makes more sense now because our hero should be at the absolute beginning \textbf{Today}. 🙌

\textit{(confession: you may notice it’s still not aligned perfectly…but it was late and I needed to publish! It’ll be better in a future iteration }😉\textit{)}
\end{minipage}
\end{center}
\vspace*{\fill}
\newpage
\renewcommand{\theposttitle}{Notemaking and the Stability of your Identity}
\renewcommand{\thepostdate}{2023-07-18}
\vspace*{\fill}
\begin{center}
{\Huge\bfseries ENCHANTING}\\[6mm]
\end{center}
\begin{center}
\includegraphics[width=\textwidth, height=0.55\textheight, keepaspectratio]{C:/GitHub/VisualIdeas/data/images/57397cf83afbd0839c40e2f4588528f7.png}
\end{center}
\vspace{2mm}
\begin{center}
\begin{minipage}{0.92\textwidth}
\small
Believing the world is broken is the fastest path (and easiest excuse) to descend into nihilism.

It keeps us from doing anything about it because…well, what’s the point.

I wonder though, what would the world look like if we fixed the lens we saw it through?
\end{minipage}
\end{center}
\vspace*{\fill}
\newpage
\renewcommand{\theposttitle}{Notemaking and the Stability of your Identity}
\renewcommand{\thepostdate}{2023-07-18}
\vspace*{\fill}
\begin{center}
{\Huge\bfseries ANALOGY}\\[6mm]
\end{center}
\begin{center}
\includegraphics[width=\textwidth, height=0.55\textheight, keepaspectratio]{C:/GitHub/VisualIdeas/data/images/e6065ef385b6fa7b015360f75677a46f.png}
\end{center}
\vspace{2mm}
\begin{center}
\begin{minipage}{0.92\textwidth}
\small
\textbf{Power of Notemaking}

I get asked a lot how I come up with ideas for visuals.

One method I use is "notemaking." It's pretty much notetaking on steroids.

Or another way to look at it: journaling on the ideas you consume.\\
Read the full post here.

Login or Subscribe to participate
\end{minipage}
\end{center}
\vspace*{\fill}
\newpage
\renewcommand{\theposttitle}{Setbacks and Sheeran's Law}
\renewcommand{\thepostdate}{2023-08-01}
\vspace*{\fill}
\begin{center}
{\Huge\bfseries INTERESTING}\\[6mm]
\end{center}
\begin{center}
\includegraphics[width=\textwidth, height=0.55\textheight, keepaspectratio]{C:/GitHub/VisualIdeas/data/images/b953a0872accb3dc6bb6f2a178a996bf.png}
\end{center}
\vspace{2mm}
\begin{center}
\begin{minipage}{0.92\textwidth}
\small
Sheeran's Law \\
("Rule of the Dirty Faucet)

A piece of advice I give my screenwriting students:\\
Write ALL your ideas down.

Especially the bad ones.

If you leave it, the gunk expands and muddies your thinking.

If you want to make the good stuff, you have to get through the bad stuff.

"When you switch on the dirty tap, sh*t water is going to flow out for a substantial amount of time. Then clean water is gonna start flowin'”\\
-Ed Sheeran
\end{minipage}
\end{center}
\vspace*{\fill}
\newpage
\renewcommand{\theposttitle}{Setbacks and Sheeran's Law}
\renewcommand{\thepostdate}{2023-08-01}
\vspace*{\fill}
\begin{center}
{\Huge\bfseries DESIGN}\\[6mm]
\end{center}
\begin{center}
\begin{minipage}{0.92\textwidth}
\small
There’s a lot that goes into \textbf{clarifying}, \textbf{simplifying} and \textbf{amplifying} an idea and turning into a visual, but ask any creator and they’ll tell you it all starts with writing down observations.

Sounds simple enough.

But one important principle most people forget:\\
\textbf{Delay judgement }until you’re in the execution phase.

One thing I tell my students all the time is \textit{“brainstorming is an idea generation process, not a decision making one.” }

Here’s one of my original notes:
\end{minipage}
\end{center}
\vspace{4mm}
\begin{center}
\includegraphics[width=\textwidth, height=0.24\textheight, keepaspectratio]{C:/GitHub/VisualIdeas/data/images/ba6e65ca8fbe885761902f18cf82fd47.jpg}
\end{center}
\vspace{2mm}
\begin{center}
\begin{minipage}{0.92\textwidth}
\small
Not pretty I know but I try to make it a point to draw an idea from my head in less than 10 seconds in notes. 

This observation later turned into this:
\end{minipage}
\end{center}
\vspace{5mm}
\begin{center}
\includegraphics[width=\textwidth, height=0.24\textheight, keepaspectratio]{C:/GitHub/VisualIdeas/data/images/ef540b5d7078d7a3bb2eccfa0a17d5b3.png}
\end{center}
\vspace{2mm}
\begin{center}
\begin{minipage}{0.92\textwidth}
\small
\textbf{Takeaways:}

Write all ideas and observations down without judging their merit.

Challenge yourself to getting it on the page in less than 10s.
\end{minipage}
\end{center}
\vspace*{\fill}
\newpage
\renewcommand{\theposttitle}{Setbacks and Sheeran's Law}
\renewcommand{\thepostdate}{2023-08-01}
\vspace*{\fill}
\begin{center}
{\Huge\bfseries ENCHANTING}\\[6mm]
\end{center}
\begin{center}
\includegraphics[width=\textwidth, height=0.55\textheight, keepaspectratio]{C:/GitHub/VisualIdeas/data/images/69c7efdb0223e6a9c5a23e04ad1c2086.png}
\end{center}
\vspace{2mm}
\begin{center}
\begin{minipage}{0.92\textwidth}
\small
"An arrow can only be shot by pulling it backward.

So when life is dragging you back with difficulties, it means that it's just going to launch you into something great.

So just focus, and keep aiming."

-Stephanie Simpson
\end{minipage}
\end{center}
\vspace*{\fill}
\newpage
\renewcommand{\theposttitle}{Setbacks and Sheeran's Law}
\renewcommand{\thepostdate}{2023-08-01}
\vspace*{\fill}
\begin{center}
{\Huge\bfseries ANALOGY}\\[6mm]
\end{center}
\begin{center}
\includegraphics[width=\textwidth, height=0.55\textheight, keepaspectratio]{C:/GitHub/VisualIdeas/data/images/3f7e88cd1ea1a94a1542a601750092ba.png}
\end{center}
\vspace{2mm}
\begin{center}
\begin{minipage}{0.92\textwidth}
\small
"Remarkable results are simply the result of unremarkable daily effort."\\
-Sahil Bloom

If you’ve never seen a stone cut this way, it’s actually quite amazing. 

In this video, a 26,000 pound block of granite is split into 2 equal parts with a 2 pound hammer:

Login or Subscribe to participate
\end{minipage}
\end{center}
\vspace*{\fill}
\newpage
\renewcommand{\theposttitle}{Maximizing vs Optimizing Your Day}
\renewcommand{\thepostdate}{2023-08-15}
\vspace*{\fill}
\begin{center}
{\Huge\bfseries INTERESTING}\\[6mm]
\end{center}
\begin{center}
\includegraphics[width=\textwidth, height=0.17\textheight, keepaspectratio]{C:/GitHub/VisualIdeas/data/images/022b6820907d52948c4d1b3d5d0b0960.png}
\end{center}
\vspace{2mm}
\begin{center}
\begin{minipage}{0.92\textwidth}
\small
Maximizing your schedule vs optimizing your schedule.

I first started thinking about this when I heard James Clear, author of “Atomic Habits,” speak on the Tim Ferris show.

He asks, when thinking about how to manage your time, \textbf{“what are you optimizing for?”}

Instead of trying to maximize your days, be a little more at peace with…well…living an optimized life unique for you.

If you’re interested in exploring these ideas further, I’d highly recommend these 2 books:

“4000 Weeks”
\end{minipage}
\end{center}
\vspace{5mm}
\begin{center}
\includegraphics[width=\textwidth, height=0.17\textheight, keepaspectratio]{C:/GitHub/VisualIdeas/data/images/d8798cfbb71e88c29c4c8f90ccd37c10.jpg}
\end{center}
\vspace{2mm}
\begin{center}
\begin{minipage}{0.92\textwidth}
\small
“From Strength to Strength”
\end{minipage}
\end{center}
\vspace{5mm}
\begin{center}
\includegraphics[width=\textwidth, height=0.17\textheight, keepaspectratio]{C:/GitHub/VisualIdeas/data/images/b9ba1a72fd2e638929a2cd017cc21ab9.jpg}
\end{center}
\vspace{2mm}
\begin{center}
\begin{minipage}{0.92\textwidth}
\small
(\textless{}= could not delete this for some reason 😅)
\end{minipage}
\end{center}
\vspace*{\fill}
\newpage
\renewcommand{\theposttitle}{Maximizing vs Optimizing Your Day}
\renewcommand{\thepostdate}{2023-08-15}
\vspace*{\fill}
\begin{center}
{\Huge\bfseries DESIGN}\\[6mm]
\end{center}
\begin{center}
\begin{minipage}{0.92\textwidth}
\small
A not-so-secret, simple hack designers use to speed up their workflow: \\
\textbf{GRIDS.} 

Grids help reliably frame and align graphics in a cohesive way.

There are lots of formats for grids, but here are simple ones I use most often:

\textbf{Horizontal 2-Shot }\\
(majority of my visuals uses this simple one)
\end{minipage}
\end{center}
\vspace{4mm}
\begin{center}

\includegraphics[width=\textwidth, height=0.13\textheight, keepaspectratio]{C:/GitHub/VisualIdeas/data/images/02617de01bb56b5e4cdbe0c2f8630a77.png}
\end{center}
\vspace{5mm}
\begin{center}

\includegraphics[width=\textwidth, height=0.13\textheight, keepaspectratio]{C:/GitHub/VisualIdeas/data/images/c4f44d23ede2767cd247a8d40bb926f1.png}
\end{center}
\vspace{5mm}
\begin{center}

\includegraphics[width=\textwidth, height=0.13\textheight, keepaspectratio]{C:/GitHub/VisualIdeas/data/images/5c6f220dcc459157954e6985f84eaed4.png}
\end{center}
\vspace{5mm}
\begin{center}
\includegraphics[width=\textwidth, height=0.13\textheight, keepaspectratio]{C:/GitHub/VisualIdeas/data/images/6adf867499e8022c8dd0d62fd782b4dd.png}
\end{center}
\vspace*{\fill}
\newpage
\renewcommand{\theposttitle}{Maximizing vs Optimizing Your Day}
\renewcommand{\thepostdate}{2023-08-15}
\vspace*{\fill}
\begin{center}
{\Huge\bfseries ENCHANTING}\\[6mm]
\end{center}
\begin{center}
\includegraphics[width=\textwidth, height=0.55\textheight, keepaspectratio]{C:/GitHub/VisualIdeas/data/images/6c1f59884d7ff9741b2f69ab58910ad9.png}
\end{center}
\vspace{2mm}
\begin{center}
\begin{minipage}{0.92\textwidth}
\small
\textbf{Friendly reminder: it's all relative.}
\end{minipage}
\end{center}
\vspace*{\fill}
\newpage
\renewcommand{\theposttitle}{Maximizing vs Optimizing Your Day}
\renewcommand{\thepostdate}{2023-08-15}
\vspace*{\fill}
\begin{center}
{\Huge\bfseries ANALOGY}\\[6mm]
\end{center}
\begin{center}
\includegraphics[width=\textwidth, height=0.55\textheight, keepaspectratio]{C:/GitHub/VisualIdeas/data/images/c98eb4a3a220732fc112b54b95fee5a3.png}
\end{center}
\vspace{2mm}
\begin{center}
\begin{minipage}{0.92\textwidth}
\small
\textbf{True definition of failure.}

Login or Subscribe to participate
\end{minipage}
\end{center}
\vspace*{\fill}
\newpage
\renewcommand{\theposttitle}{The Multitasking Myth and How Passion is Connected to Legacy}
\renewcommand{\thepostdate}{2023-08-29}
\vspace*{\fill}
\begin{center}
{\Huge\bfseries INTERESTING}\\[6mm]
\end{center}
\begin{center}
\includegraphics[width=\textwidth, height=0.55\textheight, keepaspectratio]{C:/GitHub/VisualIdeas/data/images/670288eb559b5ee6db052186071eea59.png}
\end{center}
\vspace{2mm}
\begin{center}
\begin{minipage}{0.92\textwidth}
\small
\textbf{Staying busy vs being productive.}

If you haven’t had a chance to read it yet, I’d highly recommend Cal Newport’s book, \textbf{“Deep Work Rules for Focused Success in a Distracted World.”} It’s changed the way I think about how to work on projects.
\end{minipage}
\end{center}
\vspace*{\fill}
\newpage
\renewcommand{\theposttitle}{The Multitasking Myth and How Passion is Connected to Legacy}
\renewcommand{\thepostdate}{2023-08-29}
\vspace*{\fill}
\begin{center}
{\Huge\bfseries DESIGN}\\[6mm]
\end{center}
\begin{center}
\begin{minipage}{0.92\textwidth}
\small
A question from a student \textbf{in my course}: \\
\textbf{How do you build a consistent looking brand from visual to visual? }

I wish I had a good hack…but it took me some time. Below are my first 50ish visuals I ever posted.

When I first started, I did my best to actually avoid worrying about style or consistency (as you might be able to tell from the images below). But a funny thing happened after doing it for a couple months…I started getting messages saying specifically how they liked my style! 🤯

And surprisingly, around mid-July 2022 (about 4 months and 50+ visuals after I started), you can see the style I use today emerging out of the pack.
\end{minipage}
\end{center}
\vspace{4mm}
\begin{center}
\includegraphics[width=\textwidth, height=0.13\textheight, keepaspectratio]{C:/GitHub/VisualIdeas/data/images/207e4a66ce09bd435e4244f91be4b33b.jpeg}
\end{center}
\vspace{5mm}
\begin{center}
\includegraphics[width=\textwidth, height=0.13\textheight, keepaspectratio]{C:/GitHub/VisualIdeas/data/images/e5781bcea64dc5808e70829bcdb27f58.jpeg}
\end{center}
\vspace{5mm}
\begin{center}
\includegraphics[width=\textwidth, height=0.13\textheight, keepaspectratio]{C:/GitHub/VisualIdeas/data/images/790e099f8e93a81f4889d7ff1f30ef8a.jpeg}
\end{center}
\vspace{5mm}
\begin{center}
\includegraphics[width=\textwidth, height=0.13\textheight, keepaspectratio]{C:/GitHub/VisualIdeas/data/images/9e3e74f1a336c86f6a8bc78749fb7aea.jpeg}
\end{center}
\vspace{5mm}
\begin{center}
\includegraphics[width=\textwidth, height=0.13\textheight, keepaspectratio]{C:/GitHub/VisualIdeas/data/images/858a1287c0c876dece0c25acd86315a3.jpeg}
\end{center}
\vspace{2mm}
\begin{center}
\begin{minipage}{0.92\textwidth}
\small
The lesson?

In the beginning, it’s more important to make stuff to see what you like and what’s working. Style inevitably emerges (whether you mean to or not 😅) and it evolves with your practice.

Style is more than just colors and fonts. It’s about voice and the kinds of messages you hope to portray.
\end{minipage}
\end{center}
\vspace*{\fill}
\newpage
\renewcommand{\theposttitle}{The Multitasking Myth and How Passion is Connected to Legacy}
\renewcommand{\thepostdate}{2023-08-29}
\vspace*{\fill}
\begin{center}
{\Huge\bfseries ENCHANTING}\\[6mm]
\end{center}
\begin{center}

\includegraphics[width=\textwidth, height=0.55\textheight, keepaspectratio]{C:/GitHub/VisualIdeas/data/images/8d21e9072dbdeb49a7ed266b04a87b5d.png}
\end{center}
\vspace{2mm}
\begin{center}
\begin{minipage}{0.92\textwidth}
\small
What you want to pass on to the people around you should drive your passion.
\end{minipage}
\end{center}
\vspace*{\fill}
\newpage
\renewcommand{\theposttitle}{The Multitasking Myth and How Passion is Connected to Legacy}
\renewcommand{\thepostdate}{2023-08-29}
\vspace*{\fill}
\begin{center}
{\Huge\bfseries ANALOGY}\\[6mm]
\end{center}
\begin{center}
\includegraphics[width=\textwidth, height=0.55\textheight, keepaspectratio]{C:/GitHub/VisualIdeas/data/images/548295fed295d35b22510cae3ecd872e.png}
\end{center}
\vspace{2mm}
\begin{center}
\begin{minipage}{0.92\textwidth}
\small
Tough love 🐢

Login or Subscribe to participate
\end{minipage}
\end{center}
\vspace*{\fill}
\newpage
\renewcommand{\theposttitle}{A system to capture ideas and organizing chaos with storytelling}
\renewcommand{\thepostdate}{2023-09-12}
\vspace*{\fill}
\begin{center}
{\Huge\bfseries INTERESTING}\\[6mm]
\end{center}
\begin{center}
\includegraphics[width=\textwidth, height=0.55\textheight, keepaspectratio]{C:/GitHub/VisualIdeas/data/images/cb8677d622a81faf7efe5b2367569ec5.png}
\end{center}
\vspace{2mm}
\begin{center}
\begin{minipage}{0.92\textwidth}
\small
Inspired by Tiago Forte and \textbf{his book, “Building a Second Brain.”}

A \textit{second brain} is a fancy word for an intentional notetaking system that keeps tracks of your ideas, projects, and other efforts in your life.

I have a bit of a convoluted process but it works for me:

\textbf{Quick Capture:} Drafts or Apple Notes (fleeting ideas in day)

\textbf{Slow Capture:} Notebook \& Pen (journal \& sketch)

\textbf{Long Capture:} Obsidian (examining \& colliding ideas)

The most important thing is not 'how', but 'how long.'

Avoid the temptation to switch systems once you’ve made a couple decisions. Stick to a system long enough to work out the kinks so it works for you.

There's no perfect system. It just becomes less imperfect over time. ❤️
\end{minipage}
\end{center}
\vspace*{\fill}
\newpage
\renewcommand{\theposttitle}{A system to capture ideas and organizing chaos with storytelling}
\renewcommand{\thepostdate}{2023-09-12}
\vspace*{\fill}
\begin{center}
{\Huge\bfseries DESIGN}\\[6mm]
\end{center}
\begin{center}
\begin{minipage}{0.92\textwidth}
\small
A simple tip for communicating ideas visually: \\
Use ❌ and ✅ to communicate your point of view.

What’s the number one thing your visual metaphor should answer?

That’s right!

The \textbf{Call to Action}!

In my \textbf{maximizing vs optimizing visual}\textbf{, }I present a difference…but then…

…what am I asking the reader to consider once they’ve seen this distinction?

That’s where I can say ❌ \textbf{(avoid this)}, but instead ✅ \textbf{(do this)}!

How cool is that?! 🤓
\end{minipage}
\end{center}
\vspace{4mm}
\begin{center}
\includegraphics[width=\textwidth, height=0.55\textheight, keepaspectratio]{C:/GitHub/VisualIdeas/data/images/7bb458ec3f6653cb7cb47fdec6f1055f.png}
\end{center}
\vspace{2mm}
\begin{center}
\begin{minipage}{0.92\textwidth}
\small
Before
\end{minipage}
\end{center}
\vspace*{\fill}
\newpage
\renewcommand{\theposttitle}{A system to capture ideas and organizing chaos with storytelling}
\renewcommand{\thepostdate}{2023-09-12}
\vspace*{\fill}
\begin{center}
{\Huge\bfseries ANALOGY}\\[6mm]
\end{center}
\begin{center}
\includegraphics[width=\textwidth, height=0.55\textheight, keepaspectratio]{C:/GitHub/VisualIdeas/data/images/0823c7ce1446d9148c87f68d276d7c74.png}
\end{center}
\vspace{2mm}
\begin{center}
\begin{minipage}{0.92\textwidth}
\small
In my 20th year teaching film and animation…

Wanna know the most difficult challenge my students face?

Learning to write economically. Saying more with less.

The process starts with throwing up everything onto the first draft (in screenwriting, we call this the “vomit script”…I restrained myself from visualizing this 🤣)

Then, compress as much as possible and squeeze out any excess crud!

Login or Subscribe to participate
\end{minipage}
\end{center}
\vspace*{\fill}
\newpage
\renewcommand{\theposttitle}{The difference between thinking and overthinking}
\renewcommand{\thepostdate}{2023-09-26}
\vspace*{\fill}
\begin{center}
{\Huge\bfseries INTERESTING}\\[6mm]
\end{center}
\begin{center}
\includegraphics[width=\textwidth, height=0.55\textheight, keepaspectratio]{C:/GitHub/VisualIdeas/data/images/494044774bd84df13d79269d4eee358a.png}
\end{center}
\vspace{2mm}
\begin{center}
\begin{minipage}{0.92\textwidth}
\small
The difference between thinking and overthinking. Stay sharp!
\end{minipage}
\end{center}
\vspace*{\fill}
\newpage
\renewcommand{\theposttitle}{The difference between thinking and overthinking}
\renewcommand{\thepostdate}{2023-09-26}
\vspace*{\fill}
\begin{center}
{\Huge\bfseries DESIGN}\\[6mm]
\end{center}
\begin{center}
\begin{minipage}{0.92\textwidth}
\small
If you don’t already have a sketchbook, consider starting one. 

\textbf{PRO TIP:} Use the same notebook to sketch and journal. Written insights inspire visual insights and vice versa.

\textbf{Law of Ideation}

\textit{Take your idea and make it real, no matter how rough. The final step is only attainable by taking the first step. }

One of my favorite creative practices is sketching visual strategies.

\textbf{What’s a visual strategy?} \\
It’s how an artist solves an idea problem with a visual solution.

For instance, the visual strategy below from the talented Wang Wei-Shen shows a possible solution to the problem: what happens when you cross a threshold.
\end{minipage}
\end{center}
\vspace{4mm}
\begin{center}
\includegraphics[width=\textwidth, height=0.17\textheight, keepaspectratio]{C:/GitHub/VisualIdeas/data/images/e4501303039c34891f93620af8c6b57e.jpg}
\end{center}
\vspace{2mm}
\begin{center}
\begin{minipage}{0.92\textwidth}
\small
Initial capture in my notebook. You might be able to see echoes of my journaling on the other side of the page.

This sketch ended up motivating the visual below on the left, which then iterated to the version on the right.
\end{minipage}
\end{center}
\vspace{5mm}
\begin{center}
\includegraphics[width=\textwidth, height=0.17\textheight, keepaspectratio]{C:/GitHub/VisualIdeas/data/images/2aadc144dbb14228079e7cd74dd212ae.png}
\end{center}
\vspace{5mm}
\begin{center}
\includegraphics[width=\textwidth, height=0.17\textheight, keepaspectratio]{C:/GitHub/VisualIdeas/data/images/2f4d836a93e1ae9b0b465738675a3ff5.png}
\end{center}
\vspace{2mm}
\begin{center}
\begin{minipage}{0.92\textwidth}
\small
But you know what’s funny?

I only realized this sketch was the source of inspiration months later. It burned into my subconscious image bank without my knowing. That’s the power of a strong capture habit.

Creative output is fueled by creative input.
\end{minipage}
\end{center}
\vspace*{\fill}
\newpage
\renewcommand{\theposttitle}{The difference between thinking and overthinking}
\renewcommand{\thepostdate}{2023-09-26}
\vspace*{\fill}
\begin{center}
{\Huge\bfseries ENCHANTING}\\[6mm]
\end{center}
\begin{center}
\includegraphics[width=\textwidth, height=0.55\textheight, keepaspectratio]{C:/GitHub/VisualIdeas/data/images/6b362b893fd5ee44eb19af82ec4fc9d9.png}
\end{center}
\vspace{2mm}
\begin{center}
\begin{minipage}{0.92\textwidth}
\small
Spend your time making impacts that last.
\end{minipage}
\end{center}
\vspace*{\fill}
\newpage
\renewcommand{\theposttitle}{The difference between thinking and overthinking}
\renewcommand{\thepostdate}{2023-09-26}
\vspace*{\fill}
\begin{center}
{\Huge\bfseries ANALOGY}\\[6mm]
\end{center}
\begin{center}
\includegraphics[width=\textwidth, height=0.13\textheight, keepaspectratio]{C:/GitHub/VisualIdeas/data/images/2e9ce9df999424ac058af2a48688e854.png}
\end{center}
\vspace{5mm}
\begin{center}

\includegraphics[width=\textwidth, height=0.13\textheight, keepaspectratio]{C:/GitHub/VisualIdeas/data/images/371d9a62cd1673f63833193d2345dcb8.png}
\end{center}
\vspace{5mm}
\begin{center}

\includegraphics[width=\textwidth, height=0.13\textheight, keepaspectratio]{C:/GitHub/VisualIdeas/data/images/5647859d2e86cbf8a3b0b1151780b06d.png}
\end{center}
\vspace{5mm}
\begin{center}
\includegraphics[width=\textwidth, height=0.13\textheight, keepaspectratio]{C:/GitHub/VisualIdeas/data/images/9236c6492fa4eb12467fed6711957866.png}
\end{center}
\vspace{2mm}
\begin{center}
\begin{minipage}{0.92\textwidth}
\small
\textbf{The Lines in Sand Effect:}\\
A simple mental model to get past overthinking and the fear of making decisions.

Any time I’ve had to make a choice where I recognize I’m procrastinating because of the fear of making a mistake, I imagine this metaphor.

I’ve done this to decide:

niches

type of clients I serve

tech platforms

even habits \& routines

It comforts me to remind myself these decisions aren’t set it stone.

Login or Subscribe to participate
\end{minipage}
\end{center}
\vspace*{\fill}
\newpage
\renewcommand{\theposttitle}{2 Types of Burnout}
\renewcommand{\thepostdate}{2023-10-10}
\vspace*{\fill}
\begin{center}
{\Huge\bfseries INTERESTING}\\[6mm]
\end{center}
\begin{center}
\includegraphics[width=\textwidth, height=0.55\textheight, keepaspectratio]{C:/GitHub/VisualIdeas/data/images/d3c342fc707ffcfa87e63aef80783852.png}
\end{center}
\vspace{2mm}
\begin{center}
\begin{minipage}{0.92\textwidth}
\small
Burning out can happen at 2 intensities:

Most people understand \textbf{“fast burnout”}—when you push yourself beyond capacity towards a goal, sacrificing your physical and mental health.

But equally dangerous is \textbf{“slow burnout”}—when you feel trapped in a situation and your motivation to do anything slowly suffocates. 

We often suffer silently from a misplaced belief that it's all part of "enduring the grind."

To dive deeper, my friend Roberto Ferraro shared this link:

\textbf{Harvard Business Review: }\\
"Is Your Burnout From Too Much Work or Too Little Impact?"
\end{minipage}
\end{center}
\vspace*{\fill}
\newpage
\renewcommand{\theposttitle}{2 Types of Burnout}
\renewcommand{\thepostdate}{2023-10-10}
\vspace*{\fill}
\begin{center}
{\Huge\bfseries DESIGN}\\[6mm]
\end{center}
\begin{center}
\begin{minipage}{0.92\textwidth}
\small
Most people misunderstand the term “\textbf{steal like an artist.}” They just “steal” and forget the key part of the phrase: \textbf{“like an artist.”}

\textbf{Artists combine.}

An audience should be able to see influence but not a direct carbon copy. \textbf{ }

As Austin Kleon, the author of \textbf{Steal Like an Artist }notes in his book:

“If you steal from one author, it’s plagiarism; if you steal from many, it’s research.”

Here’s a string of influences as an example:
\end{minipage}
\end{center}
\vspace{4mm}
\begin{center}
\includegraphics[width=\textwidth, height=0.17\textheight, keepaspectratio]{C:/GitHub/VisualIdeas/data/images/1a29501342754e2a5ab8755d58f6cdd2.png}
\end{center}
\vspace{2mm}
\begin{center}
\begin{minipage}{0.92\textwidth}
\small
The above image was influenced by my friend Roberto Ferraro who was originally inspired by combining my illustration of Russian Dolls and an idea from Jeff Mckeon. You can see the evolution of images and thinking.

That’s the power of ‘stealing like an artist.’
\end{minipage}
\end{center}
\vspace{5mm}
\begin{center}
\includegraphics[width=\textwidth, height=0.17\textheight, keepaspectratio]{C:/GitHub/VisualIdeas/data/images/15429511ba767a3d315fbf7b5d99d6a2.jpeg}
\end{center}
\vspace{2mm}
\begin{center}
\begin{minipage}{0.92\textwidth}
\small
Was inspired by my friend Roberto Ferraro who was inspired by 👉
\end{minipage}
\end{center}
\vspace{5mm}
\begin{center}
\includegraphics[width=\textwidth, height=0.17\textheight, keepaspectratio]{C:/GitHub/VisualIdeas/data/images/d73592712824adbf33478ffed441400e.png}
\end{center}
\vspace{2mm}
\begin{center}
\begin{minipage}{0.92\textwidth}
\small
The image I made that inspired Roberto
\end{minipage}
\end{center}
\vspace*{\fill}
\newpage
\renewcommand{\theposttitle}{2 Types of Burnout}
\renewcommand{\thepostdate}{2023-10-10}
\vspace*{\fill}
\begin{center}
{\Huge\bfseries ENCHANTING}\\[6mm]
\end{center}
\begin{center}
\includegraphics[width=\textwidth, height=0.55\textheight, keepaspectratio]{C:/GitHub/VisualIdeas/data/images/680c5d0c4f89b9d6f94e26019d957778.png}
\end{center}
\vspace{2mm}
\begin{center}
\begin{minipage}{0.92\textwidth}
\small
One brick at a time.

\textbf{Sidenote: }Equally important is to enjoy your chapter 20 once you get to that part of your story. There’s a time to build and a time to appreciate what you’ve built.
\end{minipage}
\end{center}
\vspace*{\fill}
\newpage
\renewcommand{\theposttitle}{2 Types of Burnout}
\renewcommand{\thepostdate}{2023-10-10}
\vspace*{\fill}
\begin{center}
{\Huge\bfseries ANALOGY}\\[6mm]
\end{center}
\begin{center}

\includegraphics[width=\textwidth, height=0.55\textheight, keepaspectratio]{C:/GitHub/VisualIdeas/data/images/8af2dfe8f8ae8cdccca6de9e490c0ecf.png}
\end{center}
\vspace{2mm}
\begin{center}
\begin{minipage}{0.92\textwidth}
\small
\textbf{Every book is an open invitation to make a transformation.}

Login or Subscribe to participate
\end{minipage}
\end{center}
\vspace*{\fill}
\newpage
\renewcommand{\theposttitle}{In the commitment game, speed is your ally (but not in the way you think)}
\renewcommand{\thepostdate}{2023-10-24}
\vspace*{\fill}
\begin{center}
{\Huge\bfseries INTERESTING}\\[6mm]
\end{center}
\begin{center}
\includegraphics[width=\textwidth, height=0.55\textheight, keepaspectratio]{C:/GitHub/VisualIdeas/data/images/7996382e97a6768dab599a668cb33679.png}
\end{center}
\vspace{2mm}
\begin{center}
\begin{minipage}{0.92\textwidth}
\small
"We need to learn the slow 'yes' and the quick 'no.'"\\
-Greg McKeown, author of “Essentialism: The Disciplined Pursuit of Less”

I recognize the word ‘no’ is sometimes a luxury most people in their jobs can’t afford. But we can ask the person making the request to balance out the equation: if we’re going to START doing this, what are we going to STOP doing?
\end{minipage}
\end{center}
\vspace*{\fill}
\newpage
\renewcommand{\theposttitle}{In the commitment game, speed is your ally (but not in the way you think)}
\renewcommand{\thepostdate}{2023-10-24}
\vspace*{\fill}
\begin{center}
{\Huge\bfseries DESIGN}\\[6mm]
\end{center}
\begin{center}
\begin{minipage}{0.92\textwidth}
\small
\textbf{The Leaky Bucket Metaphor Framework}
\end{minipage}
\end{center}
\vspace{4mm}
\begin{center}
\includegraphics[width=\textwidth, height=0.55\textheight, keepaspectratio]{C:/GitHub/VisualIdeas/data/images/30dc79f4e9f8526aa93e9a3eac4f23e4.png}
\end{center}
\vspace{2mm}
\begin{center}
\begin{minipage}{0.92\textwidth}
\small
You can use the Leaky Bucket visual metaphor as a “correction insight”—where you highlight a behavior that may be going overlooked and is preventing a positive outcome.
\end{minipage}
\end{center}
\vspace*{\fill}
\newpage
\renewcommand{\theposttitle}{In the commitment game, speed is your ally (but not in the way you think)}
\renewcommand{\thepostdate}{2023-10-24}
\vspace*{\fill}
\begin{center}
{\Huge\bfseries ENCHANTING}\\[6mm]
\end{center}
\begin{center}
\includegraphics[width=\textwidth, height=0.55\textheight, keepaspectratio]{C:/GitHub/VisualIdeas/data/images/859d58fff569bda38051a26c447ed423.png}
\end{center}
\vspace{2mm}
\begin{center}
\begin{minipage}{0.92\textwidth}
\small
\textbf{Success comes in many shapes. }

The older I get, the more I realize how important it is to savor each daily bite of life.
\end{minipage}
\end{center}
\vspace*{\fill}
\newpage
\renewcommand{\theposttitle}{In the commitment game, speed is your ally (but not in the way you think)}
\renewcommand{\thepostdate}{2023-10-24}
\vspace*{\fill}
\begin{center}
{\Huge\bfseries ANALOGY}\\[6mm]
\end{center}
\begin{center}
\includegraphics[width=\textwidth, height=0.55\textheight, keepaspectratio]{C:/GitHub/VisualIdeas/data/images/be1c67c53489565dffbae712c43c679f.png}
\end{center}
\vspace{2mm}
\begin{center}
\begin{minipage}{0.92\textwidth}
\small
\textbf{Rocking Horse Syndrome}: Moving a lot, but not progressing. \\
(inspired by Justin Welsh)

Before I had my first kid, I remember reading piles of parenting books \& articles to get prepared for what I knew was gonna be D-Day in the Milani household. (“D” standing for diapers of course 😆)

And you know what? I don’t remember a single piece of advice I read. 

Not a single one.

You know why? They all contradicted each other.

Turns out the only true way to learn parenting is by being a parent. Only then does the advice have a place to be relevant.

It’s the same with business (or probably any skill). Prep all you want, but it’s like using a bike manual to study how to ride a bike. The hope being the first time you get on the bike, you don’t fall.

But everybody falls! In bikes, in parenting, in business.

Time to get on a real horse. Let’s ride.

Login or Subscribe to participate
\end{minipage}
\end{center}
\vspace*{\fill}
\newpage
\renewcommand{\theposttitle}{Teaching, Breakdowns, True Wealth, \& Insights}
\renewcommand{\thepostdate}{2023-11-07}
\vspace*{\fill}
\begin{center}
{\Huge\bfseries INTERESTING}\\[6mm]
\end{center}
\begin{center}
\includegraphics[width=\textwidth, height=0.55\textheight, keepaspectratio]{C:/GitHub/VisualIdeas/data/images/d5e3615459bdbfc8807ca218c5f8574d.png}
\end{center}
\vspace{2mm}
\begin{center}
\begin{minipage}{0.92\textwidth}
\small
\textbf{The mediocre teacher tells. }\\
\textbf{The good teacher explains. }\\
\textbf{The superior teacher demonstrates. }\\
\textbf{The great teacher inspires.}\\
-\textbf{ }William Arthur Ward

This illustration about teaching sums up both my teaching philosophy and my creative philosophy.

Instead of feeding someone what to think, I want to leave space to invite them to participate in the meaning.

I’m conscious that makes my work not for everyone…but I hope the people who it is for appreciate the space. 😊
\end{minipage}
\end{center}
\vspace*{\fill}
\newpage
\renewcommand{\theposttitle}{Teaching, Breakdowns, True Wealth, \& Insights}
\renewcommand{\thepostdate}{2023-11-07}
\vspace*{\fill}
\begin{center}
{\Huge\bfseries DESIGN}\\[6mm]
\end{center}
\begin{center}
\begin{minipage}{0.92\textwidth}
\small
Once you understand the relationship you’re presenting in a metaphor, you open the door for your imagination to create substitutions with the same concept.

For example, the iceberg metaphor is considered the most commonly used visual metaphor with the following relationship:
\end{minipage}
\end{center}
\vspace{4mm}
\begin{center}
\includegraphics[width=\textwidth, height=0.24\textheight, keepaspectratio]{C:/GitHub/VisualIdeas/data/images/61c404c5e46dbc29a8bfe65ab06f98ca.png}
\end{center}
\vspace{2mm}
\begin{center}
\begin{minipage}{0.92\textwidth}
\small
Standard use-case for an iceberg metaphor

Once we understand the above surface vs below surface relationship, we can present it in a different context like this:
\end{minipage}
\end{center}
\vspace{5mm}
\begin{center}
\includegraphics[width=\textwidth, height=0.24\textheight, keepaspectratio]{C:/GitHub/VisualIdeas/data/images/7fbb3f280fc0b100379774bb53d82efd.png}
\end{center}
\vspace{2mm}
\begin{center}
\begin{minipage}{0.92\textwidth}
\small
Can you think of a metaphor using this concept? \\
Shoot me a reply if you do. I’d love to hear it! 

And if you wanna learn more about how to create visual metaphors, \textbf{check out my course, Thinking in Visual Metaphors. }
\end{minipage}
\end{center}
\vspace*{\fill}
\newpage
\renewcommand{\theposttitle}{Teaching, Breakdowns, True Wealth, \& Insights}
\renewcommand{\thepostdate}{2023-11-07}
\vspace*{\fill}
\begin{center}
{\Huge\bfseries ENCHANTING}\\[6mm]
\end{center}
\begin{center}
\includegraphics[width=\textwidth, height=0.55\textheight, keepaspectratio]{C:/GitHub/VisualIdeas/data/images/4e5298c0b0f65953688f782ed6873a53.png}
\end{center}
\vspace{2mm}
\begin{center}
\begin{minipage}{0.92\textwidth}
\small
\textbf{Where breakthroughs often come from.}

If I look back at the biggest changes I’ve made in my life, they often took root in times of darkness.

It’s as if you grow more sensitive to the light, because your eyes had to adjust in the dark. It’s unpleasant to be sure…but it’s also no fun staying alone in the dark.
\end{minipage}
\end{center}
\vspace*{\fill}
\newpage
\renewcommand{\theposttitle}{Teaching, Breakdowns, True Wealth, \& Insights}
\renewcommand{\thepostdate}{2023-11-07}
\vspace*{\fill}
\begin{center}
{\Huge\bfseries ANALOGY}\\[6mm]
\end{center}
\begin{center}
\includegraphics[width=\textwidth, height=0.55\textheight, keepaspectratio]{C:/GitHub/VisualIdeas/data/images/56c55ae390f7b98ec0631d35bf2a54c1.png}
\end{center}
\vspace{2mm}
\begin{center}
\begin{minipage}{0.92\textwidth}
\small
\textbf{There's more than one way to invest in your future.}

I used to think a couple thousand bucks was crazy to pay for some online course…but then I remembered how much it cost to go to college. 😆

And investing time in leveling up your skills is one of the best ways to compound both your wealth of knowledge and your earning potential.
\end{minipage}
\end{center}
\vspace*{\fill}
\newpage
\renewcommand{\theposttitle}{Habit tracking, comparisons, and why accuracy is overrated}
\renewcommand{\thepostdate}{2023-11-21}
\vspace*{\fill}
\begin{center}
{\Huge\bfseries INTERESTING}\\[6mm]
\end{center}
\begin{center}
\includegraphics[width=\textwidth, height=0.55\textheight, keepaspectratio]{C:/GitHub/VisualIdeas/data/images/636ecfc27c95a03bbba72e20f8fef0e1.png}
\end{center}
\vspace{2mm}
\begin{center}
\begin{minipage}{0.92\textwidth}
\small
You've read Atomic Habits. You're convinced.

Systems over goals. Habits over motivation.

Bright eyed, full of hope, you decide tomorrow you're getting it together. And miracle of miracles, you do!

You go by the book. You nail your work out! You think, "why the hell did I wait so long to start?!"

You open your calendar and mark a big X over the day. This must be what Seinfeld was talking about: “Don't break the chain!”

Easy enough.

But 6 days in...you hit a brick wall:\\
• you get sick\\
• you work late\\
• you live a normal day of life when the universe conspires against your best intentions. 

\textbf{Read full post}
\end{minipage}
\end{center}
\vspace*{\fill}
\newpage
\renewcommand{\theposttitle}{Habit tracking, comparisons, and why accuracy is overrated}
\renewcommand{\thepostdate}{2023-11-21}
\vspace*{\fill}
\begin{center}
{\Huge\bfseries DESIGN}\\[6mm]
\end{center}
\begin{center}
\begin{minipage}{0.92\textwidth}
\small
\textbf{How a simple color change opens the door for your audience to step into imagination}

My kids loved coloring Marvel superheroes growing up. But what’s funny is my son always needed to know EXACTLY what color to use for each character.

Red, white, \& blue for Captain America. \\
Red \& yellow for Iron Man. \\
But when I colored with him, I would use every color of the rainbow.

And of course he’d roll his eyes and say I was doing it wrong.

But I told him I’m just being creative. And he’d stare at me a moment…and just repeat I was wrong 😑 (some lessons take time 😆)

When you express yourself in art, it’s less important to make it accurate and more important to \textbf{make it interesting.}

Make your work unique to you. This can take shape simply in the colors you choose.

Turtles are green. \\
Basketballs are orange. \\
We’ve seen this accuracy in colors thousands of times.

But \textbf{change what your audience expects} and you earn a moment of pause. And pauses are valuable in a noisy world.
\end{minipage}
\end{center}
\vspace{4mm}
\begin{center}
\includegraphics[width=\textwidth, height=0.24\textheight, keepaspectratio]{C:/GitHub/VisualIdeas/data/images/1afb25b6fde30ffd46e0a22b7004e1a3.png}
\end{center}
\vspace{2mm}
\begin{center}
\begin{minipage}{0.92\textwidth}
\small
The goal is not to portray reality. We can show accurate colors with other mediums and tools.
\end{minipage}
\end{center}
\vspace{5mm}
\begin{center}
\includegraphics[width=\textwidth, height=0.24\textheight, keepaspectratio]{C:/GitHub/VisualIdeas/data/images/52174a5b10a1c0d100c7586b6738d5b9.png}
\end{center}
\vspace{2mm}
\begin{center}
\begin{minipage}{0.92\textwidth}
\small
The goal is to share a sense that this is coming from an imagination. Even a simple color change is a start.
\end{minipage}
\end{center}
\vspace*{\fill}
\newpage
\renewcommand{\theposttitle}{Habit tracking, comparisons, and why accuracy is overrated}
\renewcommand{\thepostdate}{2023-11-21}
\vspace*{\fill}
\begin{center}
{\Huge\bfseries ENCHANTING}\\[6mm]
\end{center}
\begin{center}
\includegraphics[width=\textwidth, height=0.55\textheight, keepaspectratio]{C:/GitHub/VisualIdeas/data/images/ba4b39f9843518fbc0a74184b308b1c7.png}
\end{center}
\vspace{2mm}
\begin{center}
\begin{minipage}{0.92\textwidth}
\small
\textbf{A flower does not think of competing with the flower next to it. }\\
\textbf{It just blooms.}\\
– Zen Proverb
\end{minipage}
\end{center}
\vspace*{\fill}
\newpage
\renewcommand{\theposttitle}{Habit tracking, comparisons, and why accuracy is overrated}
\renewcommand{\thepostdate}{2023-11-21}
\vspace*{\fill}
\begin{center}
{\Huge\bfseries ANALOGY}\\[6mm]
\end{center}
\begin{center}
\includegraphics[width=\textwidth, height=0.55\textheight, keepaspectratio]{C:/GitHub/VisualIdeas/data/images/d8a8338874c51a0dce49431b2e13d9a2.png}
\end{center}
\vspace{2mm}
\begin{center}
\begin{minipage}{0.92\textwidth}
\small
\textbf{"In nature, nothing is perfect and everything is perfect. Trees can be contorted, bent in weird ways, and they're still beautiful."}\\
-Alice Walker

Login or Subscribe to participate
\end{minipage}
\end{center}
\vspace*{\fill}
\newpage
\renewcommand{\theposttitle}{When you feel lost, think about where you're going}
\renewcommand{\thepostdate}{2023-12-05}
\vspace*{\fill}
\begin{center}
{\Huge\bfseries INTERESTING}\\[6mm]
\end{center}
\begin{center}
\includegraphics[width=\textwidth, height=0.55\textheight, keepaspectratio]{C:/GitHub/VisualIdeas/data/images/2bcf422a2647fc7b577ce853437d6c34.png}
\end{center}
\vspace{2mm}
\begin{center}
\begin{minipage}{0.92\textwidth}
\small
\textbf{"Things which matter most must never be at the mercy of things which matter least." }\\
–Johann Wolfgang von Goethe

One of my biggest “aha” moments came from the book, \textbf{“Essentialism” by Greg Mckeown}.

And the insight was simple: \textbf{the word priority is singular. }

Priority used to mean \textit{“the one thing that’s most important.”}

There was no such thing as priorit\textbf{ies} (plural)!

These days, we all want to do everything, everywhere, all at once. But that’s why we feel spread thin…why we’re everywhere and nowhere!

And personally, I know I often spend far too much time focusing on areas I should never have been in the first place!

Which reminds me of another favorite warning when it comes to where to focus your time:

\textbf{“There is nothing quite so useless, as doing with great efficiency, something that should not be done at all.”}
\end{minipage}
\end{center}
\vspace*{\fill}
\newpage
\renewcommand{\theposttitle}{When you feel lost, think about where you're going}
\renewcommand{\thepostdate}{2023-12-05}
\vspace*{\fill}
\begin{center}
{\Huge\bfseries DESIGN}\\[6mm]
\end{center}
\begin{center}
\begin{minipage}{0.92\textwidth}
\small
It’s common to get stuck when you’re brainstorming how to visualize an idea. That’s why I have a set of reliable frameworks.

For example, consider the following:
\end{minipage}
\end{center}
\vspace{4mm}
\begin{center}
\includegraphics[width=\textwidth, height=0.24\textheight, keepaspectratio]{C:/GitHub/VisualIdeas/data/images/0458fc1f74df6b8619f5ce614f655f5e.png}
\end{center}
\vspace{2mm}
\begin{center}
\begin{minipage}{0.92\textwidth}
\small
Inspired by the famous quote from Maya Angelou
\end{minipage}
\end{center}
\vspace{5mm}
\begin{center}
\includegraphics[width=\textwidth, height=0.24\textheight, keepaspectratio]{C:/GitHub/VisualIdeas/data/images/aca5fa69baacc0e4f42e18cd70f09951.png}
\end{center}
\vspace{2mm}
\begin{center}
\begin{minipage}{0.92\textwidth}
\small
This framework is inspired by fellow visual creator, Janis Ozolins.

\textbf{Frameworks are thinking vehicles} to get you from point A to point B.

And the more vehicles you know how to drive, the more likely you can travel anywhere you want to go with your messages.

If you want to learn more frameworks like this, you might be interested in the next \textbf{Thinking in Visual Metaphors course}.
\end{minipage}
\end{center}
\vspace*{\fill}
\newpage
\renewcommand{\theposttitle}{When you feel lost, think about where you're going}
\renewcommand{\thepostdate}{2023-12-05}
\vspace*{\fill}
\begin{center}
{\Huge\bfseries ENCHANTING}\\[6mm]
\end{center}
\begin{center}

\includegraphics[width=\textwidth, height=0.55\textheight, keepaspectratio]{C:/GitHub/VisualIdeas/data/images/e1465b6cdbb5aa73f8e1521561ad2869.png}
\end{center}
\vspace{2mm}
\begin{center}
\begin{minipage}{0.92\textwidth}
\small
\textbf{“The art of teaching is the art of assisting discovery.”}\\
– Mark Van Doren
\end{minipage}
\end{center}
\vspace*{\fill}
\newpage
\renewcommand{\theposttitle}{When you feel lost, think about where you're going}
\renewcommand{\thepostdate}{2023-12-05}
\vspace*{\fill}
\begin{center}
{\Huge\bfseries ANALOGY}\\[6mm]
\end{center}
\begin{center}
\includegraphics[width=\textwidth, height=0.55\textheight, keepaspectratio]{C:/GitHub/VisualIdeas/data/images/024f81ca093ec5da5f08264ff4fae777.png}
\end{center}
\vspace{2mm}
\begin{center}
\begin{minipage}{0.92\textwidth}
\small
\textbf{Just when you thought you figured it out.}

I loved Rubik’s Cubes as a kid, but honestly I never fully solved one. 😆

Best I could do was solve a single side and mess up everything else in the process.

This typically is my experience when I finally get work stuff in order—the rest of life turns into chaos (health, family, relationships).

P.S. My son can solve a Rubik’s Cube in less than 26 seconds…there’s a lesson in there somewhere. 🤔

Login or Subscribe to participate
\end{minipage}
\end{center}
\vspace*{\fill}
\newpage
\renewcommand{\theposttitle}{How expertise hurts your creativity and the ancient Zen Buddhist mindset that changes it}
\renewcommand{\thepostdate}{2023-12-19}
\vspace*{\fill}
\begin{center}
{\Huge\bfseries ENCHANTING}\\[6mm]
\end{center}
\begin{center}
\includegraphics[width=\textwidth, height=0.55\textheight, keepaspectratio]{C:/GitHub/VisualIdeas/data/images/c9b024bb932aaba345699a745ef657f6.png}
\end{center}
\vspace{2mm}
\begin{center}
\begin{minipage}{0.92\textwidth}
\small
\textbf{The definition of integrity. }

One of my favorite things to journal on is defining what a word means to me.

I started doing this after studying for the GRE (a test students have to take in the states to go to graduate school).

It made me realize there were lots of words I knew how to use in a sentence, but could only explain vaguely.

For example, take a simple word like “love”.

If you’ve ever read \textbf{The Five Love Languages}, then you know people have different (often subconscious) ways they define love.

And if you define what love means to you and see how it aligns with your partner’s definition, it’s pretty helpful to communicate more effectively in your relationship.

In the same vein, I’ve found clarifying how you understand other vague words (like integrity) helps align a deeper relationship with yourself.

I’m curious, is your definition of integrity the same or different?
\end{minipage}
\end{center}
\vspace*{\fill}
\newpage
\renewcommand{\theposttitle}{How expertise hurts your creativity and the ancient Zen Buddhist mindset that changes it}
\renewcommand{\thepostdate}{2023-12-19}
\vspace*{\fill}
\begin{center}
{\Huge\bfseries ANALOGY}\\[6mm]
\end{center}
\begin{center}
\includegraphics[width=\textwidth, height=0.55\textheight, keepaspectratio]{C:/GitHub/VisualIdeas/data/images/ec4502cf38ad861fe1c29b36bca938de.png}
\end{center}
\vspace{2mm}
\begin{center}
\begin{minipage}{0.92\textwidth}
\small
\textbf{Stumbling blocks can become stepping stones. }

Inspired by my friend, Corey Wilks (who messaged me curious to see how I would visually interpret this idea).

Login or Subscribe to participate
\end{minipage}
\end{center}
\vspace*{\fill}
\newpage
\renewcommand{\theposttitle}{How fixing what's broke is not about knowing what to do}
\renewcommand{\thepostdate}{2024-01-02}
\vspace*{\fill}
\begin{center}
{\Huge\bfseries INTERESTING}\\[6mm]
\end{center}
\begin{center}
\includegraphics[width=\textwidth, height=0.55\textheight, keepaspectratio]{C:/GitHub/VisualIdeas/data/images/bc882005892740c73ae00766ed15eab8.png}
\end{center}
\vspace{2mm}
\begin{center}
\begin{minipage}{0.92\textwidth}
\small
If you want to be more creative, give your ideas time to hatch.

A common mistake I see my students make is giving up too soon.

Try this next time you’re stuck:

go for a walk in silence

take a coffee break

let your subconscious bake your ideas!

This is the Law of Incubation.

Inspired by \textbf{Joey Cofone's book: Laws of Creativity.}
\end{minipage}
\end{center}
\vspace*{\fill}
\newpage
\renewcommand{\theposttitle}{How fixing what's broke is not about knowing what to do}
\renewcommand{\thepostdate}{2024-01-02}
\vspace*{\fill}
\begin{center}
{\Huge\bfseries DESIGN}\\[6mm]
\end{center}
\begin{center}
\begin{minipage}{0.92\textwidth}
\small
\textbf{Clear, Not Clever}

When designing a visual metaphor, one of the most important elements to nail down is the labels you use.

Labels are words, arrows, even punctuation that communicate what the metaphor is.

And the key to using effective labels is clarity. Clarity is step 1 of \textbf{my 3-Step framework to making a visual}.
\end{minipage}
\end{center}
\vspace{4mm}
\begin{center}
\includegraphics[width=\textwidth, height=0.55\textheight, keepaspectratio]{C:/GitHub/VisualIdeas/data/images/c98fa1468cc374dab1f9b70f931c1bd6.png}
\end{center}
\vspace{2mm}
\begin{center}
\begin{minipage}{0.92\textwidth}
\small
What’s critical to remember—clarity is about removing as much friction as possible between you and your audience to get your message across.

(And that means you can’t hide behind fancy words if your message is unclear to begin with)

When you write in a fancy way to impress people, you make them do extra work just so you can seem cool

This issue of Visual I.D.E.A.s is sponsored by Superside:
\end{minipage}
\end{center}
\vspace*{\fill}
\newpage
\renewcommand{\theposttitle}{How fixing what's broke is not about knowing what to do}
\renewcommand{\thepostdate}{2024-01-02}
\vspace*{\fill}
\begin{center}
{\Huge\bfseries ENCHANTING}\\[6mm]
\end{center}
\begin{center}
\includegraphics[width=\textwidth, height=0.55\textheight, keepaspectratio]{C:/GitHub/VisualIdeas/data/images/21d1e95bc2fad39169590c5df6bcc2eb.png}
\end{center}
\vspace{2mm}
\begin{center}
\begin{minipage}{0.92\textwidth}
\small
\textbf{Maximizing your day vs Optimizing your day}

(this is coffee by the way 😆 — I didn’t realize how the holiday season might make this look like hot chocolate…but then again…maybe it still works…what do you think?)
\end{minipage}
\end{center}
\vspace*{\fill}
\newpage
\renewcommand{\theposttitle}{How fixing what's broke is not about knowing what to do}
\renewcommand{\thepostdate}{2024-01-02}
\vspace*{\fill}
\begin{center}
{\Huge\bfseries ANALOGY}\\[6mm]
\end{center}
\begin{center}
\includegraphics[width=\textwidth, height=0.55\textheight, keepaspectratio]{C:/GitHub/VisualIdeas/data/images/429c92d376540e016ff6b0a85f4b07e2.png}
\end{center}
\vspace{2mm}
\begin{center}
\begin{minipage}{0.92\textwidth}
\small
\textbf{“Be big enough to admit mistakes,}\\
\textbf{smart enough to profit from them,}\\
\textbf{strong enough to correct them.”}\\
-John C. Maxwell

Login or Subscribe to participate
\end{minipage}
\end{center}
\vspace*{\fill}
\newpage
\renewcommand{\theposttitle}{How we frame failure paves our road to success}
\renewcommand{\thepostdate}{2024-01-16}
\vspace*{\fill}
\begin{center}
{\Huge\bfseries INTERESTING}\\[6mm]
\end{center}
\begin{center}
\includegraphics[width=\textwidth, height=0.55\textheight, keepaspectratio]{C:/GitHub/VisualIdeas/data/images/28f2ba379c0f25f5b675e15ffa6db341.png}
\end{center}
\vspace{2mm}
\begin{center}
\begin{minipage}{0.92\textwidth}
\small
The mental model that stops me from having an opinion about everything:

\textbf{Dunning-Kruger Effect}--\\
The less you know, the less you question.\\
The more you know, the more you question everything.

I’m fascinated by self-deception and our lack of awareness of how we think and make decisions. And the Dunning-Kruger Effect is a self-deception we all can fall prey to…because well, you don’t know what you don’t know.

One of my goals this year is to reread (and take notes this time!) the book, \textbf{“Thinking Fast and Thinking Slow” by Daniel Kahneman, }because it’s mentioned in practically every book on psychology and sociology ever written since it’s inception. Watch this space for more visuals on cognitive biases 👀
\end{minipage}
\end{center}
\vspace*{\fill}
\newpage
\renewcommand{\theposttitle}{How we frame failure paves our road to success}
\renewcommand{\thepostdate}{2024-01-16}
\vspace*{\fill}
\begin{center}
{\Huge\bfseries DESIGN}\\[6mm]
\end{center}
\begin{center}
\begin{minipage}{0.92\textwidth}
\small
In the last couple weeks I’ve posted a few more process timelapse videos on how I make my visuals.

You might observe \textbf{my process} \textbf{is messy…}and that’s because there’s no better cure for a perfectionist than to let chaos reign! I gave myself permission to be messy 2 years ago when I started making visuals.

I wanted to allow myself to work out ideas through the act of making the thing, not thinking about making the thing.
\end{minipage}
\end{center}
\vspace{4mm}
\begin{center}
\includegraphics[width=\textwidth, height=0.13\textheight, keepaspectratio]{C:/GitHub/VisualIdeas/data/images/8e6607662f41d782005a20d0225ef631.png}
\end{center}
\vspace{2mm}
\begin{center}
\begin{minipage}{0.92\textwidth}
\small
One thing you’ll notice in my videos is I write as I draw. This is important when creating visual metaphors since the labels are a critical part of the meaning.
\end{minipage}
\end{center}
\vspace{5mm}
\begin{center}
\includegraphics[width=\textwidth, height=0.13\textheight, keepaspectratio]{C:/GitHub/VisualIdeas/data/images/61b77be40e7f10992e2b55243fdb5e3d.jpg}
\end{center}
\vspace{5mm}
\begin{center}
\includegraphics[width=\textwidth, height=0.13\textheight, keepaspectratio]{C:/GitHub/VisualIdeas/data/images/8aa2ce386730545ac76fa537fafa137a.jpg}
\end{center}
\vspace{2mm}
\begin{center}
\begin{minipage}{0.92\textwidth}
\small
Which brings up the piece of advice I’d like to leave you today when you’re working through a problem. \textbf{Write ALL ideas down}. Especially the bad ones.

\textbf{Bad ideas are like pinballs that bounce around in your head till you get it out. }Often it’s not till you take the bad ideas out of your head that there’s enough space for the good ideas to rise to the surface!

Cohort 3 of \textbf{Thinking in Visual Metaphors }is live! And it’s the biggest cohort I’ve hosted yet but still keeping it nice and cozy! 🙌
\end{minipage}
\end{center}
\vspace{5mm}
\begin{center}
\includegraphics[width=\textwidth, height=0.13\textheight, keepaspectratio]{C:/GitHub/VisualIdeas/data/images/254b65ca23ba95fc2a22c4fcde79d890.png}
\end{center}
\vspace{2mm}
\begin{center}
\begin{minipage}{0.92\textwidth}
\small
My favorite comment from the cohort so far: \\
\textbf{“It is possible to get started and get going without advanced drawing skills.”}\\
-Kalyana

If you’re curious to learn how to make visual metaphors, you can \textbf{sign up for the waitlist for the next cohort here}. \textbf{ }
\end{minipage}
\end{center}
\vspace*{\fill}
\newpage
\renewcommand{\theposttitle}{How we frame failure paves our road to success}
\renewcommand{\thepostdate}{2024-01-16}
\vspace*{\fill}
\begin{center}
{\Huge\bfseries ENCHANTING}\\[6mm]
\end{center}
\begin{center}
\includegraphics[width=\textwidth, height=0.55\textheight, keepaspectratio]{C:/GitHub/VisualIdeas/data/images/ef4820bbb3fcdc0a4958fa2dd6037490.png}
\end{center}
\vspace{2mm}
\begin{center}
\begin{minipage}{0.92\textwidth}
\small
\textbf{The tighter you squeeze, the less you have.}\\
-Zen Saying
\end{minipage}
\end{center}
\vspace*{\fill}
\newpage
\renewcommand{\theposttitle}{How we frame failure paves our road to success}
\renewcommand{\thepostdate}{2024-01-16}
\vspace*{\fill}
\begin{center}
{\Huge\bfseries ANALOGY}\\[6mm]
\end{center}
\begin{center}
\includegraphics[width=\textwidth, height=0.55\textheight, keepaspectratio]{C:/GitHub/VisualIdeas/data/images/491579c78846557d5e4efd8d5a724b15.png}
\end{center}
\vspace{2mm}
\begin{center}
\begin{minipage}{0.92\textwidth}
\small
Growth in a skill is so slow it's invisible, but if you keep it up, \textbf{it's inevitable}.

Login or Subscribe to participate
\end{minipage}
\end{center}
\vspace*{\fill}
\newpage
\renewcommand{\theposttitle}{The simple shift in my thinking that got me started posting online}
\renewcommand{\thepostdate}{2024-01-30}
\vspace*{\fill}
\begin{center}
{\Huge\bfseries INTERESTING}\\[6mm]
\end{center}
\begin{center}
\includegraphics[width=\textwidth, height=0.55\textheight, keepaspectratio]{C:/GitHub/VisualIdeas/data/images/8da1dfc21e282da50dc5c5662a6618f2.png}
\end{center}
\vspace{2mm}
\begin{center}
\begin{minipage}{0.92\textwidth}
\small
A large portion of the visuals I create are \textbf{“Friendly Reminders”} to not only the people who support my work, but to myself. They echo ideas we know to be true, but are lost in the noise of everyday distractions.

This piece highlights the 3 areas in Milani Creative I’m focusing on in 2024.

\textbf{What are you hoping to do less of (more of) this year?}
\end{minipage}
\end{center}
\vspace*{\fill}
\newpage
\renewcommand{\theposttitle}{The simple shift in my thinking that got me started posting online}
\renewcommand{\thepostdate}{2024-01-30}
\vspace*{\fill}
\begin{center}
{\Huge\bfseries DESIGN}\\[6mm]
\end{center}
\begin{center}
\begin{minipage}{0.92\textwidth}
\small
Today’s design tip is more of an encouragement.

An encouragement to remember—we all start somewhere.

For all the aspiring writers, artists, and creators out there: I’d like to share one of my favorite pieces on the creative struggle from the creator of “This American Life” podcast, Ira Glass:

“Nobody tells this to people who are beginners…

All of us who do creative work, we get into it because we have good taste. \textbf{But there is this gap. }

For the first couple years you make stuff, it’s just not that good.

It’s trying to be good. \\
It has potential. \\
But it’s not.

But your taste, the thing that got you into the game, is still killer. \\
And \textbf{your taste is why your work disappoints you. }

A lot of people never get past this phase. They quit…We all go through this.

And if you are just starting out or you are still in this phase, you gotta know \textbf{its normal} and the most important thing you can do is \textbf{do a lot of work.}”

Which inspired this:
\end{minipage}
\end{center}
\vspace{4mm}
\begin{center}
\includegraphics[width=\textwidth, height=0.55\textheight, keepaspectratio]{C:/GitHub/VisualIdeas/data/images/79b24c2feb3455a8a2bc02e75744698a.png}
\end{center}
\vspace{2mm}
\begin{center}
\begin{minipage}{0.92\textwidth}
\small
I’d highly recommend listening to the full quote (it’s short) and watching this\textbf{ beautiful short film that visualizes it}. I show it to all my high school students who want to step into the creative arena to do battle!
\end{minipage}
\end{center}
\vspace*{\fill}
\newpage
\renewcommand{\theposttitle}{The simple shift in my thinking that got me started posting online}
\renewcommand{\thepostdate}{2024-01-30}
\vspace*{\fill}
\begin{center}
{\Huge\bfseries ENCHANTING}\\[6mm]
\end{center}
\begin{center}
\includegraphics[width=\textwidth, height=0.55\textheight, keepaspectratio]{C:/GitHub/VisualIdeas/data/images/2ca107ddae4b21d6c581b103b30c809c.png}
\end{center}
\vspace{2mm}
\begin{center}
\begin{minipage}{0.92\textwidth}
\small
The growth mindset is an infinite game.
\end{minipage}
\end{center}
\vspace*{\fill}
\newpage
\renewcommand{\theposttitle}{The simple shift in my thinking that got me started posting online}
\renewcommand{\thepostdate}{2024-01-30}
\vspace*{\fill}
\begin{center}
{\Huge\bfseries ANALOGY}\\[6mm]
\end{center}
\begin{center}
\includegraphics[width=\textwidth, height=0.55\textheight, keepaspectratio]{C:/GitHub/VisualIdeas/data/images/d51e292bbe9202c33b028cb941941bb9.png}
\end{center}
\vspace{2mm}
\begin{center}
\begin{minipage}{0.92\textwidth}
\small
In many ways, this visual is connected to “The Gap” idea above.

I procrastinated posting online for over a decade. I kept thinking I’d get started as soon as I knew exactly what I wanted to do. And I didn’t want to post “crap.” I wanted my stuff to be good!

Turns out…this was a great strategy to get older, but not wiser.

I didn’t do anything—for literally over a decade! Because the sad truth was I had no idea what “good” was and what “crap” was.

But then I read a book by Derek Sivers called \textbf{“Anything You Want: 40 Lessons for a New Kind of Entrepreneur.” }

And I came across the following excerpt:

“If you want to be useful, you can always start now, with only 1 percent of what you have in your grand vision.

It'll be a humble prototype version of your grand vision, but you'll be in the game.

You'll be ahead of the rest, because you actually started, while others are waiting for the finish line to magically appear at the starting line.“

This one idea became my action manifesto.

\textbf{Take the small bet version before putting all your chips on the table. }

Test first, double down what works.

It lifted the paralysis that kept me from progressing…for over a decade!
\end{minipage}
\end{center}
\vspace*{\fill}
\newpage
\renewcommand{\theposttitle}{What most people get wrong about iteration and change}
\renewcommand{\thepostdate}{2024-02-13}
\vspace*{\fill}
\begin{center}
{\Huge\bfseries INTERESTING}\\[6mm]
\end{center}
\begin{center}
\begin{minipage}{0.92\textwidth}
\small
I have to remind myself of this idea everyday:

\textbf{The choices you make every day matter. }
\end{minipage}
\end{center}
\vspace{4mm}
\begin{center}
\includegraphics[width=\textwidth, height=0.24\textheight, keepaspectratio]{C:/GitHub/VisualIdeas/data/images/89825126d1bd6ee488c4cba032b8bddf.png}
\end{center}
\vspace{5mm}
\begin{center}

\includegraphics[width=\textwidth, height=0.24\textheight, keepaspectratio]{C:/GitHub/VisualIdeas/data/images/43370857b5a4ea6f557287991c787ce7.png}
\end{center}
\vspace{2mm}
\begin{center}
\begin{minipage}{0.92\textwidth}
\small
The snowball compounding impact of daily decisions is challenging to wrap our heads around (especially when we make decisions in the counter direction we’re aiming for).

Recently I’ve been falling into the “I’ll do better tomorrow” trap…especially with diet.

My kryptonite—chocolate chip cookies…and guess what they serve everyday at the school cafeteria? 😅

“Surely this cookie today isn’t that big of a deal,” I catch myself saying for the 3rd week in a row…

What are you working towards this season? I’d love to hear! \\
(I may not always get a chance to reply but I read every email ❤️)

\textbf{1\% better: 1.01}\textbf{365 }\textbf{= 37.78}\\
\textbf{1\% worse: 0.99}\textbf{365 }\textbf{= 0.03}

Inspired by the book, \textbf{“Atomic Habits” by James Clear}\\
*I reread James Clear’s article on this idea regularly: \textbf{Continuous Improvement: How It Works and How to Master It.}
\end{minipage}
\end{center}
\vspace*{\fill}
\newpage
\renewcommand{\theposttitle}{What most people get wrong about iteration and change}
\renewcommand{\thepostdate}{2024-02-13}
\vspace*{\fill}
\begin{center}
{\Huge\bfseries DESIGN}\\[6mm]
\end{center}
\begin{center}
\begin{minipage}{0.92\textwidth}
\small
I’ve spoken before about the importance of iteration to bridge the gap between your ability and taste.

Today, I want to illustrate (🥁) \textbf{a common mistake} people make when they iterate: \textbf{focusing too heavily on “improving” the image}.

Of course we want the image to be better, but starting with that mindset actually blocks the creative process.

I tell my students all the time that \textbf{brainstorming is idea generation, not decision making. }

Starting with improvement in mind is decision making before the ideas have generated.

❌ how can I improve it? (explores only ideas you assume are good)\\
✅ how can articulate this idea in another way? (lets in all ideas, good and bad—we want ALL ideas to flow in our stream of consciousness)

\textbf{Improving is more a byproduct of trying out different articulations.}

Take this recent example of an iteration:
\end{minipage}
\end{center}
\vspace{4mm}
\begin{center}
\includegraphics[width=\textwidth, height=0.24\textheight, keepaspectratio]{C:/GitHub/VisualIdeas/data/images/480892248086a40f0e4ee5610b5bcee0.png}
\end{center}
\vspace{2mm}
\begin{center}
\begin{minipage}{0.92\textwidth}
\small
Made in August 2022
\end{minipage}
\end{center}
\vspace{5mm}
\begin{center}
\includegraphics[width=\textwidth, height=0.24\textheight, keepaspectratio]{C:/GitHub/VisualIdeas/data/images/8fd1a2ed542c39db5205c6430fd692ba.png}
\end{center}
\vspace{2mm}
\begin{center}
\begin{minipage}{0.92\textwidth}
\small
The nuance of these images are completely different but the core message is the same: address problems as they form to avoid it getting so convoluted that they become difficult to manage.

For me, one of the best parts of thinking in visual metaphors is the forcing function of having to get crystal clear on what you’re trying to say in your content.

\textbf{CLARITY = iteration x time}

Ultimately, whether version 2 is better or worse is irrelevant. The act of creating a different articulation opens another door to how deeply you think about the idea.

When you’re ready and if you’re curious to learn how to make visual metaphors like I do, \textbf{sign up for the waitlist for the next Thinking in Visual Metaphors cohort. }
\end{minipage}
\end{center}
\vspace*{\fill}
\newpage
\renewcommand{\theposttitle}{What most people get wrong about iteration and change}
\renewcommand{\thepostdate}{2024-02-13}
\vspace*{\fill}
\begin{center}
{\Huge\bfseries ENCHANTING}\\[6mm]
\end{center}
\begin{center}

\includegraphics[width=\textwidth, height=0.55\textheight, keepaspectratio]{C:/GitHub/VisualIdeas/data/images/cd3d14fe5694642597e94e75c3e06510.png}
\end{center}
\vspace{2mm}
\begin{center}
\begin{minipage}{0.92\textwidth}
\small
This Japanese concept is the secret to staying consistent in your pursuits:

\textbf{Ikigai} (生き甲斐)\\
⦿ \textit{iki} - means to live\\
⦿ \textit{gai} - means value or worth (literally translated to "beautiful shell")

Ikigai roughly means your "reason for being"--your purpose.

Find your purpose and you find the foundation to build your dream.

If you’re interested in learning more, I’d recommend checking out \textbf{“Ikigai: The Japanese Secret to a Long and Happy Life” by Héctor García. }

It took me a while to align what I do online and what I do in my day job as a teacher, but I finally found the words to articulate my ikigai:

\textbf{Encourage people to live a more creative life to make the world a more creative place. }

What's your ikigai?
\end{minipage}
\end{center}
\vspace*{\fill}
\newpage
\renewcommand{\theposttitle}{What most people get wrong about iteration and change}
\renewcommand{\thepostdate}{2024-02-13}
\vspace*{\fill}
\begin{center}
{\Huge\bfseries ANALOGY}\\[6mm]
\end{center}
\begin{center}
\includegraphics[width=\textwidth, height=0.55\textheight, keepaspectratio]{C:/GitHub/VisualIdeas/data/images/31c9a2117244c7e28486b7fd5b13b7ff.png}
\end{center}
\vspace{2mm}
\begin{center}
\begin{minipage}{0.92\textwidth}
\small
\textbf{"Change is hard at first, messy in the middle and gorgeous at the end." }\\
–Robin Sharma

Login or Subscribe to participate
\end{minipage}
\end{center}
\vspace*{\fill}
\newpage
\renewcommand{\theposttitle}{Why you accomplish less than you want...but more than you think}
\renewcommand{\thepostdate}{2024-02-27}
\vspace*{\fill}
\begin{center}
{\Huge\bfseries INTERESTING}\\[6mm]
\end{center}
\begin{center}
\includegraphics[width=\textwidth, height=0.24\textheight, keepaspectratio]{C:/GitHub/VisualIdeas/data/images/c8a27074c4860aa285ee2545018821e1.png}
\end{center}
\vspace{2mm}
\begin{center}
\begin{minipage}{0.92\textwidth}
\small
An unfortunate truth:\\
\textbf{Good judgment comes from experience.}\\
\textbf{Experience comes from bad judgment.}

I used to write this in on the first page of each journal I would gift to my filmmaking students.

\textbf{The biggest lessons are a result of making an expensive mistake} (expensive in time, money…heartbreak).

The earlier in our lives we make these mistakes, the more time we have to enjoy the compounding benefits of the lessons they teach.
\end{minipage}
\end{center}
\vspace{5mm}
\begin{center}
\includegraphics[width=\textwidth, height=0.24\textheight, keepaspectratio]{C:/GitHub/VisualIdeas/data/images/b6546c7157e28e749231e098a49bcad7.png}
\end{center}
\vspace{2mm}
\begin{center}
\begin{minipage}{0.92\textwidth}
\small
\textbf{What people "miss" about success.}
\end{minipage}
\end{center}
\vspace*{\fill}
\newpage
\renewcommand{\theposttitle}{Why you accomplish less than you want...but more than you think}
\renewcommand{\thepostdate}{2024-02-27}
\vspace*{\fill}
\begin{center}
{\Huge\bfseries DESIGN}\\[6mm]
\end{center}
\begin{center}
\begin{minipage}{0.92\textwidth}
\small
A simple trick to elevate the impact of a bar graph: \\
\textbf{Substitute the bars in the graph with a symbol. }
\end{minipage}
\end{center}
\vspace{4mm}
\begin{center}
\includegraphics[width=\textwidth, height=0.24\textheight, keepaspectratio]{C:/GitHub/VisualIdeas/data/images/45ec8553203d0b9d00bb824e2d917fcb.png}
\end{center}
\vspace{2mm}
\begin{center}
\begin{minipage}{0.92\textwidth}
\small
Version 1
\end{minipage}
\end{center}
\vspace{5mm}
\begin{center}
\includegraphics[width=\textwidth, height=0.24\textheight, keepaspectratio]{C:/GitHub/VisualIdeas/data/images/f8d2d5bb1e2895700c7a713743ab5558.png}
\end{center}
\vspace{2mm}
\begin{center}
\begin{minipage}{0.92\textwidth}
\small
In version 1, I clarified this idea I had about how deeply you have to understand something to teach it. It got the message across but one of my core creative practices is to \textbf{keep pushing}.

How can I make this: \\
• clearer?\\
• more insightful? \\
• more memorable?

That’s when I came up with version 2 to add another layer of nuance to the visual: substitute the bars with books.

This technique comes from my filmmaking background and I call it: \textbf{superimposition}. It’s one of the tools I teach in my \textbf{Thinking in Visual Metaphors} course.
\end{minipage}
\end{center}
\vspace*{\fill}
\newpage
\renewcommand{\theposttitle}{Why you accomplish less than you want...but more than you think}
\renewcommand{\thepostdate}{2024-02-27}
\vspace*{\fill}
\begin{center}
{\Huge\bfseries ENCHANTING}\\[6mm]
\end{center}
\begin{center}
\includegraphics[width=\textwidth, height=0.55\textheight, keepaspectratio]{C:/GitHub/VisualIdeas/data/images/bd2465e126f5b77ea8800bb82bdb5b4b.png}
\end{center}
\vspace{2mm}
\begin{center}
\begin{minipage}{0.92\textwidth}
\small
Join the conversation on LinkedIn

People overestimate what they can get done in a day and underestimate what they can get done in a year.
\end{minipage}
\end{center}
\vspace*{\fill}
\newpage
\renewcommand{\theposttitle}{Why you accomplish less than you want...but more than you think}
\renewcommand{\thepostdate}{2024-02-27}
\vspace*{\fill}
\begin{center}
{\Huge\bfseries ANALOGY}\\[6mm]
\end{center}
\begin{center}
\includegraphics[width=\textwidth, height=0.55\textheight, keepaspectratio]{C:/GitHub/VisualIdeas/data/images/e8e2fa67d6b5bfca081a5a337f506b00.png}
\end{center}
\vspace{2mm}
\begin{center}
\begin{minipage}{0.92\textwidth}
\small
If you think others have it figured out while you're still struggling, remember this:

\textbf{The Swan Paradox:}\\
It takes an obscene amount of effort to make what you're doing look effortless.

And the key thing to remember is to stay afloat, the effort's ongoing.\\
It's what it means to be consistent.

You'll drown if you rely only on burst efforts.

Pace yourself.
\end{minipage}
\end{center}
\vspace*{\fill}
\newpage
\renewcommand{\theposttitle}{The Japanese philosophy that got me over perfectionism}
\renewcommand{\thepostdate}{2024-03-12}
\vspace*{\fill}
\begin{center}
{\Huge\bfseries DESIGN}\\[6mm]
\end{center}
\begin{center}
\begin{minipage}{0.92\textwidth}
\small
This idea is at the core of my creative philosophy…and it might help if you struggle with perfectionism as much as I do:
\end{minipage}
\end{center}
\vspace{4mm}
\begin{center}
\includegraphics[width=\textwidth, height=0.24\textheight, keepaspectratio]{C:/GitHub/VisualIdeas/data/images/fb443804a8195eceef0f096d2c4c3f99.png}
\end{center}
\vspace{2mm}
\begin{center}
\begin{minipage}{0.92\textwidth}
\small
Quote from Ernest Heminway, “A Farewell To Arms”

\textbf{Wabi-Sabi}…

…a Zen aesthetic ideal that celebrates flaws, recognizing the beauty in imperfection and incompleteness.

The practice of wabi-sabi is the perfect (🥁) counternote to perfectionism.

You’ll find this philosophy practiced in the Japanese artform, \textbf{kintsugi}.

In kintsugi, pottery is intentionally broken, then repaired with a golden lacquer to highlight its damage.

One thing I’ve learned from 20 years of teaching and making 400+ visual metaphors is that art isn’t about being accurate. It’s about being interesting.

And what makes you interesting is your unique combination of imperfections and flaws.

Instead of trying to airbrush them away, highlight them.

Your imperfections are the feature, not the bug.

When you’re ready, consider checking out \textbf{Thinking in Visual Metaphors}. Registration for Cohort 4 closes in less than 2 weeks.
\end{minipage}
\end{center}
\vspace{5mm}
\begin{center}

\includegraphics[width=\textwidth, height=0.24\textheight, keepaspectratio]{C:/GitHub/VisualIdeas/data/images/e9f5cefec10bcfaed5660eaf3c9f0d6f.png}
\end{center}
\vspace{2mm}
\begin{center}
\begin{minipage}{0.92\textwidth}
\small
\textbf{Learn how to make AI work for you.}

AI won’t take your job, but a person using AI might. That’s why 500,000+ professionals read The Rundown – the free newsletter that keeps you updated on the latest AI news and teaches you how to apply it in just 5 minutes a day.

Subscribe with one click.
\end{minipage}
\end{center}
\vspace*{\fill}
\newpage
\renewcommand{\theposttitle}{The Japanese philosophy that got me over perfectionism}
\renewcommand{\thepostdate}{2024-03-12}
\vspace*{\fill}
\begin{center}
{\Huge\bfseries ENCHANTING}\\[6mm]
\end{center}
\begin{center}
\includegraphics[width=\textwidth, height=0.55\textheight, keepaspectratio]{C:/GitHub/VisualIdeas/data/images/d2f8970ba0bfd2b28bf1fe2a1cfe7c0a.png}
\end{center}
\vspace{2mm}
\begin{center}
\begin{minipage}{0.92\textwidth}
\small
\textbf{Reflect on your potential.}

This image brought on some thoughtful reflections (🥁) on LinkedIn…including the dangers of narcism.

I try and steer away from leading people on how to interpret my visuals, but I do love seeing thoughtful conversation that arise from the topics themselves.
\end{minipage}
\end{center}
\vspace*{\fill}
\newpage
\renewcommand{\theposttitle}{The Japanese philosophy that got me over perfectionism}
\renewcommand{\thepostdate}{2024-03-12}
\vspace*{\fill}
\begin{center}
{\Huge\bfseries ANALOGY}\\[6mm]
\end{center}
\begin{center}
\includegraphics[width=\textwidth, height=0.55\textheight, keepaspectratio]{C:/GitHub/VisualIdeas/data/images/382034b73ac816cae100d1cf7783dadd.png}
\end{center}
\vspace{2mm}
\begin{center}
\begin{minipage}{0.92\textwidth}
\small
It's never a clean fit (and that's okay)
\end{minipage}
\end{center}
\vspace*{\fill}
\newpage
\renewcommand{\theposttitle}{The question most people don't ask me and why it matters}
\renewcommand{\thepostdate}{2024-03-26}
\vspace*{\fill}
\begin{center}
{\Huge\bfseries INTERESTING}\\[6mm]
\end{center}
\begin{center}
\includegraphics[width=\textwidth, height=0.55\textheight, keepaspectratio]{C:/GitHub/VisualIdeas/data/images/51ec5c2a1471533e53ea27a556dbf422.png}
\end{center}
\vspace{2mm}
\begin{center}
\begin{minipage}{0.92\textwidth}
\small
\textbf{Comfort zones are tiny prisons that keep us }\\
\textbf{from experiencing the vastness of life.}
\end{minipage}
\end{center}
\vspace*{\fill}
\newpage
\renewcommand{\theposttitle}{The question most people don't ask me and why it matters}
\renewcommand{\thepostdate}{2024-03-26}
\vspace*{\fill}
\begin{center}
{\Huge\bfseries DESIGN}\\[6mm]
\end{center}
\begin{center}
\begin{minipage}{0.92\textwidth}
\small
I get asked all the time what software* I use to make my visuals.

I understand why it’s asked but this question frames the wrong problem.

Software enables thinking, but without the thinking, a piece of software is just an empty shell.

The more important question to ask is: \\
how do you translate an idea into a visual metaphor?

For me, it’s a 3-step process.
\end{minipage}
\end{center}
\vspace{4mm}
\begin{center}
\includegraphics[width=\textwidth, height=0.17\textheight, keepaspectratio]{C:/GitHub/VisualIdeas/data/images/17c5ff7d616dd93bd357fc9f7f5419cd.png}
\end{center}
\vspace{2mm}
\begin{center}
\begin{minipage}{0.92\textwidth}
\small
Everything is downstream of clarity. Without it, no amount of graphics will resonate with the audience.
\end{minipage}
\end{center}
\vspace{5mm}
\begin{center}
\includegraphics[width=\textwidth, height=0.17\textheight, keepaspectratio]{C:/GitHub/VisualIdeas/data/images/bdd08457ad1465c80eb3125f4af5d13a.png}
\end{center}
\vspace{2mm}
\begin{center}
\begin{minipage}{0.92\textwidth}
\small
You need to give your message shape so your audience can process it fast and easily.
\end{minipage}
\end{center}
\vspace{5mm}
\begin{center}
\includegraphics[width=\textwidth, height=0.17\textheight, keepaspectratio]{C:/GitHub/VisualIdeas/data/images/4b8caf9e9763bd0019d0f9866b047eb8.png}
\end{center}
\vspace{2mm}
\begin{center}
\begin{minipage}{0.92\textwidth}
\small
This step is critical. Most visual thinkers believe you have to blatantly hammer your reader over the head with the message. But this makes the reader passive.

I believe in trusting the reader.

Leave space in the message for the audience to fill in the gap themselves. It’s not about you looking smart. It’s about making your reader feel smart.

We’re just touching the surface here, but this is the general overview of the process: \textbf{Clarify =\textgreater{} Simplify =\textgreater{} Amplify}

*For the record, the software I use is Procreate on the iPad with an Apple Pencil. 😊

P.S. If you want to start translating your own ideas into visuals, it’s not too late to sign up for \textbf{Cohort 4 of Thinking in Visual Metaphors}\textbf{ }but time’s running out for enrollment. Class begins Wednesday, March 27th!
\end{minipage}
\end{center}
\vspace*{\fill}
\newpage
\renewcommand{\theposttitle}{The question most people don't ask me and why it matters}
\renewcommand{\thepostdate}{2024-03-26}
\vspace*{\fill}
\begin{center}
{\Huge\bfseries ENCHANTING}\\[6mm]
\end{center}
\begin{center}
\includegraphics[width=\textwidth, height=0.55\textheight, keepaspectratio]{C:/GitHub/VisualIdeas/data/images/60b148139364d09369d15ea685ecad41.png}
\end{center}
\vspace{2mm}
\begin{center}
\begin{minipage}{0.92\textwidth}
\small
\textbf{3 best returns on investments 👆}
\end{minipage}
\end{center}
\vspace*{\fill}
\newpage
\renewcommand{\theposttitle}{The question most people don't ask me and why it matters}
\renewcommand{\thepostdate}{2024-03-26}
\vspace*{\fill}
\begin{center}
{\Huge\bfseries ANALOGY}\\[6mm]
\end{center}
\begin{center}
\includegraphics[width=\textwidth, height=0.55\textheight, keepaspectratio]{C:/GitHub/VisualIdeas/data/images/a3ed29c2bd9b16daeec93ee4276a0976.png}
\end{center}
\vspace{2mm}
\begin{center}
\begin{minipage}{0.92\textwidth}
\small
\textbf{A friendly reminder to appreciate the effort it takes }\\
\textbf{to make something look effortless. }
\end{minipage}
\end{center}
\vspace*{\fill}
\newpage
\renewcommand{\theposttitle}{Where to focus your attention to improve your quality of life}
\renewcommand{\thepostdate}{2024-04-09}
\vspace*{\fill}
\begin{center}
{\Huge\bfseries DESIGN}\\[6mm]
\end{center}
\begin{center}
\begin{minipage}{0.92\textwidth}
\small
If you’ve followed me for a while, you might notice I’ll often revisit an idea—to find a different articulation of that message.

There’re 3 reasons why:

Each articulation forces diving into \textbf{deeper understanding}.

\textbf{Improves pattern recognition} of how things are connected.

Different metaphors will \textbf{impact people differently}.

There’s also a more personal reason worth mentioning—it keeps my perfectionist monster at bay.

If I know the visual I’m creating is just ‘one of many,’ then I’m less precious about finding the ‘perfect analogy.’

Here’s an example of a recent articulation of “fancy language vs simple language.”
\end{minipage}
\end{center}
\vspace{4mm}
\begin{center}
\includegraphics[width=\textwidth, height=0.24\textheight, keepaspectratio]{C:/GitHub/VisualIdeas/data/images/63378dafd8099973767c6dd5237dd4a8.png}
\end{center}
\vspace{2mm}
\begin{center}
\begin{minipage}{0.92\textwidth}
\small
And here are two earlier articulations of that same idea:
\end{minipage}
\end{center}
\vspace{5mm}
\begin{center}
\includegraphics[width=\textwidth, height=0.24\textheight, keepaspectratio]{C:/GitHub/VisualIdeas/data/images/396d389c48dc7bb797a3d26fa9c01ab2.png}
\end{center}
\vspace*{\fill}
\newpage
\renewcommand{\theposttitle}{Where to focus your attention to improve your quality of life}
\renewcommand{\thepostdate}{2024-04-09}
\vspace*{\fill}
\begin{center}
{\Huge\bfseries ENCHANTING}\\[6mm]
\end{center}
\begin{center}
\includegraphics[width=\textwidth, height=0.55\textheight, keepaspectratio]{C:/GitHub/VisualIdeas/data/images/7a19899502a84c63a7d2ae0fd0394e56.png}
\end{center}
\vspace{2mm}
\begin{center}
\begin{minipage}{0.92\textwidth}
\small
\textbf{The role of a teacher. }
\end{minipage}
\end{center}
\vspace*{\fill}
\newpage
\renewcommand{\theposttitle}{Where to focus your attention to improve your quality of life}
\renewcommand{\thepostdate}{2024-04-09}
\vspace*{\fill}
\begin{center}
{\Huge\bfseries ANALOGY}\\[6mm]
\end{center}
\begin{center}
\includegraphics[width=\textwidth, height=0.55\textheight, keepaspectratio]{C:/GitHub/VisualIdeas/data/images/35d3deee1d8e4992f40e5e397938f2fd.png}
\end{center}
\vspace{2mm}
\begin{center}
\begin{minipage}{0.92\textwidth}
\small
"The master has failed more times than the beginner has even tried."

Most people compare destinations without comparing journeys.

I'll give you an example…

For 13 years, Mike Winkelmann (aka "beeple") made a piece of art, everyday (aptly called his “everydays”).

Day 5,000, Mike takes each of his everydays and stitches them together into an NFT of a mosaic called: "The First 5000 Days."

And...

October 2021, this NFT sells for…wait for it—\\
\$69 million.

Media labeled it an "overnight success."

But there's no such thing as an overnight success. Only over time success.

13 years of shooting daily shots.

Beeple still creates today. 6183 days and counting!

This is what it means to be a master.

Hit or miss, you keep shooting.
\end{minipage}
\end{center}
\vspace*{\fill}
\newpage
\renewcommand{\theposttitle}{Why focusing on one thing at a time is a superpower}
\renewcommand{\thepostdate}{2024-04-25}
\vspace*{\fill}
\begin{center}
{\Huge\bfseries INTERESTING}\\[6mm]
\end{center}
\begin{center}
\includegraphics[width=\textwidth, height=0.55\textheight, keepaspectratio]{C:/GitHub/VisualIdeas/data/images/1173a0db1c267ca1d61b6f0d13c9b4d2.png}
\end{center}
\vspace{2mm}
\begin{center}
\begin{minipage}{0.92\textwidth}
\small
“You can do anything, but not everything.” (…all at once)

I admit, before I started posting online a couple years ago, I was notoriously bad at this. And my wife argues I still have quite the room to grow in other domains (like tidying up the bedroom!).

But I make conscious (and often painful) efforts to focus on one cup at a time when it comes to being a creator.

Add your \textbf{“ONE thing you're prioritizing this week” to the conversation here. }

My priority this week is to reflect on my father’s life through journaling.
\end{minipage}
\end{center}
\vspace*{\fill}
\newpage
\renewcommand{\theposttitle}{Why focusing on one thing at a time is a superpower}
\renewcommand{\thepostdate}{2024-04-25}
\vspace*{\fill}
\begin{center}
{\Huge\bfseries DESIGN}\\[6mm]
\end{center}
\begin{center}
\begin{minipage}{0.92\textwidth}
\small
A question came up in my \textbf{Thinking in Visual Metaphors} course about \textbf{avatars}.

Do you create a reference sheet for your avatar with different moods and poses?

This isn’t a bad idea, but I don’t recommend spending the time to do this when you’re first starting out.

Why?

For years, I struggled to make anything because I thought I needed to do x, y, and z to START. Turned out it was just productive procrastination.

It wasn’t until I realized, you know what? It’s totally fine to build as you go…make what you need as you need it (applies to many things, not just visuals).

An hour spent making 20 different facial expressions is an hour not spent on creating the visual that needs only 2. And once you have those 2, you can reuse them like this:
\end{minipage}
\end{center}
\vspace{4mm}
\begin{center}
\includegraphics[width=\textwidth, height=0.17\textheight, keepaspectratio]{C:/GitHub/VisualIdeas/data/images/11e30b310e17184ed0f28b0d384f7fb9.png}
\end{center}
\vspace{5mm}
\begin{center}
\includegraphics[width=\textwidth, height=0.17\textheight, keepaspectratio]{C:/GitHub/VisualIdeas/data/images/ceb9a65f3b0673eb60cb464a6e36f474.png}
\end{center}
\vspace{2mm}
\begin{center}
\begin{minipage}{0.92\textwidth}
\small
The 80/20 rule applies here: You’ll probably use 20\% of those poses for 80\% of your visuals.

So build that reference sheet, but only as you’re building your visuals.
\end{minipage}
\end{center}
\vspace{5mm}
\begin{center}
\includegraphics[width=\textwidth, height=0.17\textheight, keepaspectratio]{C:/GitHub/VisualIdeas/data/images/ca498544e646fb882398d1e0034e12e5.png}
\end{center}
\vspace{2mm}
\begin{center}
\begin{minipage}{0.92\textwidth}
\small
Build as you go is my philosophy.

P.S. If you’re ready to start \textbf{Thinking in Visual Metaphors}, the next cohort starts late May.
\end{minipage}
\end{center}
\vspace*{\fill}
\newpage
\renewcommand{\theposttitle}{Why focusing on one thing at a time is a superpower}
\renewcommand{\thepostdate}{2024-04-25}
\vspace*{\fill}
\begin{center}
{\Huge\bfseries ENCHANTING}\\[6mm]
\end{center}
\begin{center}
\includegraphics[width=\textwidth, height=0.55\textheight, keepaspectratio]{C:/GitHub/VisualIdeas/data/images/12a9f6ed36bc29cc7bb5fa7c88c04483.png}
\end{center}
\vspace{2mm}
\begin{center}
\begin{minipage}{0.92\textwidth}
\small
\textbf{The longest journey you'll ever make.}
\end{minipage}
\end{center}
\vspace*{\fill}
\newpage
\renewcommand{\theposttitle}{Why focusing on one thing at a time is a superpower}
\renewcommand{\thepostdate}{2024-04-25}
\vspace*{\fill}
\begin{center}
{\Huge\bfseries ANALOGY}\\[6mm]
\end{center}
\begin{center}
\includegraphics[width=\textwidth, height=0.55\textheight, keepaspectratio]{C:/GitHub/VisualIdeas/data/images/cbd04756ee263ef493152c95148bd741.jpeg}
\end{center}
\vspace{2mm}
\begin{center}
\begin{minipage}{0.92\textwidth}
\small
Most of us (myself included 😅) get caught up in shiny object syndrome which splinters our focus a mile wide, but an inch deep.

It’s the paradox of the creator. We seek novelty through experimentation, but need to improve what we do through repetition.

If we truly want to upgrade our thinking and skillsets, we need to dive an inch wide, a mile deep.

The shiny objects will always be waiting for us at the surface when we come back up.
\end{minipage}
\end{center}
\vspace*{\fill}
\newpage
\renewcommand{\theposttitle}{What most people get wrong about how to stand out as a creator}
\renewcommand{\thepostdate}{2024-05-14}
\vspace*{\fill}
\begin{center}
{\Huge\bfseries INTERESTING}\\[6mm]
\end{center}
\begin{center}
\includegraphics[width=\textwidth, height=0.55\textheight, keepaspectratio]{C:/GitHub/VisualIdeas/data/images/b143a59387b055a4c003e87200a60d48.png}
\end{center}
\vspace{2mm}
\begin{center}
\begin{minipage}{0.92\textwidth}
\small
\textbf{Friendly reminder: Perfection shouldn’t even be the goal. }

This is at the heart of my creator philosophy:

One perfect circle looks just like any other perfect circle. Most people think they can’t draw because they’re comparing their drawings to what they believe it SHOULD look like.

But if you want to make stuff that stands out, it’s less important to be accurate and more important to be interesting.
\end{minipage}
\end{center}
\vspace*{\fill}
\newpage
\renewcommand{\theposttitle}{What most people get wrong about how to stand out as a creator}
\renewcommand{\thepostdate}{2024-05-14}
\vspace*{\fill}
\begin{center}
{\Huge\bfseries DESIGN}\\[6mm]
\end{center}
\begin{center}
\begin{minipage}{0.92\textwidth}
\small
I’ve been teaching film \& animation for 20 years, and I try to bring filmmaking concepts into my visuals whenever I can. One of my favorite things to play with is \textbf{movement! }

Now how do we add movement to a static image? \\
\textbf{Motion Lines. }

There are 3 main types:
\end{minipage}
\end{center}
\vspace{4mm}
\begin{center}
\includegraphics[width=\textwidth, height=0.13\textheight, keepaspectratio]{C:/GitHub/VisualIdeas/data/images/10351bae12fbf3ca9801464bdb323c5f.png}
\end{center}
\vspace{2mm}
\begin{center}
\begin{minipage}{0.92\textwidth}
\small
Try it yourself! Here’s a simple exercise to make straight motion:
\end{minipage}
\end{center}
\vspace{5mm}
\begin{center}
\includegraphics[width=\textwidth, height=0.13\textheight, keepaspectratio]{C:/GitHub/VisualIdeas/data/images/a5b2ac22a98d07d49ee2819371e604f3.png}
\end{center}
\vspace{2mm}
\begin{center}
\begin{minipage}{0.92\textwidth}
\small
Draw a simple circle
\end{minipage}
\end{center}
\vspace{5mm}
\begin{center}
\includegraphics[width=\textwidth, height=0.13\textheight, keepaspectratio]{C:/GitHub/VisualIdeas/data/images/07128bf0fadd1c067fc77d361f55fe53.png}
\end{center}
\vspace{2mm}
\begin{center}
\begin{minipage}{0.92\textwidth}
\small
And make some straight lines over it!

It’s as simple as that!

And by the way, if you’re keen to learn how to make visual metaphors, enrollment for Thinking in Visual Metaphors is open for another week! 28 students have already enrolled in Cohort 5!
\end{minipage}
\end{center}
\vspace{5mm}
\begin{center}
\includegraphics[width=\textwidth, height=0.13\textheight, keepaspectratio]{C:/GitHub/VisualIdeas/data/images/f3c8c479a4856823f006dd4dbdec07a0.png}
\end{center}
\vspace*{\fill}
\newpage
\renewcommand{\theposttitle}{What most people get wrong about how to stand out as a creator}
\renewcommand{\thepostdate}{2024-05-14}
\vspace*{\fill}
\begin{center}
{\Huge\bfseries ANALOGY}\\[6mm]
\end{center}
\begin{center}
\includegraphics[width=\textwidth, height=0.55\textheight, keepaspectratio]{C:/GitHub/VisualIdeas/data/images/a63a99c60eb0ddf4b9ba25859ff33661.png}
\end{center}
\vspace{2mm}
\begin{center}
\begin{minipage}{0.92\textwidth}
\small
\textbf{The price of admission for success.}
\end{minipage}
\end{center}
\vspace*{\fill}
\newpage
\renewcommand{\theposttitle}{How to overcome your impossible standards}
\renewcommand{\thepostdate}{2024-05-29}
\vspace*{\fill}
\begin{center}
{\Huge\bfseries INTERESTING}\\[6mm]
\end{center}
\begin{center}
\includegraphics[width=\textwidth, height=0.55\textheight, keepaspectratio]{C:/GitHub/VisualIdeas/data/images/503f6b35b967b6ae47a9c1f99a5367e3.png}
\end{center}
\vspace{2mm}
\begin{center}
\begin{minipage}{0.92\textwidth}
\small
\textbf{Avoid the second mistake.}\\
Inspired by James Clear’s book, \textbf{Atomic Habits}.

Now the trick is how?

The best piece of advice I’ve heard from James Clear is, “reduce the scope to maintain the habit.”

Most of my life I’ve held almost impossible to meet standards, for myself and others. While this was useful to me for a season, I realized over time those standards distorted my creative output into an all or nothing approach.

The harsh truth is that \textbf{a balanced life doesn’t mean balanced days.} Balance comes in the form of an averaging of your days.

There will be days when I only get a few push-ups or a brisk walk in for exercise. There are days when I get home so late that I only have enough time with the kids for a short cuddle and bedtime story.

The point is to embrace the imperfect messiness of it all. Miss a day, you’re human. Try your best to get back on the bike day 2, even if it means you don’t go as far forward as you’d like.

A little progress is still progress.
\end{minipage}
\end{center}
\vspace*{\fill}
\newpage
\renewcommand{\theposttitle}{How to overcome your impossible standards}
\renewcommand{\thepostdate}{2024-05-29}
\vspace*{\fill}
\begin{center}
{\Huge\bfseries DESIGN}\\[6mm]
\end{center}
\begin{center}
\begin{minipage}{0.92\textwidth}
\small
I have a simple hack to fight my perfectionist tendencies:

\textbf{Draw on paper…with a pen.}

This practice forces you to:

Stop being precious

Be more efficient with your lines

Avoid wasting time erasing to make things perfect
\end{minipage}
\end{center}
\vspace{4mm}
\begin{center}
%\includegraphics[width=\textwidth, height=0.55\textheight, keepaspectratio]{C:/GitHub/VisualIdeas/data/images/ef1ea0d4c3e11ddb083690f8403143b8.gif}
\end{center}
\vspace*{\fill}
\newpage
\renewcommand{\theposttitle}{How to overcome your impossible standards}
\renewcommand{\thepostdate}{2024-05-29}
\vspace*{\fill}
\begin{center}
{\Huge\bfseries ENCHANTING}\\[6mm]
\end{center}
\begin{center}
\includegraphics[width=\textwidth, height=0.55\textheight, keepaspectratio]{C:/GitHub/VisualIdeas/data/images/a805c243f8b393c74310115be856084e.png}
\end{center}
\vspace{2mm}
\begin{center}
\begin{minipage}{0.92\textwidth}
\small
\textbf{Friendly reminder of the day.}

I try to think of my visual metaphors as “friendly reminders” to the reader (and myself). I hope they FEEL true, like a familiar echo…that somehow got lost in the sea of relentless noise.

This one is a simple reminder to myself that the negative stuff we each face day in and day out are what help us appreciate the positive stuff.

If you like friendly reminders, I posted a bunch more\textbf{ to start the week here. }
\end{minipage}
\end{center}
\vspace*{\fill}
\newpage
\renewcommand{\theposttitle}{How to overcome your impossible standards}
\renewcommand{\thepostdate}{2024-05-29}
\vspace*{\fill}
\begin{center}
{\Huge\bfseries ANALOGY}\\[6mm]
\end{center}
\begin{center}
\includegraphics[width=\textwidth, height=0.55\textheight, keepaspectratio]{C:/GitHub/VisualIdeas/data/images/630ecc5198b8512a2ad8272038ed66bb.png}
\end{center}
\vspace{2mm}
\begin{center}
\begin{minipage}{0.92\textwidth}
\small
\textbf{Perception vs reality of routines}

I used to hate routines. But nothing gets you to love routines faster than having kids. 😆
\end{minipage}
\end{center}
\vspace*{\fill}
\newpage
\renewcommand{\theposttitle}{What learning the guitar taught me about progress}
\renewcommand{\thepostdate}{2024-06-18}
\vspace*{\fill}
\begin{center}
{\Huge\bfseries INTERESTING}\\[6mm]
\end{center}
\begin{center}
\includegraphics[width=\textwidth, height=0.55\textheight, keepaspectratio]{C:/GitHub/VisualIdeas/data/images/e4cdf7e00d13e62329d9b688265dc5a4.png}
\end{center}
\vspace{2mm}
\begin{center}
\begin{minipage}{0.92\textwidth}
\small
Most progress is invisible until it's not.

When I first learned to play the guitar—

E-chord? No problem. \\
C-chord? No problem. \\
But then I got to the G-chord, and I was like, what the hell is this?!

I thought, "how in the world do fingers make this shape?"

Day to day, I tried to play that freaking G-chord.

But no sound. All I felt was the pain of metal strings carving into my fingers. I couldn't see progress but for some reason I trusted the process.

And eventually one day out of no where, I heard \textit{strum}...\\
the sweet release of that G-chord! 🫠

That was the first time in my life I had proof that just because you can't see progress doesn't mean you're not making progress.
\end{minipage}
\end{center}
\vspace*{\fill}
\newpage
\renewcommand{\theposttitle}{What learning the guitar taught me about progress}
\renewcommand{\thepostdate}{2024-06-18}
\vspace*{\fill}
\begin{center}
{\Huge\bfseries ENCHANTING}\\[6mm]
\end{center}
\begin{center}
\includegraphics[width=\textwidth, height=0.55\textheight, keepaspectratio]{C:/GitHub/VisualIdeas/data/images/3b92ec2c227d466815fbca1a7ffabad1.png}
\end{center}
\vspace{2mm}
\begin{center}
\begin{minipage}{0.92\textwidth}
\small
\textbf{A critical shift in thinking to make for leaders. }
\end{minipage}
\end{center}
\vspace*{\fill}
\newpage
\renewcommand{\theposttitle}{What learning the guitar taught me about progress}
\renewcommand{\thepostdate}{2024-06-18}
\vspace*{\fill}
\begin{center}
{\Huge\bfseries ANALOGY}\\[6mm]
\end{center}
\begin{center}
\includegraphics[width=\textwidth, height=0.24\textheight, keepaspectratio]{C:/GitHub/VisualIdeas/data/images/ecf16e85deb3bf9c9e27c714dfc09374.png}
\end{center}
\vspace{2mm}
\begin{center}
\begin{minipage}{0.92\textwidth}
\small
\textbf{It's in your own hands.}

\textbf{Life is short. }
\end{minipage}
\end{center}
\vspace{5mm}
\begin{center}
\includegraphics[width=\textwidth, height=0.24\textheight, keepaspectratio]{C:/GitHub/VisualIdeas/data/images/3e4fee0486287afc4ec882cb8411c3d5.png}
\end{center}
\vspace*{\fill}
\newpage
\renewcommand{\theposttitle}{Visual IDEAs: Why making daily visuals is a superpower}
\renewcommand{\thepostdate}{2024-08-06}
\vspace*{\fill}
\begin{center}
{\Huge\bfseries INTERESTING}\\[6mm]
\end{center}
\begin{center}
\includegraphics[width=\textwidth, height=0.55\textheight, keepaspectratio]{C:/GitHub/VisualIdeas/data/images/5ae69c5bc648b9f5cd740f5e0080bb94.png}
\end{center}
\vspace{2mm}
\begin{center}
\begin{minipage}{0.92\textwidth}
\small
\textbf{It gets messier before it doesn't. }

One of the most beneficial aspects of writing is that it gives your memory and thinking structure. Random events and emotions can be organized with a beginning, middle and end.

But often, I find myself creating more chaos in the process before finding the thread to tie everything together.

Ah but the satisfaction and clarity once that baggage you’ve been carrying is organized feels oh so good! 😁

“This is the only real concern of the artist, \\
to recreate out of the disorder of life \\
that order which is art.”
\end{minipage}
\end{center}
\vspace*{\fill}
\newpage
\renewcommand{\theposttitle}{Visual IDEAs: Why making daily visuals is a superpower}
\renewcommand{\thepostdate}{2024-08-06}
\vspace*{\fill}
\begin{center}
{\Huge\bfseries DESIGN}\\[6mm]
\end{center}
\begin{center}
\begin{minipage}{0.92\textwidth}
\small
The beauty of creating visuals daily is the ability to \textbf{repurpose} the assets you build over time.

For example:
\end{minipage}
\end{center}
\vspace{4mm}
\begin{center}
\includegraphics[width=\textwidth, height=0.17\textheight, keepaspectratio]{C:/GitHub/VisualIdeas/data/images/755767468ae14c102ee8154d60a9e02f.jpeg}
\end{center}
\vspace{2mm}
\begin{center}
\begin{minipage}{0.92\textwidth}
\small
This is the 21st visual I ever made (very different than today I know 😅)
\end{minipage}
\end{center}
\vspace{5mm}
\begin{center}
\includegraphics[width=\textwidth, height=0.17\textheight, keepaspectratio]{C:/GitHub/VisualIdeas/data/images/cdfb206b7fab4ea04315384ede477878.png}
\end{center}
\vspace{2mm}
\begin{center}
\begin{minipage}{0.92\textwidth}
\small
Here’s that same firefly repurposed a year later!

This is a superpower because the more you create, the more you CAN create!

For example, this \textbf{recent carousel on the difference between amateurs vs professionals}\textbf{*} that went viral were mostly visuals I gave new life by repurposing them.

This means even though in the beginning you’re doing a lot of heavy lifting, the more reps you do, the bigger you grow your personal library of assets. And that means you don’t always have to start from scratch!

P.S. If you’re ready to start building your own visual library, you may want to join the next cohort of \textbf{Thinking in Visual Metaphors} in September. Be sure to sign up for the waitlist to be one of the first to know when registration starts!
\end{minipage}
\end{center}
\vspace{5mm}
\begin{center}
\includegraphics[width=\textwidth, height=0.17\textheight, keepaspectratio]{C:/GitHub/VisualIdeas/data/images/11e7b17799d0c82d0c1d0f04b2b1860e.png}
\end{center}
\vspace{2mm}
\begin{center}
\begin{minipage}{0.92\textwidth}
\small
*I should note to be an amateur is a beautiful thing. “Amateur” comes from the Latin word for lover (amator). I’m a proud amateur guitarist! 😬 The problem comes \textit{if I want to become a pro} guitarist. Pursuing it as a career forces changes in my mindset that need to happen if I want to avoid disappointment.
\end{minipage}
\end{center}
\vspace*{\fill}
\newpage
\renewcommand{\theposttitle}{Visual IDEAs: Why making daily visuals is a superpower}
\renewcommand{\thepostdate}{2024-08-06}
\vspace*{\fill}
\begin{center}
{\Huge\bfseries ENCHANTING}\\[6mm]
\end{center}
\begin{center}
\includegraphics[width=\textwidth, height=0.55\textheight, keepaspectratio]{C:/GitHub/VisualIdeas/data/images/d10b157208e23ee754caa35794bf3965.png}
\end{center}
\vspace{2mm}
\begin{center}
\begin{minipage}{0.92\textwidth}
\small
\textbf{The Swan Paradox: }\\
It takes an obscene amount of effort to make what you're doing look effortless.\\
To stay afloat, the effort never stops.

This is what it means to be "consistent."

But be careful. You'll drown if you only rely on burst efforts. \\
Pace yourself.
\end{minipage}
\end{center}
\vspace*{\fill}
\newpage
\renewcommand{\theposttitle}{Visual IDEAs: Why making daily visuals is a superpower}
\renewcommand{\thepostdate}{2024-08-06}
\vspace*{\fill}
\begin{center}
{\Huge\bfseries ANALOGY}\\[6mm]
\end{center}
\begin{center}
\includegraphics[width=\textwidth, height=0.55\textheight, keepaspectratio]{C:/GitHub/VisualIdeas/data/images/b43f0eee90b1f626551c065cb16051d5.png}
\end{center}
\vspace{2mm}
\begin{center}
\begin{minipage}{0.92\textwidth}
\small
\textbf{Definition of failure. }
\end{minipage}
\end{center}
\vspace*{\fill}
\newpage
\renewcommand{\theposttitle}{How to find happiness where you are}
\renewcommand{\thepostdate}{2024-08-20}
\vspace*{\fill}
\begin{center}
{\Huge\bfseries INTERESTING}\\[6mm]
\end{center}
\begin{center}
\includegraphics[width=\textwidth, height=0.55\textheight, keepaspectratio]{C:/GitHub/VisualIdeas/data/images/c272761b4b39b14bc97b6e36580a9cbe.png}
\end{center}
\vspace{2mm}
\begin{center}
\begin{minipage}{0.92\textwidth}
\small
\textbf{The one word never to include in an apology.}
\end{minipage}
\end{center}
\vspace*{\fill}
\newpage
\renewcommand{\theposttitle}{How to find happiness where you are}
\renewcommand{\thepostdate}{2024-08-20}
\vspace*{\fill}
\begin{center}
{\Huge\bfseries DESIGN}\\[6mm]
\end{center}
\begin{center}
\begin{minipage}{0.92\textwidth}
\small
I’m currently reading \textbf{Dr. Peter Attia’s book, Outlive}, which is a very thorough and detailed book summarizing most of Dr. Attia’s best practices.

The one thing I noticed reading the book though is that its jargon filled and could probably benefit from some visual thinking to cement the messages.

The example below is a chart Dr. Attia fought tooth and nail to put into the book but was overruled by his editors. Now I don’t know about you, but I can see why. I had to listen to the explanation twice and then spend another few minutes really looking over this chart to get the point.
\end{minipage}
\end{center}
\vspace{4mm}
\begin{center}
\includegraphics[width=\textwidth, height=0.24\textheight, keepaspectratio]{C:/GitHub/VisualIdeas/data/images/34dfbfbf68ffcffd6e56a4a8ba7d96c8.png}
\end{center}
\vspace{2mm}
\begin{center}
\begin{minipage}{0.92\textwidth}
\small
Today’s quick lesson if you are trying to convey a message in a visual:\\
\textbf{Stop presenting data. Instead, present the point of the data. }

For example, here’s ONE of the interpretations of the data (there are many other helpful interpretations that EACH deserve their own visual):
\end{minipage}
\end{center}
\vspace{5mm}
\begin{center}
\includegraphics[width=\textwidth, height=0.24\textheight, keepaspectratio]{C:/GitHub/VisualIdeas/data/images/206b9417657fc1bf2309b48d062be4ce.png}
\end{center}
\vspace{2mm}
\begin{center}
\begin{minipage}{0.92\textwidth}
\small
Based on Dr. Peter Attia’s Data

It’s true we lose some nuance from simplifying, but we also save our reader from spending as much time as I did to interpret the data.

You can always complement the visual with the hardcore data in an appendix (or if giving a live-presentation, a QR code to a website with the charts)

P.S. If you want to go deeper to learn how to create visual metaphors, \textbf{click here to sign up for the waitlist. }
\end{minipage}
\end{center}
\vspace*{\fill}
\newpage
\renewcommand{\theposttitle}{How to find happiness where you are}
\renewcommand{\thepostdate}{2024-08-20}
\vspace*{\fill}
\begin{center}
{\Huge\bfseries ENCHANTING}\\[6mm]
\end{center}
\begin{center}
\includegraphics[width=\textwidth, height=0.55\textheight, keepaspectratio]{C:/GitHub/VisualIdeas/data/images/3f66053afd01580c3f08eac02a873fed.png}
\end{center}
\vspace{2mm}
\begin{center}
\begin{minipage}{0.92\textwidth}
\small
\textbf{The pursuit of having what you want vs wanting what you have.}

It’s all too easy to miss all the things you have right where you are.
\end{minipage}
\end{center}
\vspace*{\fill}
\newpage
\renewcommand{\theposttitle}{How to find happiness where you are}
\renewcommand{\thepostdate}{2024-08-20}
\vspace*{\fill}
\begin{center}
{\Huge\bfseries ANALOGY}\\[6mm]
\end{center}
\begin{center}
\includegraphics[width=\textwidth, height=0.55\textheight, keepaspectratio]{C:/GitHub/VisualIdeas/data/images/d3b0b6c5904017573a85006d404c3019.png}
\end{center}
\vspace{2mm}
\begin{center}
\begin{minipage}{0.92\textwidth}
\small
\textbf{The distance between your dreams and reality is not as far as you think.}

I overplan. Overanalyze. Overthink most things. It may be why my \textbf{overthinking visual carousel }struck a chord with so many folks a couple weeks ago.

Before I made my first online post, I kept thinking I had to figure everything out before starting.

Questions like:

What's my niche?

Who am I to talk about this?

And hasn't everything that needs to be said about everything already been said?!

So I kept thinking...and waiting...for over 10+ years!

The biggest lesson I learned after starting was you can't steer a parked car (as the saying goes).

99\% of that analysis paralysis was just in my head.
\end{minipage}
\end{center}
\vspace*{\fill}
\newpage
\renewcommand{\theposttitle}{How to understand more deeply}
\renewcommand{\thepostdate}{2024-09-03}
\vspace*{\fill}
\begin{center}
{\Huge\bfseries ENCHANTING}\\[6mm]
\end{center}
\begin{center}
\includegraphics[width=\textwidth, height=0.55\textheight, keepaspectratio]{C:/GitHub/VisualIdeas/data/images/9b252e1d36c5774778e0063beb8ce7b2.png}
\end{center}
\vspace{2mm}
\begin{center}
\begin{minipage}{0.92\textwidth}
\small
Everything you've ever wanted is sitting on the other side of fear.

\textbf{Let's walk through that curtain.}
\end{minipage}
\end{center}
\vspace*{\fill}
\newpage
\renewcommand{\theposttitle}{How to understand more deeply}
\renewcommand{\thepostdate}{2024-09-03}
\vspace*{\fill}
\begin{center}
{\Huge\bfseries ANALOGY}\\[6mm]
\end{center}
\begin{center}
\includegraphics[width=\textwidth, height=0.55\textheight, keepaspectratio]{C:/GitHub/VisualIdeas/data/images/6072056478874332b36e4fd95207e563.png}
\end{center}
\vspace{2mm}
\begin{center}
\begin{minipage}{0.92\textwidth}
\small
When you talk, you are only repeating

what you already know.

But if you listen, you may learn something new.

Inspired by Carolyn Coughlin, co-founder of Cultivating Leadership
\end{minipage}
\end{center}
\vspace*{\fill}
\newpage
\renewcommand{\theposttitle}{Why you should embrace FOMO}
\renewcommand{\thepostdate}{2024-09-17}
\vspace*{\fill}
\begin{center}
{\Huge\bfseries DESIGN}\\[6mm]
\end{center}
\begin{center}
\begin{minipage}{0.92\textwidth}
\small
A quick tip to simplify a metaphor:\\
Play with the \textbf{AMOUNT} of an element in the visual.

It’s a helpful way to frame an idea so its meaning is more easily understood by the reader.

To start, try this:
\end{minipage}
\end{center}
\vspace{4mm}
\begin{center}
\includegraphics[width=\textwidth, height=0.24\textheight, keepaspectratio]{C:/GitHub/VisualIdeas/data/images/bb9c3411777c15a7bf4ff67806015ccb.png}
\end{center}
\vspace{2mm}
\begin{center}
\begin{minipage}{0.92\textwidth}
\small
Multiply the quantity of the object…
\end{minipage}
\end{center}
\vspace{5mm}
\begin{center}
\includegraphics[width=\textwidth, height=0.24\textheight, keepaspectratio]{C:/GitHub/VisualIdeas/data/images/f46edc8abf8f86f01ae85681802fc6a5.png}
\end{center}
\vspace{2mm}
\begin{center}
\begin{minipage}{0.92\textwidth}
\small
…or scale its size.

P.S. When you’re ready to start creating your own visuals, \textbf{sign up for the next cohort of Thinking in Visual Metaphors ↗} .
\end{minipage}
\end{center}
\vspace*{\fill}
\newpage
\renewcommand{\theposttitle}{Why you should embrace FOMO}
\renewcommand{\thepostdate}{2024-09-17}
\vspace*{\fill}
\begin{center}
{\Huge\bfseries ENCHANTING}\\[6mm]
\end{center}
\begin{center}
\includegraphics[width=\textwidth, height=0.55\textheight, keepaspectratio]{C:/GitHub/VisualIdeas/data/images/79488b2225fb6784f879b6cbbe80c608.png}
\end{center}
\vspace{2mm}
\begin{center}
\begin{minipage}{0.92\textwidth}
\small
It's not what you gather, but what you scatter that tells what kind of life you have lived.

No one will remember how much you made. \\
But the people that matter most will remember your integrity.
\end{minipage}
\end{center}
\vspace*{\fill}
\newpage
\renewcommand{\theposttitle}{Why you should embrace FOMO}
\renewcommand{\thepostdate}{2024-09-17}
\vspace*{\fill}
\begin{center}
{\Huge\bfseries ANALOGY}\\[6mm]
\end{center}
\begin{center}
\includegraphics[width=\textwidth, height=0.55\textheight, keepaspectratio]{C:/GitHub/VisualIdeas/data/images/504cd821383602c220a0c9f984868b3f.png}
\end{center}
\vspace{2mm}
\begin{center}
\begin{minipage}{0.92\textwidth}
\small
Missing out’ is what makes our choices meaningful.

One of my favorite quotes from one of my favorite books, \textbf{“4000 Weeks: Time Management for Mortals” by Oliver Burkeman. }

Embrace the joy of missing out by focusing on what truly matters and saving your energy for what counts.
\end{minipage}
\end{center}
\vspace*{\fill}
\newpage
\renewcommand{\theposttitle}{Why there's no shortcut to getting wiser}
\renewcommand{\thepostdate}{2024-10-01}
\vspace*{\fill}
\begin{center}
{\Huge\bfseries INTERESTING}\\[6mm]
\end{center}
\begin{center}
\includegraphics[width=\textwidth, height=0.55\textheight, keepaspectratio]{C:/GitHub/VisualIdeas/data/images/61b25f9cc21dc98f828d34e41261ac0e.png}
\end{center}
\vspace{2mm}
\begin{center}
\begin{minipage}{0.92\textwidth}
\small
Mistakes are the portals of discovery.

A key lesson I've learned after teaching high school for over 20 years:\\
\textbf{You can’t shortcut the experience that comes through trial \& error.}

Knowledge is learned. Wisdom is earned.
\end{minipage}
\end{center}
\vspace*{\fill}
\newpage
\renewcommand{\theposttitle}{Why there's no shortcut to getting wiser}
\renewcommand{\thepostdate}{2024-10-01}
\vspace*{\fill}
\begin{center}
{\Huge\bfseries DESIGN}\\[6mm]
\end{center}
\begin{center}
\begin{minipage}{0.92\textwidth}
\small
If you’ve been following me for a while, you know to \textbf{focus on iteration over perfection.}

Knowing you can reverse a decision is an antidote to overthinking.

Iterations can be BIG. Or they can be small. The point is you can \textbf{perfect in public. }

Here’s an example of a small iteration:
\end{minipage}
\end{center}
\vspace{4mm}
\begin{center}
\includegraphics[width=\textwidth, height=0.24\textheight, keepaspectratio]{C:/GitHub/VisualIdeas/data/images/c732bc52667602096ee0894ed1509a18.png}
\end{center}
\vspace{2mm}
\begin{center}
\begin{minipage}{0.92\textwidth}
\small
Version 1: Published 3/20/24
\end{minipage}
\end{center}
\vspace{5mm}
\begin{center}
\includegraphics[width=\textwidth, height=0.24\textheight, keepaspectratio]{C:/GitHub/VisualIdeas/data/images/c07cd5fce9ca5b81ff9380f52775b6bf.png}
\end{center}
\vspace{2mm}
\begin{center}
\begin{minipage}{0.92\textwidth}
\small
Version 2: Published 9/29/24

Login or Subscribe to participate

The visual is far from perfect. But the beauty (and irritating thing 😆) about social media is that people forget!

In a few months, you can make a significant or subtle change and try again.
\end{minipage}
\end{center}
\vspace*{\fill}
\newpage
\renewcommand{\theposttitle}{Why there's no shortcut to getting wiser}
\renewcommand{\thepostdate}{2024-10-01}
\vspace*{\fill}
\begin{center}
{\Huge\bfseries ENCHANTING}\\[6mm]
\end{center}
\begin{center}
\includegraphics[width=\textwidth, height=0.55\textheight, keepaspectratio]{C:/GitHub/VisualIdeas/data/images/75e01585917ca476866eec3e4de50eaa.png}
\end{center}
\vspace{2mm}
\begin{center}
\begin{minipage}{0.92\textwidth}
\small
\textbf{You can't pour from an empty cup.}\\
When you take care of yourself, you’re better able to help others.
\end{minipage}
\end{center}
\vspace*{\fill}
\newpage
\renewcommand{\theposttitle}{Why there's no shortcut to getting wiser}
\renewcommand{\thepostdate}{2024-10-01}
\vspace*{\fill}
\begin{center}
{\Huge\bfseries ANALOGY}\\[6mm]
\end{center}
\begin{center}
\includegraphics[width=\textwidth, height=0.55\textheight, keepaspectratio]{C:/GitHub/VisualIdeas/data/images/1ddc562f6df45b23634a5d064f3ece99.png}
\end{center}
\vspace{2mm}
\begin{center}
\begin{minipage}{0.92\textwidth}
\small
"Success is not to be pursued; it is to be attracted by the person you become."

\textbf{Focus on your roots and success will blossom.}
\end{minipage}
\end{center}
\vspace*{\fill}
\newpage
\renewcommand{\theposttitle}{What tea can teach us about overthinking}
\renewcommand{\thepostdate}{2024-10-15}
\vspace*{\fill}
\begin{center}
{\Huge\bfseries INTERESTING}\\[6mm]
\end{center}
\begin{center}
\includegraphics[width=\textwidth, height=0.55\textheight, keepaspectratio]{C:/GitHub/VisualIdeas/data/images/3653bb3f6faaccaf8f34032f1d2ff363.png}
\end{center}
\vspace{2mm}
\begin{center}
\begin{minipage}{0.92\textwidth}
\small
\textbf{The sweet spot for how much to think before acting.}\\
(this may only make sense if you drink tea 😆)
\end{minipage}
\end{center}
\vspace*{\fill}
\newpage
\renewcommand{\theposttitle}{What tea can teach us about overthinking}
\renewcommand{\thepostdate}{2024-10-15}
\vspace*{\fill}
\begin{center}
{\Huge\bfseries DESIGN}\\[6mm]
\end{center}
\begin{center}
\begin{minipage}{0.92\textwidth}
\small
Last week, I ran a poll on 2 versions of a visual:
\end{minipage}
\end{center}
\vspace{4mm}
\begin{center}
\includegraphics[width=\textwidth, height=0.55\textheight, keepaspectratio]{C:/GitHub/VisualIdeas/data/images/24dfcfcb99de04d8cd1cfce6310a6f0b.png}
\end{center}
\vspace{2mm}
\begin{center}
\begin{minipage}{0.92\textwidth}
\small
\textbf{What’s the lesson?}

For my fellow recovering perfectionists, it’s difficult to know what works until you share your work. When you get caught in the weeds of messing with subtle details, it doesn’t always mean you’re making an improvement! (a humbling lesson)

Back to the drawing board! 😆
\end{minipage}
\end{center}
\vspace*{\fill}
\newpage
\renewcommand{\theposttitle}{What tea can teach us about overthinking}
\renewcommand{\thepostdate}{2024-10-15}
\vspace*{\fill}
\begin{center}
{\Huge\bfseries ENCHANTING}\\[6mm]
\end{center}
\begin{center}
\includegraphics[width=\textwidth, height=0.55\textheight, keepaspectratio]{C:/GitHub/VisualIdeas/data/images/6e5a6891b641419163fa291467e87ae7.png}
\end{center}
\vspace{2mm}
\begin{center}
\begin{minipage}{0.92\textwidth}
\small
A single act of kindness throws out roots in all directions, and the roots spring up and make new trees.
\end{minipage}
\end{center}
\vspace*{\fill}
\newpage
\renewcommand{\theposttitle}{What tea can teach us about overthinking}
\renewcommand{\thepostdate}{2024-10-15}
\vspace*{\fill}
\begin{center}
{\Huge\bfseries ANALOGY}\\[6mm]
\end{center}
\begin{center}
\includegraphics[width=\textwidth, height=0.55\textheight, keepaspectratio]{C:/GitHub/VisualIdeas/data/images/db30550bb50e4ed15838165603c484eb.png}
\end{center}
\vspace{2mm}
\begin{center}
\begin{minipage}{0.92\textwidth}
\small
Skill is only developed by hours and hours of work.
\end{minipage}
\end{center}
\vspace*{\fill}
\newpage
\renewcommand{\theposttitle}{Are you reaching your full potential?}
\renewcommand{\thepostdate}{2024-10-29}
\vspace*{\fill}
\begin{center}
{\Huge\bfseries INTERESTING}\\[6mm]
\end{center}
\begin{center}
\includegraphics[width=\textwidth, height=0.24\textheight, keepaspectratio]{C:/GitHub/VisualIdeas/data/images/e2279be7724c25f806a8a3882e69d1e4.png}
\end{center}
\vspace{2mm}
\begin{center}
\begin{minipage}{0.92\textwidth}
\small
You do not rise to the level of your goals. \\
You fall to the level of your systems.

One of my favorite quotes from the book, \textbf{“Atomic Habits.”}
\end{minipage}
\end{center}
\vspace{5mm}
\begin{center}
\includegraphics[width=\textwidth, height=0.24\textheight, keepaspectratio]{C:/GitHub/VisualIdeas/data/images/555ab948b93e67ef2bdb9a54be662152.jpg}
\end{center}
\vspace{2mm}
\begin{center}
\begin{minipage}{0.92\textwidth}
\small
Join over 4 million Americans who start their day with 1440 – your daily digest for unbiased, fact-centric news. From politics to sports, we cover it all by analyzing over 100 sources. Our concise, 5-minute read lands in your inbox each morning at no cost. Experience news without the noise; let 1440 help you make up your own mind. Sign up now and invite your friends and family to be part of the informed.

Sign up today!
\end{minipage}
\end{center}
\vspace*{\fill}
\newpage
\renewcommand{\theposttitle}{Are you reaching your full potential?}
\renewcommand{\thepostdate}{2024-10-29}
\vspace*{\fill}
\begin{center}
{\Huge\bfseries DESIGN}\\[6mm]
\end{center}
\begin{center}
\begin{minipage}{0.92\textwidth}
\small
Recently, I shared a first animation of one of my visual metaphors. There’s plenty of room to grow and experiment in this space, but it did make me reflect…what is lost and what is gained in the reader’s experience when we add motion?
\end{minipage}
\end{center}
\vspace{4mm}
\begin{center}
%\includegraphics[width=\textwidth, height=0.55\textheight, keepaspectratio]{C:/GitHub/VisualIdeas/data/images/29c4478b0fea5695857346841ecd1f24.gif}
\end{center}
\vspace{2mm}
\begin{center}
\begin{minipage}{0.92\textwidth}
\small
Login or Subscribe to participate

Or shoot a reply to this email if you have more thoughts to add.
\end{minipage}
\end{center}
\vspace*{\fill}
\newpage
\renewcommand{\theposttitle}{Are you reaching your full potential?}
\renewcommand{\thepostdate}{2024-10-29}
\vspace*{\fill}
\begin{center}
{\Huge\bfseries ENCHANTING}\\[6mm]
\end{center}
\begin{center}
\includegraphics[width=\textwidth, height=0.55\textheight, keepaspectratio]{C:/GitHub/VisualIdeas/data/images/91895fea85321b431cb13b85407009c7.png}
\end{center}
\vspace{2mm}
\begin{center}
\begin{minipage}{0.92\textwidth}
\small
\textbf{Motion = staying busy }\\
\textbf{Action = focused execution}

I used to wear “being busy” as a badge of honor. But a lesson I’m relearning is doing a lot and doing what matters are not the same thing.

What’s worse is the busier I get, the less I focus on the latter.

Don’t make this same mistake. Create space for a few moments today to think about how you might position this week to take action on the things that matter most.
\end{minipage}
\end{center}
\vspace*{\fill}
\newpage
\renewcommand{\theposttitle}{Are you reaching your full potential?}
\renewcommand{\thepostdate}{2024-10-29}
\vspace*{\fill}
\begin{center}
{\Huge\bfseries ANALOGY}\\[6mm]
\end{center}
\begin{center}
\includegraphics[width=\textwidth, height=0.55\textheight, keepaspectratio]{C:/GitHub/VisualIdeas/data/images/4c03ea19980b50b3818f3bddea921e89.png}
\end{center}
\vspace{2mm}
\begin{center}
\begin{minipage}{0.92\textwidth}
\small
Leadership is not about being in charge. \\
It’s about taking care of those in your charge.
\end{minipage}
\end{center}
\vspace*{\fill}
\newpage
\renewcommand{\theposttitle}{The art of remixing ideas}
\renewcommand{\thepostdate}{2024-11-12}
\vspace*{\fill}
\begin{center}
{\Huge\bfseries INTERESTING}\\[6mm]
\end{center}
\begin{center}
\includegraphics[width=\textwidth, height=0.55\textheight, keepaspectratio]{C:/GitHub/VisualIdeas/data/images/04b8e093afdbfbde13e2415e4f7923ef.png}
\end{center}
\vspace{2mm}
\begin{center}
\begin{minipage}{0.92\textwidth}
\small
\textbf{Speaking to impress vs listening to understand.}
\end{minipage}
\end{center}
\vspace*{\fill}
\newpage
\renewcommand{\theposttitle}{The art of remixing ideas}
\renewcommand{\thepostdate}{2024-11-12}
\vspace*{\fill}
\begin{center}
{\Huge\bfseries ANALOGY}\\[6mm]
\end{center}
\begin{center}
\includegraphics[width=\textwidth, height=0.55\textheight, keepaspectratio]{C:/GitHub/VisualIdeas/data/images/3943cb56e0fc88deaf54286864b250d3.png}
\end{center}
\vspace{2mm}
\begin{center}
\begin{minipage}{0.92\textwidth}
\small
When you're surrounded by the right people, \\
hours feel like minutes.

When you're surrounded by the wrong people, \\
minutes feel like hours.

Login or Subscribe to participate
\end{minipage}
\end{center}
\vspace*{\fill}
\newpage
\renewcommand{\theposttitle}{The problem with low-hanging fruit}
\renewcommand{\thepostdate}{2024-11-26}
\vspace*{\fill}
\begin{center}
{\Huge\bfseries INTERESTING}\\[6mm]
\end{center}
\begin{center}
\includegraphics[width=\textwidth, height=0.55\textheight, keepaspectratio]{C:/GitHub/VisualIdeas/data/images/ba7397c58347b3031304748f055dc731.png}
\end{center}
\vspace{2mm}
\begin{center}
\begin{minipage}{0.92\textwidth}
\small
A mistake a lot of folks make:

Falling into "the collector's fallacy" trap--\\
Believing that by collecting more information, we somehow absorb it to grow.

It feels productive to:

highlight more books

bookmark more articles

listen to more podcasts

“If more information was the answer, then we'd all be billionaires with perfect abs.”

Real progress happens through decisions, action and commitment.
\end{minipage}
\end{center}
\vspace*{\fill}
\newpage
\renewcommand{\theposttitle}{The problem with low-hanging fruit}
\renewcommand{\thepostdate}{2024-11-26}
\vspace*{\fill}
\begin{center}
{\Huge\bfseries ENCHANTING}\\[6mm]
\end{center}
\begin{center}
\includegraphics[width=\textwidth, height=0.55\textheight, keepaspectratio]{C:/GitHub/VisualIdeas/data/images/b00a1548d330ba327803eb0a29582711.png}
\end{center}
\vspace{2mm}
\begin{center}
\begin{minipage}{0.92\textwidth}
\small
“If we only do what we’re familiar with, we might miss what we’ve been made for.”
\end{minipage}
\end{center}
\vspace*{\fill}
\newpage
\renewcommand{\theposttitle}{The problem with low-hanging fruit}
\renewcommand{\thepostdate}{2024-11-26}
\vspace*{\fill}
\begin{center}
{\Huge\bfseries ANALOGY}\\[6mm]
\end{center}
\begin{center}
\includegraphics[width=\textwidth, height=0.55\textheight, keepaspectratio]{C:/GitHub/VisualIdeas/data/images/b895e36605db1947861ca587163275b7.png}
\end{center}
\vspace{2mm}
\begin{center}
\begin{minipage}{0.92\textwidth}
\small
\textbf{Friendly reminder:} Always check you’re climbing the right mountain.
\end{minipage}
\end{center}
\vspace*{\fill}
\newpage
\renewcommand{\theposttitle}{Building a creative habit in 2025}
\renewcommand{\thepostdate}{2024-12-10}
\vspace*{\fill}
\begin{center}
{\Huge\bfseries INTERESTING}\\[6mm]
\end{center}
\begin{center}
\includegraphics[width=\textwidth, height=0.55\textheight, keepaspectratio]{C:/GitHub/VisualIdeas/data/images/1f461a4e51b14ffcdfe9a0f758ff9b37.png}
\end{center}
\vspace{2mm}
\begin{center}
\begin{minipage}{0.92\textwidth}
\small
One of my favorite idea learned this year from Adam Grant: \\
the word "\textbf{sonder"} from the \textbf{Dictionary of Obscure Sorrows by John Koenig}

Some of my other favorites from the dictionary:

\textbf{HARMONIA}\\
n.an itchy sense of dread when life feels just a hint too peaceful—when everyone seems to get along suspiciously well, with an eerie stillness that makes you want to brace for the inevitable collapse, or burn it down yourself.\\
(I feel like I’m in an unhealthy constant state of harmonia. Am I the only one? 🤔)

\textbf{ANCHORAGE}\\
n.the desire to hold on to time as it passes, like trying to keep your grip on a rock in the middle of a river, feeling the weight of the current against your chest while your elders float on downstream, calling over the roar of the rapids, “Just let go—it’s okay—let go.”
\end{minipage}
\end{center}
\vspace*{\fill}
\newpage
\renewcommand{\theposttitle}{Building a creative habit in 2025}
\renewcommand{\thepostdate}{2024-12-10}
\vspace*{\fill}
\begin{center}
{\Huge\bfseries DESIGN}\\[6mm]
\end{center}
\begin{center}
\begin{minipage}{0.92\textwidth}
\small
Building a creative practice isn’t just about making illustrations. Here’s a short highlight from my chat with Roberto Ferraro on the value of starting and keeping a consistent creative habit.
\end{minipage}
\end{center}
\vspace{4mm}
\begin{center}
\includegraphics[width=\textwidth, height=0.55\textheight, keepaspectratio]{C:/GitHub/VisualIdeas/data/images/dcebf335927442da25662d3df825dfed.png}
\end{center}
\vspace{2mm}
\begin{center}
\begin{minipage}{0.92\textwidth}
\small
The start of a new year is a perfect time to dive into making stuff of your own!
\end{minipage}
\end{center}
\vspace*{\fill}
\newpage
\renewcommand{\theposttitle}{Building a creative habit in 2025}
\renewcommand{\thepostdate}{2024-12-10}
\vspace*{\fill}
\begin{center}
{\Huge\bfseries ENCHANTING}\\[6mm]
\end{center}
\begin{center}
\includegraphics[width=\textwidth, height=0.55\textheight, keepaspectratio]{C:/GitHub/VisualIdeas/data/images/33867e956aa95ebd6c768779a125c99f.png}
\end{center}
\vspace{2mm}
\begin{center}
\begin{minipage}{0.92\textwidth}
\small
\textbf{There's a reason why dough must rest before it can rise.}

This is something I'm still working on.

Resting may seem commonsensical, but when you've built habits around rolling forward at full momentum, it can be tough slowing down.

Earlier this year, I scheduled a weekend to relax, but I was shocked how little I could shut off without feeling guilty or anxious I should be doing more. It wasn't until last month when I got sick I realized there's only so far you can go if you don't learn to rest.
\end{minipage}
\end{center}
\vspace*{\fill}
\newpage
\renewcommand{\theposttitle}{Building a creative habit in 2025}
\renewcommand{\thepostdate}{2024-12-10}
\vspace*{\fill}
\begin{center}
{\Huge\bfseries ANALOGY}\\[6mm]
\end{center}
\begin{center}
\includegraphics[width=\textwidth, height=0.55\textheight, keepaspectratio]{C:/GitHub/VisualIdeas/data/images/3e6774281f7367e1bf846208d417dfcc.png}
\end{center}
\vspace{2mm}
\begin{center}
\begin{minipage}{0.92\textwidth}
\small
The difference between thinking and overthinking.
\end{minipage}
\end{center}
\vspace*{\fill}
\newpage
\renewcommand{\theposttitle}{Putting it all into perspective}
\renewcommand{\thepostdate}{2024-12-27}
\vspace*{\fill}
\begin{center}
{\Huge\bfseries INTERESTING}\\[6mm]
\end{center}
\begin{center}
\includegraphics[width=\textwidth, height=0.55\textheight, keepaspectratio]{C:/GitHub/VisualIdeas/data/images/3af18e080df5ebfafe96ac38e039d4d6.png}
\end{center}
\vspace{2mm}
\begin{center}
\begin{minipage}{0.92\textwidth}
\small
If you're a recovering perfectionist, remember this: daily 1\% improvement feels slow, but it's far faster than waiting for your work to be perfect before you share.

\textbf{Read the full post here.}
\end{minipage}
\end{center}
\vspace*{\fill}
\newpage
\renewcommand{\theposttitle}{Putting it all into perspective}
\renewcommand{\thepostdate}{2024-12-27}
\vspace*{\fill}
\begin{center}
{\Huge\bfseries DESIGN}\\[6mm]
\end{center}
\begin{center}
\begin{minipage}{0.92\textwidth}
\small
Why it’s dangerous to measure your success through the lens of social media growth and how to look at it instead:
\end{minipage}
\end{center}
\vspace{4mm}
\begin{center}
\includegraphics[width=\textwidth, height=0.55\textheight, keepaspectratio]{C:/GitHub/VisualIdeas/data/images/23e0c29803b5e6cb185ffe3b7b7666ed.png}
\end{center}
\vspace{2mm}
\begin{center}
\begin{minipage}{0.92\textwidth}
\small
P.S. If you’re ready for a new beginning, enrollment for Cohort 9 (January) of Thinking in Visual Metaphors is now open.
\end{minipage}
\end{center}
\vspace*{\fill}
\newpage
\renewcommand{\theposttitle}{Putting it all into perspective}
\renewcommand{\thepostdate}{2024-12-27}
\vspace*{\fill}
\begin{center}
{\Huge\bfseries ENCHANTING}\\[6mm]
\end{center}
\begin{center}
\includegraphics[width=\textwidth, height=0.55\textheight, keepaspectratio]{C:/GitHub/VisualIdeas/data/images/d13d2644bdafa7b183c6b74a49b9735d.png}
\end{center}
\vspace{2mm}
\begin{center}
\begin{minipage}{0.92\textwidth}
\small
The past in perspective.
\end{minipage}
\end{center}
\vspace*{\fill}
\newpage
\renewcommand{\theposttitle}{Putting it all into perspective}
\renewcommand{\thepostdate}{2024-12-27}
\vspace*{\fill}
\begin{center}
{\Huge\bfseries ANALOGY}\\[6mm]
\end{center}
\begin{center}
\includegraphics[width=\textwidth, height=0.55\textheight, keepaspectratio]{C:/GitHub/VisualIdeas/data/images/5c4db4ad69ef5db01064cc77d59e0ac6.png}
\end{center}
\vspace{2mm}
\begin{center}
\begin{minipage}{0.92\textwidth}
\small
The most interesting thing about a fixed vs growth mindset is where and how it's applied.

\textbf{Example:} We can be growth minded about getting good at basketball, but fixed in our math ability. We each have self-belief blind spots where we see things as fixed.

The good news is self-awareness of these spots is a useful first step to getting closer to where we want to be.
\end{minipage}
\end{center}
\vspace*{\fill}
\newpage
\renewcommand{\theposttitle}{How happiness is like a butterfly}
\renewcommand{\thepostdate}{2025-01-08}
\vspace*{\fill}
\begin{center}
{\Huge\bfseries INTERESTING}\\[6mm]
\end{center}
\begin{center}
\includegraphics[width=\textwidth, height=0.55\textheight, keepaspectratio]{C:/GitHub/VisualIdeas/data/images/d66cad85795e029ea2089f2057b5b808.png}
\end{center}
\vspace{2mm}
\begin{center}
\begin{minipage}{0.92\textwidth}
\small
There is no one-size fits all formula for work-life balance.

It’s personal to you. A moving target based on the stage of your life.

Arrange and rearrange the pieces until you feel aligned.

And remember, you may need to reshape the balance as you shift from one stage of life to another.

What felt right in your 20s will most likely not fit where you are in your 40s.
\end{minipage}
\end{center}
\vspace*{\fill}
\newpage
\renewcommand{\theposttitle}{How happiness is like a butterfly}
\renewcommand{\thepostdate}{2025-01-08}
\vspace*{\fill}
\begin{center}
{\Huge\bfseries DESIGN}\\[6mm]
\end{center}
\begin{center}
\includegraphics[width=\textwidth, height=0.55\textheight, keepaspectratio]{C:/GitHub/VisualIdeas/data/images/b85524e1c6a3ca5cf0150bb343efa023.png}
\end{center}
\vspace{2mm}
\begin{center}
\begin{minipage}{0.92\textwidth}
\small
Social media today is flooded everywhere you look with AI writing and images. In this short highlight from my chat with Roberto Ferraro, I talk about the danger of using AI too early in your creative career.

By the way, if you’re interested in how to make visuals yourself, it’s not too late to sign up for Cohort 9.
\end{minipage}
\end{center}
\vspace*{\fill}
\newpage
\renewcommand{\theposttitle}{How happiness is like a butterfly}
\renewcommand{\thepostdate}{2025-01-08}
\vspace*{\fill}
\begin{center}
{\Huge\bfseries ENCHANTING}\\[6mm]
\end{center}
\begin{center}
\includegraphics[width=\textwidth, height=0.55\textheight, keepaspectratio]{C:/GitHub/VisualIdeas/data/images/edc6e8ad231e1dea4bcab87f569ac05e.png}
\end{center}
\vspace{2mm}
\begin{center}
\begin{minipage}{0.92\textwidth}
\small
Happiness is like a butterfly. The more you chase it, the more it evades you.

But if you notice the other things around you, it will gently come and sit on your shoulder.
\end{minipage}
\end{center}
\vspace*{\fill}
\newpage
\renewcommand{\theposttitle}{How happiness is like a butterfly}
\renewcommand{\thepostdate}{2025-01-08}
\vspace*{\fill}
\begin{center}
{\Huge\bfseries ANALOGY}\\[6mm]
\end{center}
\begin{center}
\includegraphics[width=\textwidth, height=0.55\textheight, keepaspectratio]{C:/GitHub/VisualIdeas/data/images/e0a3d65a443cd964153807f543e7a5df.png}
\end{center}
\vspace{2mm}
\begin{center}
\begin{minipage}{0.92\textwidth}
\small
\textbf{Be sure you're where you need to be to meet your potential.}
\end{minipage}
\end{center}
\vspace*{\fill}
\newpage
\renewcommand{\theposttitle}{What does it mean to have character?}
\renewcommand{\thepostdate}{2025-01-21}
\vspace*{\fill}
\begin{center}
{\Huge\bfseries INTERESTING}\\[6mm]
\end{center}
\begin{center}
\includegraphics[width=\textwidth, height=0.55\textheight, keepaspectratio]{C:/GitHub/VisualIdeas/data/images/42e4b6034d27649c388dbfcaee01c0f1.png}
\end{center}
\vspace{2mm}
\begin{center}
\begin{minipage}{0.92\textwidth}
\small
\textbf{A hit or a miss doesn’t define your character.}\\
\textbf{It’s how often you come back to keep shooting. }
\end{minipage}
\end{center}
\vspace*{\fill}
\newpage
\renewcommand{\theposttitle}{What does it mean to have character?}
\renewcommand{\thepostdate}{2025-01-21}
\vspace*{\fill}
\begin{center}
{\Huge\bfseries DESIGN}\\[6mm]
\end{center}
\begin{center}
\includegraphics[width=\textwidth, height=0.24\textheight, keepaspectratio]{C:/GitHub/VisualIdeas/data/images/269d94216d16857df501012224f5b697.png}
\end{center}
\vspace{2mm}
\begin{center}
\begin{minipage}{0.92\textwidth}
\small
Imagine if you could write, record, and make meaningful work on demand. The thing is most people have ideas worth sharing, but the content they consume poisons their creative well.

You may have heard of the famous Jim Rohn quote: “You’re the average of the 5 people you spend most thing with.” I’d add that your work is the average of the 5 types of content you most often consume.

By the way, if you’re interested in how to make visuals yourself, be one of the first to know when Cohort 10 is open for enrollment.
\end{minipage}
\end{center}
\vspace{5mm}
\begin{center}
\includegraphics[width=\textwidth, height=0.24\textheight, keepaspectratio]{C:/GitHub/VisualIdeas/data/images/26f060aeca0b6171044788ecddda76c3.png}
\end{center}
\vspace{2mm}
\begin{center}
\begin{minipage}{0.92\textwidth}
\small
\textbf{WHAT STUDENTS ARE SAYING:}\\
\textbf{“It’s a transformative course, and I couldn’t recommend it enough!” }\\
-Eva

\textbf{“I can see how much effort has gone into creating it and the time PJ has taken to be available to his students, both on live calls and in between.”} \\
-Rosie

\textbf{“I can’t believe that after only 3 weeks my brain is already starting to think differently about how to more clearly communicate information.” }\\
-Jen

\textbf{“This was the perfect course for me, at the right time in my life. }\\
\textbf{I only wish I had found it sooner.”}\\
-Kaamna

\textbf{Registration for the next cohort starts soon}\\
\textbf{↳ Translate your ideas into visuals!}
\end{minipage}
\end{center}
\vspace*{\fill}
\newpage
\renewcommand{\theposttitle}{What does it mean to have character?}
\renewcommand{\thepostdate}{2025-01-21}
\vspace*{\fill}
\begin{center}
{\Huge\bfseries ENCHANTING}\\[6mm]
\end{center}
\begin{center}
\includegraphics[width=\textwidth, height=0.55\textheight, keepaspectratio]{C:/GitHub/VisualIdeas/data/images/eeca2daa34f72b1604ca2686d355ac36.png}
\end{center}
\vspace{2mm}
\begin{center}
\begin{minipage}{0.92\textwidth}
\small
We are more alike, my friends, than we are unalike.
\end{minipage}
\end{center}
\vspace*{\fill}
\newpage
\renewcommand{\theposttitle}{What does it mean to have character?}
\renewcommand{\thepostdate}{2025-01-21}
\vspace*{\fill}
\begin{center}
{\Huge\bfseries ANALOGY}\\[6mm]
\end{center}
\begin{center}
\includegraphics[width=\textwidth, height=0.55\textheight, keepaspectratio]{C:/GitHub/VisualIdeas/data/images/31ec9b4524a0bc81fa2cb7a3cf3f87a4.png}
\end{center}
\vspace{2mm}
\begin{center}
\begin{minipage}{0.92\textwidth}
\small
\textbf{Definition of Failure}

For every idea shared, dozens never see the light of day.

Messing up and failing isn't a sign you're doing it wrong. It's a lesson to help you course correct.
\end{minipage}
\end{center}
\vspace*{\fill}
\newpage
\renewcommand{\theposttitle}{Why friendships are like trees}
\renewcommand{\thepostdate}{2025-02-04}
\vspace*{\fill}
\begin{center}
{\Huge\bfseries DESIGN}\\[6mm]
\end{center}
\begin{center}
\includegraphics[width=\textwidth, height=0.55\textheight, keepaspectratio]{C:/GitHub/VisualIdeas/data/images/11dc35d015f61f80ec08aec430f6d29a.png}
\end{center}
\vspace{2mm}
\begin{center}
\begin{minipage}{0.92\textwidth}
\small
Arrows are a big part of visual metaphors so I made this easy to follow tutorial on how to make them in Figma (a free design app).
\end{minipage}
\end{center}
\vspace*{\fill}
\newpage
\renewcommand{\theposttitle}{Why friendships are like trees}
\renewcommand{\thepostdate}{2025-02-04}
\vspace*{\fill}
\begin{center}
{\Huge\bfseries ENCHANTING}\\[6mm]
\end{center}
\begin{center}
\includegraphics[width=\textwidth, height=0.55\textheight, keepaspectratio]{C:/GitHub/VisualIdeas/data/images/02d326254577b2b3a1a8a26bece2831a.png}
\end{center}
\vspace{2mm}
\begin{center}
\begin{minipage}{0.92\textwidth}
\small
Darkness cannot drive out darkness; \\
only light can do that. \\
Hate cannot drive out hate; \\
only love can do that.
\end{minipage}
\end{center}
\vspace*{\fill}
\newpage
\renewcommand{\theposttitle}{Why friendships are like trees}
\renewcommand{\thepostdate}{2025-02-04}
\vspace*{\fill}
\begin{center}
{\Huge\bfseries ANALOGY}\\[6mm]
\end{center}
\begin{center}
\includegraphics[width=\textwidth, height=0.55\textheight, keepaspectratio]{C:/GitHub/VisualIdeas/data/images/98df2aa56d3f8b94498c3482c54b2c25.png}
\end{center}
\vspace{2mm}
\begin{center}
\begin{minipage}{0.92\textwidth}
\small
There's beauty in having all 3 friendships in our lives, because they each serve a different purpose. But the most distress I’ve felt in my personal friendships is confusing one for another so it’s helpful to know which friends are which.

Leaves flow in and out of our lives as new leaves fill their place (this is natural so we shouldn’t be too sad when it happens).

Branches—we see daily. Some are stronger than others. They can endure pressure, but only so much before breaking.

Root friends keep us grounded. Often unseen, but we feel their impact and stability.
\end{minipage}
\end{center}
\vspace*{\fill}
\newpage
\renewcommand{\theposttitle}{How money is like water}
\renewcommand{\thepostdate}{2025-02-18}
\vspace*{\fill}
\begin{center}
{\Huge\bfseries INTERESTING}\\[6mm]
\end{center}
\begin{center}
\includegraphics[width=\textwidth, height=0.55\textheight, keepaspectratio]{C:/GitHub/VisualIdeas/data/images/7a84f76ff0797606697143505a271415.png}
\end{center}
\vspace{2mm}
\begin{center}
\begin{minipage}{0.92\textwidth}
\small
I wish I knew this in my 20s:

There'll never be a perfect someday.\\
But there's always an imperfect now.

So how do you start?

With these 3 simple words: \textbf{“I am doing.”}

Not “I will…” (procrastinating)

Not “I should’ve…” (regret)

But \textbf{“I am…”} (right here, right now)
\end{minipage}
\end{center}
\vspace*{\fill}
\newpage
\renewcommand{\theposttitle}{How money is like water}
\renewcommand{\thepostdate}{2025-02-18}
\vspace*{\fill}
\begin{center}
{\Huge\bfseries DESIGN}\\[6mm]
\end{center}
\begin{center}
\includegraphics[width=\textwidth, height=0.55\textheight, keepaspectratio]{C:/GitHub/VisualIdeas/data/images/fef1b00f2c0f51e9489e3b4cd2fb007c.png}
\end{center}
\vspace{2mm}
\begin{center}
\begin{minipage}{0.92\textwidth}
\small
Readers often ask how do you come up with visual metaphors, and this short video is as high-level as you can get with your approach in how to think about it.

And if you’re curious to dive deeper, join the waitlist here for Cohort 10 of Thinking in Visual Metaphors 😊
\end{minipage}
\end{center}
\vspace*{\fill}
\newpage
\renewcommand{\theposttitle}{How money is like water}
\renewcommand{\thepostdate}{2025-02-18}
\vspace*{\fill}
\begin{center}
{\Huge\bfseries ENCHANTING}\\[6mm]
\end{center}
\begin{center}
\includegraphics[width=\textwidth, height=0.55\textheight, keepaspectratio]{C:/GitHub/VisualIdeas/data/images/75b3cfaec48ffafee713f331d39ce7f9.png}
\end{center}
\vspace{2mm}
\begin{center}
\begin{minipage}{0.92\textwidth}
\small
\textbf{Failure doesn't block the way. }\\
\textbf{It creates the path.}
\end{minipage}
\end{center}
\vspace*{\fill}
\newpage
\renewcommand{\theposttitle}{How money is like water}
\renewcommand{\thepostdate}{2025-02-18}
\vspace*{\fill}
\begin{center}
{\Huge\bfseries ANALOGY}\\[6mm]
\end{center}
\begin{center}
\includegraphics[width=\textwidth, height=0.24\textheight, keepaspectratio]{C:/GitHub/VisualIdeas/data/images/3bbe5f2db63cd88de2db77695b545f27.png}
\end{center}
\vspace{2mm}
\begin{center}
\begin{minipage}{0.92\textwidth}
\small
This is the biggest regret most "successful people" have.

The reason is simple:

Money is a means to an end (not the end itself).

But often in the pursuit of more, we sacrifice the very things money was supposed to help solve (or at least make easier).

When the bucket is full, nourish your flowers.

Inspired by Sahil Bloom's book, \textbf{"The 5 Types of Wealth”}
\end{minipage}
\end{center}
\vspace{5mm}
\begin{center}
\includegraphics[width=\textwidth, height=0.24\textheight, keepaspectratio]{C:/GitHub/VisualIdeas/data/images/35b8bd78f61272ff40df6dfe9596afc2.jpg}
\end{center}
\vspace{2mm}
\begin{center}
\begin{minipage}{0.92\textwidth}
\small
People really seemed to like this behind-the-scenes first draft on LinkedIn so sharing it here 😊
\end{minipage}
\end{center}
\vspace*{\fill}
\newpage
\renewcommand{\theposttitle}{How risks are like being stranded on a boat}
\renewcommand{\thepostdate}{2025-03-04}
\vspace*{\fill}
\begin{center}
{\Huge\bfseries INTERESTING}\\[6mm]
\end{center}
\begin{center}
\includegraphics[width=\textwidth, height=0.55\textheight, keepaspectratio]{C:/GitHub/VisualIdeas/data/images/f8f9c706d866a024de7c3835900403c0.png}
\end{center}
\vspace{2mm}
\begin{center}
\begin{minipage}{0.92\textwidth}
\small
Change is risky, but so is standing still.

Most of my career, my whole identity was “Mr. Milani”—\\
A high school film \& animation teacher.

Now I love being a teacher. But whenever students asked, “What are you making?” I froze. Teaching creativity, yet creating nothing of my own, left me feeling like an imposter.

And then the pandemic hit—and confronted me with a scary question: If I can’t teach, then who am I?

One night as I journaled about my need to start creating, I had déjà-vu. So I flipped through old notebooks and found the same words—written 11 years ago!

“I need to at least start.” (\textit{oof})

It seems when you stand in a leaky boat long enough, you stop noticing the water.

So doubts and all, I took a leap—and, yes, I stumbled right into the sharks. 😂\\
(No one gets it right the first try)

The lesson: when you procrastinate a decision, you’re still making a choice to stay where you are.

If you’ve been wanting to make a change, maybe it’s time to test the waters.
\end{minipage}
\end{center}
\vspace*{\fill}
\newpage
\renewcommand{\theposttitle}{How risks are like being stranded on a boat}
\renewcommand{\thepostdate}{2025-03-04}
\vspace*{\fill}
\begin{center}
{\Huge\bfseries DESIGN}\\[6mm]
\end{center}
\begin{center}
\includegraphics[width=\textwidth, height=0.55\textheight, keepaspectratio]{C:/GitHub/VisualIdeas/data/images/30ef4489e43e23f20f928aeac4749334.png}
\end{center}
\vspace{2mm}
\begin{center}
\begin{minipage}{0.92\textwidth}
\small
↳ The simplest hack to overcome “disappointments of outcome” when you post online.
\end{minipage}
\end{center}
\vspace*{\fill}
\newpage
\renewcommand{\theposttitle}{How risks are like being stranded on a boat}
\renewcommand{\thepostdate}{2025-03-04}
\vspace*{\fill}
\begin{center}
{\Huge\bfseries ENCHANTING}\\[6mm]
\end{center}
\begin{center}
\includegraphics[width=\textwidth, height=0.55\textheight, keepaspectratio]{C:/GitHub/VisualIdeas/data/images/69482c3f6fc9958f4b2d0ce90322dfc8.png}
\end{center}
\vspace{2mm}
\begin{center}
\begin{minipage}{0.92\textwidth}
\small
Friendly reminder: You can't fill every cup (and that's perfectly okay).

Work demands your focus

Your partner needs quality 1-on-1 time

Family deserves your presence

Friends want you to show up when it matters

The more cups you try to fill:\\
↳ Faster yours runs empty\\
↳ Longer it takes to fill back up.

You can’t give 100\% to everyone, every day. There's always a trade-off.

You can have it all. Just not all at once.
\end{minipage}
\end{center}
\vspace*{\fill}
\newpage
\renewcommand{\theposttitle}{How risks are like being stranded on a boat}
\renewcommand{\thepostdate}{2025-03-04}
\vspace*{\fill}
\begin{center}
{\Huge\bfseries ANALOGY}\\[6mm]
\end{center}
\begin{center}
\includegraphics[width=\textwidth, height=0.55\textheight, keepaspectratio]{C:/GitHub/VisualIdeas/data/images/a0a40d41332959b7258eca96816e0300.png}
\end{center}
\vspace{2mm}
\begin{center}
\begin{minipage}{0.92\textwidth}
\small
\textbf{Hard work won’t guarantee victory, but it WILL guarantee improvement.}

Read the full post.
\end{minipage}
\end{center}
\vspace*{\fill}
\newpage
\renewcommand{\theposttitle}{How to look at your fears differently}
\renewcommand{\thepostdate}{2025-03-19}
\vspace*{\fill}
\begin{center}
{\Huge\bfseries INTERESTING}\\[6mm]
\end{center}
\begin{center}
\includegraphics[width=\textwidth, height=0.55\textheight, keepaspectratio]{C:/GitHub/VisualIdeas/data/images/fbad289beb295f0a312569c0d1373ed9.png}
\end{center}
\vspace{2mm}
\begin{center}
\begin{minipage}{0.92\textwidth}
\small
Fear is a mile wide, an inch deep.

You stand at the edge, convincing yourself the water’s too deep.\\
You take a step, but you don’t sink! Your foot hits solid ground.\\
You question why you doubted yourself for so long.

You change—action always clears the mud in your head.\\
You remember fear grows in hesitation, shrinks in motion.\\
You move forward—realizing the gap was all in your mind.

Inaction magnifies self-doubt → Action shrinks it
\end{minipage}
\end{center}
\vspace*{\fill}
\newpage
\renewcommand{\theposttitle}{How to look at your fears differently}
\renewcommand{\thepostdate}{2025-03-19}
\vspace*{\fill}
\begin{center}
{\Huge\bfseries DESIGN}\\[6mm]
\end{center}
\begin{center}
\begin{minipage}{0.92\textwidth}
\small
The most common reason people don’t create visual metaphors is the limiting belief they “can’t draw.” But visual metaphors isn’t about “making things pretty,” it’s about making things pithy.

Today’s bite-sized lesson is how we can communicate big ideas with simple shapes and arrows!

I’ll keep it short \& sweet with 3 examples:
\end{minipage}
\end{center}
\vspace{4mm}
\begin{center}
\includegraphics[width=\textwidth, height=0.17\textheight, keepaspectratio]{C:/GitHub/VisualIdeas/data/images/efd850f10f1e9dfac287430c9ddb4f04.png}
\end{center}
\vspace{5mm}
\begin{center}
\includegraphics[width=\textwidth, height=0.17\textheight, keepaspectratio]{C:/GitHub/VisualIdeas/data/images/d72867433cfdeda67359d0a996924458.png}
\end{center}
\vspace{5mm}
\begin{center}
\includegraphics[width=\textwidth, height=0.17\textheight, keepaspectratio]{C:/GitHub/VisualIdeas/data/images/ffb3682a924d0ee8c8e80447ee3ec6e5.png}
\end{center}
\vspace{2mm}
\begin{center}
\begin{minipage}{0.92\textwidth}
\small
There’s no doubt people are more attracted to pretty things. But the good news is you don’t have to be Leonardo da Vinci to make images that get readers to think.

When you’re ready to learn more, you can \textbf{sign up for the waitlist for Thinking in Visual Metaphors. }
\end{minipage}
\end{center}
\vspace*{\fill}
\newpage
\renewcommand{\theposttitle}{How to look at your fears differently}
\renewcommand{\thepostdate}{2025-03-19}
\vspace*{\fill}
\begin{center}
{\Huge\bfseries ENCHANTING}\\[6mm]
\end{center}
\begin{center}
\includegraphics[width=\textwidth, height=0.55\textheight, keepaspectratio]{C:/GitHub/VisualIdeas/data/images/34976c0d9821e244186d7ab9f4f0af0d.png}
\end{center}
\vspace{2mm}
\begin{center}
\begin{minipage}{0.92\textwidth}
\small
When you are anxious because of your to-do list, take comfort in your have-done list.

Inspired by \textbf{Kevin Kelly's book, "Excellent Advice for Living"}
\end{minipage}
\end{center}
\vspace*{\fill}
\newpage
\renewcommand{\theposttitle}{How to look at your fears differently}
\renewcommand{\thepostdate}{2025-03-19}
\vspace*{\fill}
\begin{center}
{\Huge\bfseries ANALOGY}\\[6mm]
\end{center}
\begin{center}
\includegraphics[width=\textwidth, height=0.55\textheight, keepaspectratio]{C:/GitHub/VisualIdeas/data/images/84b460b316a5b85e1f79b2003a3cb9c3.png}
\end{center}
\vspace{2mm}
\begin{center}
\begin{minipage}{0.92\textwidth}
\small
We suffer more often in imagination than in reality.
\end{minipage}
\end{center}
\vspace*{\fill}
\newpage
\renewcommand{\theposttitle}{Why failing is not what you think}
\renewcommand{\thepostdate}{2025-04-03}
\vspace*{\fill}
\begin{center}
{\Huge\bfseries INTERESTING}\\[6mm]
\end{center}
\begin{center}
\includegraphics[width=\textwidth, height=0.55\textheight, keepaspectratio]{C:/GitHub/VisualIdeas/data/images/6ae425b8fef49eb0b14dc832e2365d1d.png}
\end{center}
\vspace{2mm}
\begin{center}
\begin{minipage}{0.92\textwidth}
\small
There are two ladders in your career.

But remember, what helped you climb the first ladder may not be useful in the second.
\end{minipage}
\end{center}
\vspace*{\fill}
\newpage
\renewcommand{\theposttitle}{Why failing is not what you think}
\renewcommand{\thepostdate}{2025-04-03}
\vspace*{\fill}
\begin{center}
{\Huge\bfseries DESIGN}\\[6mm]
\end{center}
\begin{center}
\begin{minipage}{0.92\textwidth}
\small
Engaging in creativity is like building a LEGO set without instructions.

With each new piece, the whole thing warps and twists. You have a vague sense of the shape, but there inevitably is a point where key pieces go missing and you have no idea how to move forward.

The lights turn off and you start running into brick walls that come out of nowhere. This dark void is called \textbf{“The Zone of Ambiguity,”} a temporary but necessary space.
\end{minipage}
\end{center}
\vspace{4mm}
\begin{center}
\includegraphics[width=\textwidth, height=0.55\textheight, keepaspectratio]{C:/GitHub/VisualIdeas/data/images/311e497de466902e4e9348ada58d6317.png}
\end{center}
\vspace{2mm}
\begin{center}
\begin{minipage}{0.92\textwidth}
\small
The only difference between someone who considers themselves creative and someone who doesn’t is the confidence that the \textbf{Zone of Ambiguity is a temporary space.}

And you build confidence through seeing you can in fact can find your way through. With each rep, that gap between ambiguity and insight shrinks smaller, but it will always exist.
\end{minipage}
\end{center}
\vspace*{\fill}
\newpage
\renewcommand{\theposttitle}{Why failing is not what you think}
\renewcommand{\thepostdate}{2025-04-03}
\vspace*{\fill}
\begin{center}
{\Huge\bfseries ENCHANTING}\\[6mm]
\end{center}
\begin{center}
\includegraphics[width=\textwidth, height=0.24\textheight, keepaspectratio]{C:/GitHub/VisualIdeas/data/images/5b0457238efda097a778010386e24cf5.png}
\end{center}
\vspace{5mm}
\begin{center}
\includegraphics[width=\textwidth, height=0.24\textheight, keepaspectratio]{C:/GitHub/VisualIdeas/data/images/f9fc9817ea037aa75beaf30b04f56495.png}
\end{center}
\vspace{2mm}
\begin{center}
\begin{minipage}{0.92\textwidth}
\small
I used to have this T-shirt I loved in high school from a brand called "No Fear."

On the back it said, "There's nothing in life more painful than regret."

Everything we want has a sacrifice.

Everything we don't want has a consequence.

But one thing’s for sure—the regret of "what could've been" is often heavier than the discipline it takes to change the present.
\end{minipage}
\end{center}
\vspace*{\fill}
\newpage
\renewcommand{\theposttitle}{Why failing is not what you think}
\renewcommand{\thepostdate}{2025-04-03}
\vspace*{\fill}
\begin{center}
{\Huge\bfseries ANALOGY}\\[6mm]
\end{center}
\begin{center}
\includegraphics[width=\textwidth, height=0.55\textheight, keepaspectratio]{C:/GitHub/VisualIdeas/data/images/9c770471a730d7645834fd63545c5201.png}
\end{center}
\vspace{2mm}
\begin{center}
\begin{minipage}{0.92\textwidth}
\small
To be clear, no one enjoys failing.\\
It sucks to feel like you're not good at something.

But there'll always be learning pains. Accept the miss and adjust your aim.

Don't quit before you've given yourself a chance to feel what it's like to get good.
\end{minipage}
\end{center}
\vspace*{\fill}
\newpage
\renewcommand{\theposttitle}{The Myth of Doing It All}
\renewcommand{\thepostdate}{2025-04-18}
\vspace*{\fill}
\begin{center}
{\Huge\bfseries INTERESTING}\\[6mm]
\end{center}
\begin{center}
\includegraphics[width=\textwidth, height=0.55\textheight, keepaspectratio]{C:/GitHub/VisualIdeas/data/images/9d206d03120e55f40c584d1dc761400d.png}
\end{center}
\vspace{2mm}
\begin{center}
\begin{minipage}{0.92\textwidth}
\small
Let’s be real. It’s HARD to prioritize.

If you’ve followed me for a while, this visual metaphor may feel familiar. That’s because it’s not the first time I’ve articulated this idea. \textbf{(check out this one if you’re curious)}

The thing is most people (myself included) want to do everything all at once. But that’s a surefire recipe for feeling overwhelmed and frustrated at a lack of making real progress on any one thing.

The thing no one tells when you’ve got multiple irons in the fire is you’re probably going to get burned (out).

And as difficult as it is, one thing at a time is a far more useful approach to truly build momentum.
\end{minipage}
\end{center}
\vspace*{\fill}
\newpage
\renewcommand{\theposttitle}{The Myth of Doing It All}
\renewcommand{\thepostdate}{2025-04-18}
\vspace*{\fill}
\begin{center}
{\Huge\bfseries DESIGN}\\[6mm]
\end{center}
\begin{center}
\includegraphics[width=\textwidth, height=0.55\textheight, keepaspectratio]{C:/GitHub/VisualIdeas/data/images/7a63482637c3d35c619dd48f3178769c.png}
\end{center}
\vspace{2mm}
\begin{center}
\begin{minipage}{0.92\textwidth}
\small
A quick argument on the power of clarity and why visuals are so important.

And if you want to be first to learn when registration for the next cohort of Thinking in Visual Metaphors opens,\textbf{ sign up for the waitlist here. }
\end{minipage}
\end{center}
\vspace*{\fill}
\newpage
\renewcommand{\theposttitle}{The Myth of Doing It All}
\renewcommand{\thepostdate}{2025-04-18}
\vspace*{\fill}
\begin{center}
{\Huge\bfseries ANALOGY}\\[6mm]
\end{center}
\begin{center}
\includegraphics[width=\textwidth, height=0.55\textheight, keepaspectratio]{C:/GitHub/VisualIdeas/data/images/76f584c008ce99a4c3c8e21341dde9c1.png}
\end{center}
\vspace{2mm}
\begin{center}
\begin{minipage}{0.92\textwidth}
\small
Join the conversation on LinkedIn

\textbf{People hear your words—but they remember your actions.}

The world is full of talkers, but the loudest statement you can make is in your follow-through.
\end{minipage}
\end{center}
\vspace*{\fill}
\newpage
\renewcommand{\theposttitle}{How Progress Is Like Crossing a River}
\renewcommand{\thepostdate}{2025-04-29}
\vspace*{\fill}
\begin{center}
{\Huge\bfseries INTERESTING}\\[6mm]
\end{center}
\begin{center}
\includegraphics[width=\textwidth, height=0.55\textheight, keepaspectratio]{C:/GitHub/VisualIdeas/data/images/853dd73c23018493a8347834ab4ec27e.png}
\end{center}
\vspace{2mm}
\begin{center}
\begin{minipage}{0.92\textwidth}
\small
Most people let fear stop them before they start.

"What if I fail?"\\
"What if people judge me?"\\
"What if I'm not good enough?"

These questions paralyze us.\\
Try flipping the script.

Replace "what if" with "EVEN IF."\\
"Even if I fail, I'll learn something valuable."\\
"Even if people judge me, I'll grow stronger."\\
"Even if I'm not perfect, I'll be better than yesterday."

The most successful people I know aren't fearless.\\
They simply move forward, "even if" scared.

\textbf{Read the full post here.}
\end{minipage}
\end{center}
\vspace*{\fill}
\newpage
\renewcommand{\theposttitle}{How Progress Is Like Crossing a River}
\renewcommand{\thepostdate}{2025-04-29}
\vspace*{\fill}
\begin{center}
{\Huge\bfseries DESIGN}\\[6mm]
\end{center}
\begin{center}
\begin{minipage}{0.92\textwidth}
\small
A quick and easy visual metaphor anyone can make ↴
\end{minipage}
\end{center}
\vspace{4mm}
\begin{center}
\includegraphics[width=\textwidth, height=0.55\textheight, keepaspectratio]{C:/GitHub/VisualIdeas/data/images/b8098475b98eacf77cfe5cd2e20796e6.jpeg}
\end{center}
\vspace{2mm}
\begin{center}
\begin{minipage}{0.92\textwidth}
\small
A useful way to show the Pareto Principle (80/20 Rule). You can use this for anything where you want to highlight the importance of the smaller element—like the handful of customers who drive most of your revenue.

The big square is the majority (80\%), while that small detached piece represents the minority (20\%).

Know what else can show a majority/minority relationship? That’s right! A pie chart! So you can think of this framework as a variation of the pie chart but with just a different texture.

What I love about this metaphor is its instant clarity. That small piece calls out, “I’m small, but mighty!”

P.S. When you’re ready to learn more frameworks like this, you can \textbf{secure your spot for Thinking in Visual Metaphors. }Enrollment for Cohort 11 launches today!
\end{minipage}
\end{center}
\vspace*{\fill}
\newpage
\renewcommand{\theposttitle}{How Progress Is Like Crossing a River}
\renewcommand{\thepostdate}{2025-04-29}
\vspace*{\fill}
\begin{center}
{\Huge\bfseries ENCHANTING}\\[6mm]
\end{center}
\begin{center}
\includegraphics[width=\textwidth, height=0.55\textheight, keepaspectratio]{C:/GitHub/VisualIdeas/data/images/d815d5a039c3e11b81514146a5180a40.png}
\end{center}
\vspace{2mm}
\begin{center}
\begin{minipage}{0.92\textwidth}
\small
\textbf{Remember: above every cloud, the sun remains.}
\end{minipage}
\end{center}
\vspace*{\fill}
\newpage
\renewcommand{\theposttitle}{How Progress Is Like Crossing a River}
\renewcommand{\thepostdate}{2025-04-29}
\vspace*{\fill}
\begin{center}
{\Huge\bfseries ANALOGY}\\[6mm]
\end{center}
\begin{center}
\includegraphics[width=\textwidth, height=0.55\textheight, keepaspectratio]{C:/GitHub/VisualIdeas/data/images/25b53c7eda98f17cdc8f104d69da7630.png}
\end{center}
\vspace{2mm}
\begin{center}
\begin{minipage}{0.92\textwidth}
\small
\textbf{If it feels too heavy, you’re probably trying to carry it all at once.}
\end{minipage}
\end{center}
\vspace*{\fill}
\newpage
\renewcommand{\theposttitle}{Why low-hanging fruit are dangerous}
\renewcommand{\thepostdate}{2025-05-12}
\vspace*{\fill}
\begin{center}
{\Huge\bfseries DESIGN}\\[6mm]
\end{center}
\begin{center}
\begin{minipage}{0.92\textwidth}
\small
I had an embarrassing near-miss last week that taught me 3 valuable lessons about visual metaphors.
\end{minipage}
\end{center}
\vspace{4mm}
\begin{center}
\includegraphics[width=\textwidth, height=0.13\textheight, keepaspectratio]{C:/GitHub/VisualIdeas/data/images/b0d773baa1da05ead9bc13a317251618.png}
\end{center}
\vspace{5mm}
\begin{center}
\includegraphics[width=\textwidth, height=0.13\textheight, keepaspectratio]{C:/GitHub/VisualIdeas/data/images/9e78457b5f22f61da2ac331368940c00.png}
\end{center}
\vspace{5mm}
\begin{center}
\includegraphics[width=\textwidth, height=0.13\textheight, keepaspectratio]{C:/GitHub/VisualIdeas/data/images/c25b78110c09d1a7ddc327f5ea2a4eae.png}
\end{center}
\vspace{5mm}
\begin{center}
\includegraphics[width=\textwidth, height=0.13\textheight, keepaspectratio]{C:/GitHub/VisualIdeas/data/images/6a9bb06b192b8c364e7bc879b7fb68bf.png}
\end{center}
\vspace{2mm}
\begin{center}
\begin{minipage}{0.92\textwidth}
\small
Version 1
\end{minipage}
\end{center}
\vspace*{\fill}
\newpage
\renewcommand{\theposttitle}{Why low-hanging fruit are dangerous}
\renewcommand{\thepostdate}{2025-05-12}
\vspace*{\fill}
\begin{center}
{\Huge\bfseries ENCHANTING}\\[6mm]
\end{center}
\begin{center}
\includegraphics[width=\textwidth, height=0.55\textheight, keepaspectratio]{C:/GitHub/VisualIdeas/data/images/d21e78993183bbdea6c971654ab56a64.png}
\end{center}
\vspace{2mm}
\begin{center}
\begin{minipage}{0.92\textwidth}
\small
I've been rethinking leadership lately.

Your impact isn't measured by how high you climb.\\
It's counted in how many others climbed BECAUSE of you.

\textbf{"Top-of-the-ladder" folks:}\\
→ Measure success by title and status\\
→ Chase personal achievements\\
→ Stand alone at the summit

\textbf{True leaders:}\\
→ Carve out paths for their teams\\
→ Often work quietly at ground level\\
→ Find delight in other people's wins
\end{minipage}
\end{center}
\vspace*{\fill}
\newpage
\renewcommand{\theposttitle}{Why low-hanging fruit are dangerous}
\renewcommand{\thepostdate}{2025-05-12}
\vspace*{\fill}
\begin{center}
{\Huge\bfseries ANALOGY}\\[6mm]
\end{center}
\begin{center}
\includegraphics[width=\textwidth, height=0.55\textheight, keepaspectratio]{C:/GitHub/VisualIdeas/data/images/291c9a31016921fba39b052ace8aed0b.png}
\end{center}
\vspace{2mm}
\begin{center}
\begin{minipage}{0.92\textwidth}
\small
Day 1 is all about courage.\\
Day 30 is proof of your commitment.

You'll be surprised what you can do in a month if you show up daily.
\end{minipage}
\end{center}
\vspace*{\fill}
\newpage
\renewcommand{\theposttitle}{The time-saving power of one simple word}
\renewcommand{\thepostdate}{2025-05-27}
\vspace*{\fill}
\begin{center}
{\Huge\bfseries INTERESTING}\\[6mm]
\end{center}
\begin{center}
\includegraphics[width=\textwidth, height=0.55\textheight, keepaspectratio]{C:/GitHub/VisualIdeas/data/images/3929fbe3c58b67d704cf1901204a3bd6.png}
\end{center}
\vspace{2mm}
\begin{center}
\begin{minipage}{0.92\textwidth}
\small
Success isn’t about being a genius on your first shot. Just ask the data (which tells a sobering story).

1st time startup founders success rate: 18\%

Those same founders success rate by 10th attempt: 65\%

"Better than yesterday" isn't just motivational fluff. The idea of becoming an overnight success is a fantasy. A harmful one. (and this doesn’t just relate to entrepreneurship) It's how mastery actually works.

Think about it. If you improve just 1\% daily, you'll be 37x better in one year.\\
Not 365\% better. \\
3,700\% better.

You might ask, how do you actually improve 1\% daily?

Answer 3 simple questions:

“What went slightly better today than yesterday?"

“What small improvement can I make tomorrow?"

“Where am I aiming my next arrow?”

It’s not about having perfect aim. \\
It’s not about being the best.\\
It’s just being a little better. Everyday.

P.S. I talk more about the importance of consistency in \textbf{this post. }
\end{minipage}
\end{center}
\vspace*{\fill}
\newpage
\renewcommand{\theposttitle}{The time-saving power of one simple word}
\renewcommand{\thepostdate}{2025-05-27}
\vspace*{\fill}
\begin{center}
{\Huge\bfseries DESIGN}\\[6mm]
\end{center}
\begin{center}
\begin{minipage}{0.92\textwidth}
\small
If you’re a long time reader, you may have noticed I iterated a new banner!
\end{minipage}
\end{center}
\vspace{4mm}
\begin{center}
\includegraphics[width=\textwidth, height=0.17\textheight, keepaspectratio]{C:/GitHub/VisualIdeas/data/images/0a40c1527e22a739c91d866e9203b7f9.jpg}
\end{center}
\vspace{2mm}
\begin{center}
\begin{minipage}{0.92\textwidth}
\small
Iterations!

Settled on version 10 if you're counting! 😅

Here's the funny thing about iteration. You need fresh eyes to see straight.

That's why finding your creative safety net matters. Those trusted few who'll tell you:

"This font feels off"

"Actually, the shadow is too far too the right"

"The old logo had more charm because…"

Without them, you're just rearranging pixels in the dark. Find your people.\\
(shoutout to Hannah Wilson and Michael J Boorman — Hype crew for life!)
\end{minipage}
\end{center}
\vspace{5mm}
\begin{center}
\includegraphics[width=\textwidth, height=0.17\textheight, keepaspectratio]{C:/GitHub/VisualIdeas/data/images/9d15ea599867bb73e2f0b81cc2f3a118.png}
\end{center}
\vspace{2mm}
\begin{center}
\begin{minipage}{0.92\textwidth}
\small
Before
\end{minipage}
\end{center}
\vspace{5mm}
\begin{center}
\includegraphics[width=\textwidth, height=0.17\textheight, keepaspectratio]{C:/GitHub/VisualIdeas/data/images/a4f3ca02a8974c54a3b5ddf349a91e9f.png}
\end{center}
\vspace{2mm}
\begin{center}
\begin{minipage}{0.92\textwidth}
\small
After

Login or Subscribe to participate
\end{minipage}
\end{center}
\vspace*{\fill}
\newpage
\renewcommand{\theposttitle}{The time-saving power of one simple word}
\renewcommand{\thepostdate}{2025-05-27}
\vspace*{\fill}
\begin{center}
{\Huge\bfseries ENCHANTING}\\[6mm]
\end{center}
\begin{center}
\includegraphics[width=\textwidth, height=0.55\textheight, keepaspectratio]{C:/GitHub/VisualIdeas/data/images/cc7318039901da5fddcd7dbc56a012ae.png}
\end{center}
\vspace{2mm}
\begin{center}
\begin{minipage}{0.92\textwidth}
\small
Time doesn’t disappear.\\
It’s given away, one "yes" at a time.

We make time for what we prioritize.\\
And the shape of a priority starts with the gentle art of “no.”
\end{minipage}
\end{center}
\vspace*{\fill}
\newpage
\renewcommand{\theposttitle}{Why friendships get harder as we get older}
\renewcommand{\thepostdate}{2025-06-18}
\vspace*{\fill}
\begin{center}
{\Huge\bfseries INTERESTING}\\[6mm]
\end{center}
\begin{center}
\includegraphics[width=\textwidth, height=0.55\textheight, keepaspectratio]{C:/GitHub/VisualIdeas/data/images/0907b065c4e0efa5e6c457af7e90c957.png}
\end{center}
\vspace{2mm}
\begin{center}
\begin{minipage}{0.92\textwidth}
\small
As kids, friendships were easy.\\
Same classes. Same recess. Same milestones.

As adults, friendships are complicated.\\
Because we’re all running different races at difference paces.

Getting married. Getting divorced. \\
Starting over. Settling down.\\
Having kids. Choosing not to.

I used to think growing apart meant growing away from each other.\\
Now I understand it just means growing in different directions.

At different speeds. Toward different things.

"I'm here" doesn't mean "I'm stuck"\\
"You're there" doesn't mean "You've left me"\\
It just means we're all finding our way.

But somehow...still sending memes at 2am like nothing's changed. 😄
\end{minipage}
\end{center}
\vspace*{\fill}
\newpage
\renewcommand{\theposttitle}{Why friendships get harder as we get older}
\renewcommand{\thepostdate}{2025-06-18}
\vspace*{\fill}
\begin{center}
{\Huge\bfseries DESIGN}\\[6mm]
\end{center}
\begin{center}
\begin{minipage}{0.92\textwidth}
\small
In this short clip below, I share why that frustrating "stuck" feeling might be a good thing.
\end{minipage}
\end{center}
\vspace{4mm}
\begin{center}
\includegraphics[width=\textwidth, height=0.55\textheight, keepaspectratio]{C:/GitHub/VisualIdeas/data/images/3522947c471ef9b06fba05ca3253895d.png}
\end{center}
\vspace*{\fill}
\newpage
\renewcommand{\theposttitle}{Why friendships get harder as we get older}
\renewcommand{\thepostdate}{2025-06-18}
\vspace*{\fill}
\begin{center}
{\Huge\bfseries ENCHANTING}\\[6mm]
\end{center}
\begin{center}
\includegraphics[width=\textwidth, height=0.55\textheight, keepaspectratio]{C:/GitHub/VisualIdeas/data/images/d81e749facf9ea71bc98f7b7c407b6f0.png}
\end{center}
\vspace{2mm}
\begin{center}
\begin{minipage}{0.92\textwidth}
\small
The thing about circles is...You don't always know you're in one.

You row hard, make waves.\\
Staying in motion, but repeating the same patterns.

Testing if you’re in a circle is simple:\\
When did you last feel uncomfortable?

Because comfort and growth rarely share the same boat.\\
Maybe today's the day to point somewhere new.

Even if you have to row alone for a bit.
\end{minipage}
\end{center}
\vspace*{\fill}
\newpage
\renewcommand{\theposttitle}{Why friendships get harder as we get older}
\renewcommand{\thepostdate}{2025-06-18}
\vspace*{\fill}
\begin{center}
{\Huge\bfseries ANALOGY}\\[6mm]
\end{center}
\begin{center}
\includegraphics[width=\textwidth, height=0.55\textheight, keepaspectratio]{C:/GitHub/VisualIdeas/data/images/f7e50a1f99eb902ab8a02406bd25423a.png}
\end{center}
\vspace{2mm}
\begin{center}
\begin{minipage}{0.92\textwidth}
\small
A good book is like a blueprint. But if all you do is collect blueprints, you're just rearranging furniture in a house you haven't built yet.

The thing is no one ever learned to ride a bike by reading a book. \\
You have to put your foot on the pedal and push forward.

Read the full post here.
\end{minipage}
\end{center}
\vspace*{\fill}
\newpage
\renewcommand{\theposttitle}{Why I'm not worried about AI Images}
\renewcommand{\thepostdate}{2025-07-15}
\vspace*{\fill}
\begin{center}
{\Huge\bfseries INTERESTING}\\[6mm]
\end{center}
\begin{center}
\includegraphics[width=\textwidth, height=0.55\textheight, keepaspectratio]{C:/GitHub/VisualIdeas/data/images/59fcc612d3e4146e996d051822e185d1.png}
\end{center}
\vspace{2mm}
\begin{center}
\begin{minipage}{0.92\textwidth}
\small
When you use other people's definition of success to measure your happiness, consider this:

You wouldn't use other people's house key to open your front door.\\
So why use their scorecard to measure your life?

Their success might be a corner office.Yours might be Tuesday afternoons free.\\
Their win might be scaling to seven figures.\\
Yours might be to make it to all your son's soccer games.

Both are valid.

But happiness has never been a one-size fits all game. It's a personal fit.
\end{minipage}
\end{center}
\vspace*{\fill}
\newpage
\renewcommand{\theposttitle}{Why I'm not worried about AI Images}
\renewcommand{\thepostdate}{2025-07-15}
\vspace*{\fill}
\begin{center}
{\Huge\bfseries DESIGN}\\[6mm]
\end{center}
\begin{center}
\includegraphics[width=\textwidth, height=0.55\textheight, keepaspectratio]{C:/GitHub/VisualIdeas/data/images/4ccee77f323fb924fac379cce9d706c3.png}
\end{center}
\vspace{2mm}
\begin{center}
\begin{minipage}{0.92\textwidth}
\small
If ALL you use is AI images to communicate your ideas, you are slowly blending in. There’s a place for it in your process, but how do you differentiate yourself from everyone else using AI?

My hypothesis: People will get bored of AI images (if not already). And that means there’s still opportunity. Remember: how you draw something is unique to you, imperfections and all. And it’s those very human imperfections you might fear that make your work interesting.

If you choose to lean into it, it may turn into your competitive advantage to stand out over the AI slop that is shaping the creator landscape today.
\end{minipage}
\end{center}
\vspace*{\fill}
\newpage
\renewcommand{\theposttitle}{Why I'm not worried about AI Images}
\renewcommand{\thepostdate}{2025-07-15}
\vspace*{\fill}
\begin{center}
{\Huge\bfseries ENCHANTING}\\[6mm]
\end{center}
\begin{center}
\includegraphics[width=\textwidth, height=0.55\textheight, keepaspectratio]{C:/GitHub/VisualIdeas/data/images/7034d7f3e6fc9a0f1f35c121a3ff3121.png}
\end{center}
\vspace{2mm}
\begin{center}
\begin{minipage}{0.92\textwidth}
\small
We suffer more from imagination than from reality.
\end{minipage}
\end{center}
\vspace*{\fill}
\newpage
\renewcommand{\theposttitle}{Why I'm not worried about AI Images}
\renewcommand{\thepostdate}{2025-07-15}
\vspace*{\fill}
\begin{center}
{\Huge\bfseries ANALOGY}\\[6mm]
\end{center}
\begin{center}
\includegraphics[width=\textwidth, height=0.55\textheight, keepaspectratio]{C:/GitHub/VisualIdeas/data/images/ed2dc2c0dc98135f61d724d9ab8c5a65.png}
\end{center}
\vspace{2mm}
\begin{center}
\begin{minipage}{0.92\textwidth}
\small
You're not doing anything wrong because your path is messy and loopy. \\
Because here's the thing, human beings don't grow in straight lines.

We grow in spirals.

We may circle back to places we've been before, but we're different each time. Stronger. Wiser. Better equipped.

And still moving forward.
\end{minipage}
\end{center}
\vspace*{\fill}
\newpage
\renewcommand{\theposttitle}{What is the "Big Picture Paradox?"}
\renewcommand{\thepostdate}{2025-07-29}
\vspace*{\fill}
\begin{center}
{\Huge\bfseries INTERESTING}\\[6mm]
\end{center}
\begin{center}
\includegraphics[width=\textwidth, height=0.55\textheight, keepaspectratio]{C:/GitHub/VisualIdeas/data/images/a4f8d40133c3ac324b20c15f78570b18.png}
\end{center}
\vspace{2mm}
\begin{center}
\begin{minipage}{0.92\textwidth}
\small
\textbf{Defined by helplessness or defined by determination.}

We cannot change the cards we are dealt, just how we play the hand.
\end{minipage}
\end{center}
\vspace*{\fill}
\newpage
\renewcommand{\theposttitle}{What is the "Big Picture Paradox?"}
\renewcommand{\thepostdate}{2025-07-29}
\vspace*{\fill}
\begin{center}
{\Huge\bfseries DESIGN}\\[6mm]
\end{center}
\begin{center}
\includegraphics[width=\textwidth, height=0.55\textheight, keepaspectratio]{C:/GitHub/VisualIdeas/data/images/f8deeea73e18b0907a1505dd45e7691c.png}
\end{center}
\vspace{2mm}
\begin{center}
\begin{minipage}{0.92\textwidth}
\small
I was thrilled to join my good buddy, Corey Wilks, Psy.D., on his podcast, Creator Alchemy a few weeks ago. We jump straight into this deep dive on:

creative philosophy

the psychology behind consistent output

how to overcame decades of creative paralysis

my approach to handling imposter syndrome

the constraints that fuel creativity

why "good enough" might be the secret to unlocking your creative potential.

You can also check out this awesome summary article Corey wrote with key takeaways (and visual metaphor examples!)
\end{minipage}
\end{center}
\vspace*{\fill}
\newpage
\renewcommand{\theposttitle}{What is the "Big Picture Paradox?"}
\renewcommand{\thepostdate}{2025-07-29}
\vspace*{\fill}
\begin{center}
{\Huge\bfseries ENCHANTING}\\[6mm]
\end{center}
\begin{center}
\includegraphics[width=\textwidth, height=0.55\textheight, keepaspectratio]{C:/GitHub/VisualIdeas/data/images/faec8edaa5fff5877a898f0582c8a101.png}
\end{center}
\vspace{2mm}
\begin{center}
\begin{minipage}{0.92\textwidth}
\small
The bigger your vision, the smaller your focus needs to be.

It's a bit counterintuitive, but those who achieve the impossible are not staring at the top of the mountain.

They're focusing on the next foothold.
\end{minipage}
\end{center}
\vspace*{\fill}
\newpage
\renewcommand{\theposttitle}{What is the "Big Picture Paradox?"}
\renewcommand{\thepostdate}{2025-07-29}
\vspace*{\fill}
\begin{center}
{\Huge\bfseries ANALOGY}\\[6mm]
\end{center}
\begin{center}
\includegraphics[width=\textwidth, height=0.55\textheight, keepaspectratio]{C:/GitHub/VisualIdeas/data/images/058bfbf7e9e5881b248da20e4f201fe3.png}
\end{center}
\vspace{2mm}
\begin{center}
\begin{minipage}{0.92\textwidth}
\small
\textbf{The grass is greener where you water it.}
\end{minipage}
\end{center}
\vspace*{\fill}
\newpage
\renewcommand{\theposttitle}{Why showing up doesn't mean showing up perfectly.}
\renewcommand{\thepostdate}{2025-10-21}
\vspace*{\fill}
\begin{center}
{\Huge\bfseries INTERESTING}\\[6mm]
\end{center}
\begin{center}
\includegraphics[width=\textwidth, height=0.55\textheight, keepaspectratio]{C:/GitHub/VisualIdeas/data/images/5a8f83625005abd39c142ad9d7581ab4.png}
\end{center}
\vspace{2mm}
\begin{center}
\begin{minipage}{0.92\textwidth}
\small
If you’ve been subscribed for a while, you may have noticed I’ve not sent out a newsletter in over a month. That miss has been gnawing away at me daily. A lot of it has come from setting priorities around health and family (see my last email). But if I’m honest, more of it comes from breaking the streak built over several years.

What’s funny when you aim for a goal is that the misses often teach far more than the bullseye. And my recent misses on the newsletter have retaught me the importance of consistency and more importantly momentum.

Showing up doesn’t mean showing up perfectly.

Some days will splash. Some days will drip.\\
Both water the garden.

What garden are you watering as we head into the holidays? Hit reply. I’d love to know!
\end{minipage}
\end{center}
\vspace*{\fill}
\newpage
\renewcommand{\theposttitle}{Why showing up doesn't mean showing up perfectly.}
\renewcommand{\thepostdate}{2025-10-21}
\vspace*{\fill}
\begin{center}
{\Huge\bfseries ENCHANTING}\\[6mm]
\end{center}
\begin{center}
\includegraphics[width=\textwidth, height=0.55\textheight, keepaspectratio]{C:/GitHub/VisualIdeas/data/images/b236ce30aabe9a016fc69acfb2d14ec0.png}
\end{center}
\vspace{2mm}
\begin{center}
\begin{minipage}{0.92\textwidth}
\small
The magic isn't in starting big. It's in starting.

If only I learned in my 20s that your entrance doesn't need to match your ambition. And if there’s no door? Build that first.

Dreams don't care about your starting point.
\end{minipage}
\end{center}
\vspace*{\fill}
\newpage
\renewcommand{\theposttitle}{Why showing up doesn't mean showing up perfectly.}
\renewcommand{\thepostdate}{2025-10-21}
\vspace*{\fill}
\begin{center}
{\Huge\bfseries ANALOGY}\\[6mm]
\end{center}
\begin{center}
\includegraphics[width=\textwidth, height=0.55\textheight, keepaspectratio]{C:/GitHub/VisualIdeas/data/images/241addce9d5448fdab8b8b675bd45d96.png}
\end{center}
\vspace{2mm}
\begin{center}
\begin{minipage}{0.92\textwidth}
\small
You can't sit by a fire you didn't build.

Funny how this works everywhere:

→ Sore muscles before strength\\
→ Messy journals before clear thoughts\\
→ Years of "no" before overnight success

The warmth you want lives inside the work you're avoiding. Let's get after it!
\end{minipage}
\end{center}
\vspace*{\fill}
\newpage
\renewcommand{\theposttitle}{There's a difference between what's cool and what's meaningful}
\renewcommand{\thepostdate}{2025-11-19}
\vspace*{\fill}
\begin{center}
{\Huge\bfseries INTERESTING}\\[6mm]
\end{center}
\begin{center}
\includegraphics[width=\textwidth, height=0.55\textheight, keepaspectratio]{C:/GitHub/VisualIdeas/data/images/ffc47dc13b8d21cc3bdd498cc4deca5d.png}
\end{center}
\vspace{2mm}
\begin{center}
\begin{minipage}{0.92\textwidth}
\small
\textbf{Friendly reminder: }The year is not over yet. Let's get after it. 💪
\end{minipage}
\end{center}
\vspace*{\fill}
\newpage
\renewcommand{\theposttitle}{There's a difference between what's cool and what's meaningful}
\renewcommand{\thepostdate}{2025-11-19}
\vspace*{\fill}
\begin{center}
{\Huge\bfseries ENCHANTING}\\[6mm]
\end{center}
\begin{center}
\includegraphics[width=\textwidth, height=0.55\textheight, keepaspectratio]{C:/GitHub/VisualIdeas/data/images/77f4c4e5ba3f20cf77ec363b86f1c0a8.png}
\end{center}
\vspace{2mm}
\begin{center}
\begin{minipage}{0.92\textwidth}
\small
We become what we refuse to release.

I had a friend who used to collect scars like they were stamps. As if he was curating a museum of hurts and betrayals.

Don’t get me wrong. He was justified.

But over time, it distorted into a weird paradox.\\
Holding the pain validated his story.\\
But the pain then became the story.

The thing is pain can be a teacher. But it’s meant to visit, not move in. And you can acknowledge the hurt without making it your home.

It’s possible the things we can’t release have already released us. We're just the last to know.
\end{minipage}
\end{center}
\vspace*{\fill}
\newpage
\renewcommand{\theposttitle}{There's a difference between what's cool and what's meaningful}
\renewcommand{\thepostdate}{2025-11-19}
\vspace*{\fill}
\begin{center}
{\Huge\bfseries ANALOGY}\\[6mm]
\end{center}
\begin{center}
\includegraphics[width=\textwidth, height=0.55\textheight, keepaspectratio]{C:/GitHub/VisualIdeas/data/images/4c742ffb515c128db013743aed22b111.png}
\end{center}
\vspace{2mm}
\begin{center}
\begin{minipage}{0.92\textwidth}
\small
Teaching others teaches yourself.
\end{minipage}
\end{center}
\vspace*{\fill}
\newpage
\renewcommand{\theposttitle}{Are you preparing or are you procrastinating?}
\renewcommand{\thepostdate}{2025-12-02}
\vspace*{\fill}
\begin{center}
{\Huge\bfseries INTERESTING}\\[6mm]
\end{center}
\begin{center}
\includegraphics[width=\textwidth, height=0.55\textheight, keepaspectratio]{C:/GitHub/VisualIdeas/data/images/2711741791958e42365f2d38026269e3.png}
\end{center}
\vspace{2mm}
\begin{center}
\begin{minipage}{0.92\textwidth}
\small
Wisdom isn't about having all the answers.\\
It's about gathering all the angles.

I illustrated this visual to remind us that your perspective matters. \\
But it's never the complete story.
\end{minipage}
\end{center}
\vspace*{\fill}
\newpage
\renewcommand{\theposttitle}{Are you preparing or are you procrastinating?}
\renewcommand{\thepostdate}{2025-12-02}
\vspace*{\fill}
\begin{center}
{\Huge\bfseries ENCHANTING}\\[6mm]
\end{center}
\begin{center}
\includegraphics[width=\textwidth, height=0.55\textheight, keepaspectratio]{C:/GitHub/VisualIdeas/data/images/5431ee9fbea77dc4a542e3a10ce66182.png}
\end{center}
\vspace{2mm}
\begin{center}
\begin{minipage}{0.92\textwidth}
\small
Not what you think about. Not what you plan to do. What you actually practice daily.

It becomes non-negotiable when it's tied to who you are.

\textbf{Writers write }(even when it's hard).

\textbf{Athletes train} (even when they're tired).

\textbf{Leaders lead} (even when they doubt).

We act in alignment with who we believe we are.\\
\textbf{Belief → Practice → Identity}
\end{minipage}
\end{center}
\vspace*{\fill}
\newpage
\renewcommand{\theposttitle}{Are you preparing or are you procrastinating?}
\renewcommand{\thepostdate}{2025-12-02}
\vspace*{\fill}
\begin{center}
{\Huge\bfseries ANALOGY}\\[6mm]
\end{center}
\begin{center}
\includegraphics[width=\textwidth, height=0.55\textheight, keepaspectratio]{C:/GitHub/VisualIdeas/data/images/9268573813fbc0360e983058ed92bfc9.png}
\end{center}
\vspace{2mm}
\begin{center}
\begin{minipage}{0.92\textwidth}
\small
When I turned 40, I read over two dozen books on writing and creativity.\\
Guess how many words I wrote that same year. Zero.

Books are maps, not the journey.\\
Starting points, not destinations.\\
Don't build a library where you should build experience.

It wasn’t until I committed to post online daily that I built any forward momentum.

Read one good book. Apply something fully. Then maybe read another.

\textbf{Productive:} "Let me understand the basics to start”\\
\textbf{Productive procrastination: }"Let me understand everything"

At some point, you have to close the book to open the door.
\end{minipage}
\end{center}
\vspace*{\fill}
\newpage
\renewcommand{\theposttitle}{What’s the difference between generosity and strategy?}
\renewcommand{\thepostdate}{2025-12-16}
\vspace*{\fill}
\begin{center}
{\Huge\bfseries INTERESTING}\\[6mm]
\end{center}
\begin{center}
\includegraphics[width=\textwidth, height=0.55\textheight, keepaspectratio]{C:/GitHub/VisualIdeas/data/images/bde5ff3ab4b31b99fa416e12b8dcf556.png}
\end{center}
\vspace{2mm}
\begin{center}
\begin{minipage}{0.92\textwidth}
\small
What’s the difference between generosity and strategy?

Or more precisely, what connects your actions and your expectations? (whether it’s spoken aloud or a silent understanding)

“I’ll help you ONLY WHEN you help me" → Transactional

"I'll help you BECAUSE you helped me" → Reciprocal

"I'll help you SO THAT you'll help me later" → Manipulative

"I'll help you." → Unconditional

Some relationships are meant to be transactional. That's okay. It’s how most of the world works. In fact, many of our closest and longstanding relationships fall into the reciprocal (there’s nothing wrong with that).

But the relationships that hit different are the ones that don’t keep score. \\
And it's worth knowing which kind you're in and which kind you're offering.

So here's my question for you this week: When was a time you wanted to be unconditional but then found yourself transition into one of these other modes? (or vice versa)

Reply and let me know. While I can't answer every email, I read every single response!
\end{minipage}
\end{center}
\vspace*{\fill}
\newpage
\renewcommand{\theposttitle}{What’s the difference between generosity and strategy?}
\renewcommand{\thepostdate}{2025-12-16}
\vspace*{\fill}
\begin{center}
{\Huge\bfseries DESIGN}\\[6mm]
\end{center}
\begin{center}
\begin{minipage}{0.92\textwidth}
\small
The best visual metaphors reveal an insight that \textbf{feels new, but still oddly familiar. }They remind us of something we quietly know in our mind, heart, or soul, but due to the frantic hustle \& bustle of the day to day, we forget!
\end{minipage}
\end{center}
\vspace{4mm}
\begin{center}
\includegraphics[width=\textwidth, height=0.24\textheight, keepaspectratio]{C:/GitHub/VisualIdeas/data/images/45f3c04dbcf151c6bdd7a225e2fa9cc0.png}
\end{center}
\vspace{2mm}
\begin{center}
\begin{minipage}{0.92\textwidth}
\small
The key is using what’s familiar to act as the relational bridge between what we know and what feels foreign. Usually the bridge connects (or corrects) a behavior or belief.
\end{minipage}
\end{center}
\vspace{5mm}
\begin{center}
\includegraphics[width=\textwidth, height=0.24\textheight, keepaspectratio]{C:/GitHub/VisualIdeas/data/images/8e9040af525f88582daecec111a347a2.png}
\end{center}
\vspace{2mm}
\begin{center}
\begin{minipage}{0.92\textwidth}
\small
For example, if you’re struggling to find pockets of happiness in your life and I just said “fix your environment and you might be a little happier,” it doesn’t feel very insightful.

So instead, I start with an assumption of what you might familiar (and it’s helpful to point out no assumption holds true for everyone). But what if I assumed you might know what it feels like to chase a butterfly?
\end{minipage}
\end{center}
\vspace*{\fill}
\newpage
\renewcommand{\theposttitle}{What’s the difference between generosity and strategy?}
\renewcommand{\thepostdate}{2025-12-16}
\vspace*{\fill}
\begin{center}
{\Huge\bfseries ENCHANTING}\\[6mm]
\end{center}
\begin{center}
\includegraphics[width=\textwidth, height=0.55\textheight, keepaspectratio]{C:/GitHub/VisualIdeas/data/images/047fd6cae26c23a618a71f8655269153.png}
\end{center}
\vspace*{\fill}
\newpage
\renewcommand{\theposttitle}{What’s the difference between generosity and strategy?}
\renewcommand{\thepostdate}{2025-12-16}
\vspace*{\fill}
\begin{center}
{\Huge\bfseries ANALOGY}\\[6mm]
\end{center}
\begin{center}
\includegraphics[width=\textwidth, height=0.55\textheight, keepaspectratio]{C:/GitHub/VisualIdeas/data/images/24a49f03ec3f80760a43b3eab8507af6.png}
\end{center}
\vspace*{\fill}
\newpage
\renewcommand{\theposttitle}{How to use evidence to build confidence}
\renewcommand{\thepostdate}{2025-12-30}
\vspace*{\fill}
\begin{center}
{\Huge\bfseries INTERESTING}\\[6mm]
\end{center}
\begin{center}
\includegraphics[width=\textwidth, height=0.55\textheight, keepaspectratio]{C:/GitHub/VisualIdeas/data/images/e83c09ee515f52266cffce3137af6ee0.png}
\end{center}
\vspace{2mm}
\begin{center}
\begin{minipage}{0.92\textwidth}
\small
For years, runners believed the human body couldn't run a mile in under four minutes. Doctors said it was physically impossible!

Then on May 6, 1954, British runner Roger Bannister ran it in 3:59.4.

But here’s the crazy part: 46 days later, John Landy from Australia did it in 3:58.0. By year's end, three more runners beat the “impossible” 4-minute mile.

Nothing changed physically, except belief.

I love this story because whenever someone does the "impossible,” others see themselves in that someone.

“I can't" → "they did"

"They did" → "I might"

"I might" → "I will"

It's like a \textbf{possibility chain of evidence}. Of course, your "four-minute mile" might look different. Maybe it's writing that post and hitting publish. Or speaking up in meetings. Or just facing the day.

And it works both ways. You might need someone else's proof to believe in what you can do. Someone else needs your proof to believe in what they can do.

Either way, you're part of this chain.

Break your barriers. Not just for you. \\
For everyone watching who thinks they can't.

(if you like this idea, check out my visual for The Stockdale Paradox)
\end{minipage}
\end{center}
\vspace*{\fill}
\newpage
\renewcommand{\theposttitle}{How to use evidence to build confidence}
\renewcommand{\thepostdate}{2025-12-30}
\vspace*{\fill}
\begin{center}
{\Huge\bfseries DESIGN}\\[6mm]
\end{center}
\begin{center}
\begin{minipage}{0.92\textwidth}
\small
When you brainstorm visual ideas, it’s okay to be fast, messy, and even incomplete! The most important thing → get it down on paper and out of your head. I call this the \textbf{“minimum viable capture.”}

My notebooks are full of chicken scratches often filled more with words than sketches. They look like this:
\end{minipage}
\end{center}
\vspace{4mm}
\begin{center}
\includegraphics[width=\textwidth, height=0.13\textheight, keepaspectratio]{C:/GitHub/VisualIdeas/data/images/bc8ec2dd10fe93f5eb178785ca1ac417.jpg}
\end{center}
\vspace{2mm}
\begin{center}
\begin{minipage}{0.92\textwidth}
\small
Complaining visual (brainstorm)
\end{minipage}
\end{center}
\vspace{5mm}
\begin{center}
\includegraphics[width=\textwidth, height=0.13\textheight, keepaspectratio]{C:/GitHub/VisualIdeas/data/images/4a65e3f32afcc0d8a5d3e310515875a4.jpg}
\end{center}
\vspace{2mm}
\begin{center}
\begin{minipage}{0.92\textwidth}
\small
Stockdale Paradox (brainstorm)

And they eventually turn into this:
\end{minipage}
\end{center}
\vspace{5mm}
\begin{center}
\includegraphics[width=\textwidth, height=0.13\textheight, keepaspectratio]{C:/GitHub/VisualIdeas/data/images/e72e62ddd60f5ef7257bbafb15907a28.png}
\end{center}
\vspace{5mm}
\begin{center}
\includegraphics[width=\textwidth, height=0.13\textheight, keepaspectratio]{C:/GitHub/VisualIdeas/data/images/0910457e1f1c2e832c6e35068e4d82ed.png}
\end{center}
\vspace{2mm}
\begin{center}
\begin{minipage}{0.92\textwidth}
\small
The point is to avoid being precious with your sketches (especially if it’s preventing you from working out the ideas in your head).

By the way, if you’re ready to take the first step to translate your ideas to visuals, Thinking in Visual Metaphors starts in a couple weeks! Enrollment closes soon!
\end{minipage}
\end{center}
\vspace*{\fill}
\newpage
\renewcommand{\theposttitle}{How to use evidence to build confidence}
\renewcommand{\thepostdate}{2025-12-30}
\vspace*{\fill}
\begin{center}
{\Huge\bfseries ENCHANTING}\\[6mm]
\end{center}
\begin{center}
\includegraphics[width=\textwidth, height=0.24\textheight, keepaspectratio]{C:/GitHub/VisualIdeas/data/images/3a403bc6019dd17e4813c115f7060e8d.png}
\end{center}
\vspace{2mm}
\begin{center}
\begin{minipage}{0.92\textwidth}
\small
The more efficient you get, the more you can get done, right?

But does “success” just boil down to getting the most done and as fast as possible? Hmm. Depends how you define success.

Personally, I’m burning out trying to be more efficient. So as we gear up for 2026, I’m trying something new → Slooowing dooown. I call it \textbf{“bullet-time productivity”} 😬
\end{minipage}
\end{center}
\vspace{5mm}
\begin{center}
%\includegraphics[width=\textwidth, height=0.24\textheight, keepaspectratio]{C:/GitHub/VisualIdeas/data/images/f39a268d25cc999780cd93fe634a8728.gif}
\end{center}
\vspace{2mm}
\begin{center}
\begin{minipage}{0.92\textwidth}
\small
Instead of doing more, I’m \textbf{noticing more}, \textbf{feeling more}, \textbf{appreciating more}. And the number one thing I’m \textbf{loving more} is really seeing the people who’re next to me on the journey.

The older I get, the more I just want to spend every minute around the people I love.

Wishing you and yours safety and health, wondrous journeys, and well-earned time with the people who make you happiest this New Years! ❤️
\end{minipage}
\end{center}
\vspace*{\fill}
\newpage
\renewcommand{\theposttitle}{How to use evidence to build confidence}
\renewcommand{\thepostdate}{2025-12-30}
\vspace*{\fill}
\begin{center}
{\Huge\bfseries ANALOGY}\\[6mm]
\end{center}
\begin{center}
\includegraphics[width=\textwidth, height=0.55\textheight, keepaspectratio]{C:/GitHub/VisualIdeas/data/images/d21eb22e1b9a8db5803b7d8734388728.png}
\end{center}
\vspace*{\fill}
\newpage
\renewcommand{\theposttitle}{The bumps you’re feeling doesn’t mean you’re on the wrong road}
\renewcommand{\thepostdate}{2026-01-13}
\vspace*{\fill}
\begin{center}
{\Huge\bfseries INTERESTING}\\[6mm]
\end{center}
\begin{center}
\includegraphics[width=\textwidth, height=0.55\textheight, keepaspectratio]{C:/GitHub/VisualIdeas/data/images/1c3f24277f09c12de45dcd741d0c4983.png}
\end{center}
\vspace*{\fill}
\newpage
\renewcommand{\theposttitle}{The bumps you’re feeling doesn’t mean you’re on the wrong road}
\renewcommand{\thepostdate}{2026-01-13}
\vspace*{\fill}
\begin{center}
{\Huge\bfseries DESIGN}\\[6mm]
\end{center}
\begin{center}
\begin{minipage}{0.92\textwidth}
\small
Today’s lesson at first blush feels may feel like it has nothing to do with design. But your creative habits have everything to do with what you actually make.
\end{minipage}
\end{center}
\vspace{4mm}
\begin{center}
\includegraphics[width=\textwidth, height=0.55\textheight, keepaspectratio]{C:/GitHub/VisualIdeas/data/images/a0f0ad6a595bada5af1e30ed2192f12e.png}
\end{center}
\vspace{2mm}
\begin{center}
\begin{minipage}{0.92\textwidth}
\small
The morning is when I feel most creative so for me, the number one creative habit I’m working on is waking up early. And I. Am. STRUGGLING! 🫠

But with time, it will get easier. I’m not giving up!

So here’s my question this week: \textbf{What one habit (if you could could get into it) would help you be more creative?}

Reply and let me know. While I can't write back to everyone, I read every response!
\end{minipage}
\end{center}
\vspace*{\fill}
\newpage
\renewcommand{\theposttitle}{The bumps you’re feeling doesn’t mean you’re on the wrong road}
\renewcommand{\thepostdate}{2026-01-13}
\vspace*{\fill}
\begin{center}
{\Huge\bfseries ENCHANTING}\\[6mm]
\end{center}
\begin{center}
\includegraphics[width=\textwidth, height=0.55\textheight, keepaspectratio]{C:/GitHub/VisualIdeas/data/images/b624ac186d6b2c721c0aad2839a04d36.png}
\end{center}
\vspace{2mm}
\begin{center}
\begin{minipage}{0.92\textwidth}
\small
A lesson I have to learn over and over again:

Sometimes the bumps you’re feeling doesn’t mean you’re on the wrong road. \\
It’s just that the path to get to where you want to go is a bumpy ride!
\end{minipage}
\end{center}
\vspace*{\fill}
\newpage
\renewcommand{\theposttitle}{The bumps you’re feeling doesn’t mean you’re on the wrong road}
\renewcommand{\thepostdate}{2026-01-13}
\vspace*{\fill}
\begin{center}
{\Huge\bfseries ANALOGY}\\[6mm]
\end{center}
\begin{center}
\includegraphics[width=\textwidth, height=0.55\textheight, keepaspectratio]{C:/GitHub/VisualIdeas/data/images/6e31475a54eb130650ce3a71d89fa176.png}
\end{center}
\vspace*{\fill}
\newpage
\renewcommand{\theposttitle}{You're not the problem. The problem is the problem.}
\renewcommand{\thepostdate}{2026-01-28}
\vspace*{\fill}
\begin{center}
{\Huge\bfseries INTERESTING}\\[6mm]
\end{center}
\begin{center}
\includegraphics[width=\textwidth, height=0.55\textheight, keepaspectratio]{C:/GitHub/VisualIdeas/data/images/9b62519e7c092248527622bef1bcc233.png}
\end{center}
\vspace*{\fill}
\newpage
\renewcommand{\theposttitle}{You're not the problem. The problem is the problem.}
\renewcommand{\thepostdate}{2026-01-28}
\vspace*{\fill}
\begin{center}
{\Huge\bfseries ENCHANTING}\\[6mm]
\end{center}
\begin{center}
\includegraphics[width=\textwidth, height=0.55\textheight, keepaspectratio]{C:/GitHub/VisualIdeas/data/images/928fe1332ba23bd31f0e21b3fc55dc31.png}
\end{center}
\vspace{2mm}
\begin{center}
\begin{minipage}{0.92\textwidth}
\small
\textbf{Painful how often I feel like the poor inflatable stickman here. }🫠
\end{minipage}
\end{center}
\vspace*{\fill}
\newpage
\renewcommand{\theposttitle}{Why you should always "root" for yourself}
\renewcommand{\thepostdate}{2026-02-10}
\vspace*{\fill}
\begin{center}
{\Huge\bfseries INTERESTING}\\[6mm]
\end{center}
\begin{center}
\includegraphics[width=\textwidth, height=0.55\textheight, keepaspectratio]{C:/GitHub/VisualIdeas/data/images/6e1af08dadef73ecf733b5aac609c806.png}
\end{center}
\vspace*{\fill}
\newpage
\renewcommand{\theposttitle}{Why you should always "root" for yourself}
\renewcommand{\thepostdate}{2026-02-10}
\vspace*{\fill}
\begin{center}
{\Huge\bfseries DESIGN}\\[6mm]
\end{center}
\begin{center}
\includegraphics[width=\textwidth, height=0.13\textheight, keepaspectratio]{C:/GitHub/VisualIdeas/data/images/c5eb36cc238bc769c8705a2214504767.png}
\end{center}
\vspace{2mm}
\begin{center}
\begin{minipage}{0.92\textwidth}
\small
A lot of people are surprised when I tell them \textbf{Twitter (not Instagram) is the first place I started posting visuals.} It’s also my “least successful” (I cringe at defining success by follower count…but let’s temporarily use this definition).

Over the last 4 years of regularly posting on X, my follower count grew to about 19,000. Nothing to complain about but half that growth happened in the first 6 months. It’s felt like a grind to grow the other half.

Then \textbf{last week something happened. }The algorithm gods in the matrix saw it fit to bring some attention to my X account probably during a perfect storm of the right person seeing the right content at the right time of day.
\end{minipage}
\end{center}
\vspace{5mm}
\begin{center}
\includegraphics[width=\textwidth, height=0.13\textheight, keepaspectratio]{C:/GitHub/VisualIdeas/data/images/c11e59208b12e6569b247bc9707c36ec.png}
\end{center}
\vspace{2mm}
\begin{center}
\begin{minipage}{0.92\textwidth}
\small
Dec. 6, 2025 - 4 Likes
\end{minipage}
\end{center}
\vspace{5mm}
\begin{center}
\includegraphics[width=\textwidth, height=0.13\textheight, keepaspectratio]{C:/GitHub/VisualIdeas/data/images/a9547ed9a14ce5bc8235432fdd458fc7.png}
\end{center}
\vspace{2mm}
\begin{center}
\begin{minipage}{0.92\textwidth}
\small
Feb. 5, 2026 - 16K Likes

Thankfully I learned early as a creator that social media algorithms are fickle creatures (notice that my content didn’t change, the algo did). It’s a dangerous game to measure success by their metrics. So I’ve stumbled into playing a more infinite game. The game of showing up despite the outcome.

\textbf{The 3 most important lessons I’ve learned in this game:}

You will go on a winning streak

You will go on a losing streak

You won’t know when either will stop.

And this has happened to me in every platform.
\end{minipage}
\end{center}
\vspace{5mm}
\begin{center}
\includegraphics[width=\textwidth, height=0.13\textheight, keepaspectratio]{C:/GitHub/VisualIdeas/data/images/76871455636823b53a6cfc770dfcb118.png}
\end{center}
\vspace*{\fill}
\newpage
\renewcommand{\theposttitle}{Why you should always "root" for yourself}
\renewcommand{\thepostdate}{2026-02-10}
\vspace*{\fill}
\begin{center}
{\Huge\bfseries ENCHANTING}\\[6mm]
\end{center}
\begin{center}
\includegraphics[width=\textwidth, height=0.55\textheight, keepaspectratio]{C:/GitHub/VisualIdeas/data/images/1806bf19516de7db8d1557df857bb67d.png}
\end{center}
\vspace{2mm}
\begin{center}
\begin{minipage}{0.92\textwidth}
\small
That thing that made you weird as a kid could you make you great as an adult.
\end{minipage}
\end{center}
\vspace*{\fill}
\newpage
\renewcommand{\theposttitle}{Why you should always "root" for yourself}
\renewcommand{\thepostdate}{2026-02-10}
\vspace*{\fill}
\begin{center}
{\Huge\bfseries ANALOGY}\\[6mm]
\end{center}
\begin{center}
\includegraphics[width=\textwidth, height=0.55\textheight, keepaspectratio]{C:/GitHub/VisualIdeas/data/images/a412140411b867baef8afa393f6ad360.png}
\end{center}
\vspace*{\fill}
\newpage
\renewcommand{\theposttitle}{How purpose is like a candle}
\renewcommand{\thepostdate}{2026-02-24}
\vspace*{\fill}
\begin{center}
{\Huge\bfseries INTERESTING}\\[6mm]
\end{center}
\begin{center}
\includegraphics[width=\textwidth, height=0.55\textheight, keepaspectratio]{C:/GitHub/VisualIdeas/data/images/a881a196d34bb56314ac541a95086982.png}
\end{center}
\vspace{2mm}
\begin{center}
\begin{minipage}{0.92\textwidth}
\small
Motivation is what gets you started.\\
Habit is what keeps you going.
\end{minipage}
\end{center}
\vspace*{\fill}
\newpage
\renewcommand{\theposttitle}{How purpose is like a candle}
\renewcommand{\thepostdate}{2026-02-24}
\vspace*{\fill}
\begin{center}
{\Huge\bfseries DESIGN}\\[6mm]
\end{center}
\begin{center}
\includegraphics[width=\textwidth, height=0.17\textheight, keepaspectratio]{C:/GitHub/VisualIdeas/data/images/206a6bcd11f59d99cae9e2be54ffe47b.png}
\end{center}
\vspace{2mm}
\begin{center}
\begin{minipage}{0.92\textwidth}
\small
The biggest reason most folks believe they can’t make visual metaphors may come to you as an obvious one—\textbf{”I can’t draw.”}*

But I’ll scream this till my face is as blue as my brand colors: visual metaphors are not about pretty pictures, they’re about thinking. My friend Laura Evans said it best: “It’s not art. It’s marks with meaning.“

And the meaning doesn’t come from the graphic. Sure the pictures are a vehicle. But the passengers they deliver are the message behind the pictures.

The truth is \textbf{if you can write, you’re 70\% of the way} to crafting your own visual metaphors.

Okay. Okay. You’re right. They are called VISUAL metaphors. But the above visual is an example of how the writing is MORE important than the graphic.

Check out my Instagram comparison of the same graphic, different words:
\end{minipage}
\end{center}
\vspace{5mm}
\begin{center}
\includegraphics[width=\textwidth, height=0.17\textheight, keepaspectratio]{C:/GitHub/VisualIdeas/data/images/360af0cbadb525e1100d06f1d9eff3d5.png}
\end{center}
\vspace{2mm}
\begin{center}
\begin{minipage}{0.92\textwidth}
\small
Feb. 21, 2025 — 1,525 Likes
\end{minipage}
\end{center}
\vspace{5mm}
\begin{center}
\includegraphics[width=\textwidth, height=0.17\textheight, keepaspectratio]{C:/GitHub/VisualIdeas/data/images/66ce058427b0717fdd36093cf2fb5397.png}
\end{center}
\vspace{2mm}
\begin{center}
\begin{minipage}{0.92\textwidth}
\small
Feb. 21, 2026 - 24K+ Likes

It’s true you have to understand some basic visual strategies to translate a metaphor into a visual metaphor, but if you know how to wordsmith (and if you understand that creativity is an iterative pursuit), you will stand in the top 1\% of visual creators!

\textbf{*side note:} is anyone else pissed off that A.I. has killed the human use of em-dashes?!
\end{minipage}
\end{center}
\vspace*{\fill}
\newpage
\renewcommand{\theposttitle}{How purpose is like a candle}
\renewcommand{\thepostdate}{2026-02-24}
\vspace*{\fill}
\begin{center}
{\Huge\bfseries ENCHANTING}\\[6mm]
\end{center}
\begin{center}
\includegraphics[width=\textwidth, height=0.55\textheight, keepaspectratio]{C:/GitHub/VisualIdeas/data/images/577d3e87f84997e691ef251f8cb14ec7.png}
\end{center}
\vspace{2mm}
\begin{center}
\begin{minipage}{0.92\textwidth}
\small
→ \textbf{Leaves} offer color \& change.\\
→ \textbf{Branches} offer support.\\
→ \textbf{Roots} become a part of who you are.
\end{minipage}
\end{center}
\vspace*{\fill}
\newpage
\renewcommand{\theposttitle}{How purpose is like a candle}
\renewcommand{\thepostdate}{2026-02-24}
\vspace*{\fill}
\begin{center}
{\Huge\bfseries ANALOGY}\\[6mm]
\end{center}
\begin{center}
\includegraphics[width=\textwidth, height=0.55\textheight, keepaspectratio]{C:/GitHub/VisualIdeas/data/images/110f5079ea39828ca1bc9b43746fd9de.png}
\end{center}
\vspace*{\fill}
\newpage
\renewcommand{\theposttitle}{The myth of "being ready"}
\renewcommand{\thepostdate}{2026-03-18}
\vspace*{\fill}
\begin{center}
{\Huge\bfseries INTERESTING}\\[6mm]
\end{center}
\begin{center}
\includegraphics[width=\textwidth, height=0.55\textheight, keepaspectratio]{C:/GitHub/VisualIdeas/data/images/be7e66611183a6ed070257eac4e3196c.png}
\end{center}
\vspace{2mm}
\begin{center}
\begin{minipage}{0.92\textwidth}
\small
Sure we’re already halfway through March, but 2026 is still full of opportunities. Let's get after it 💪
\end{minipage}
\end{center}
\vspace*{\fill}
\newpage
\renewcommand{\theposttitle}{The myth of "being ready"}
\renewcommand{\thepostdate}{2026-03-18}
\vspace*{\fill}
\begin{center}
{\Huge\bfseries DESIGN}\\[6mm]
\end{center}
\begin{center}
\includegraphics[width=\textwidth, height=0.13\textheight, keepaspectratio]{C:/GitHub/VisualIdeas/data/images/a56b1486419543093ba667336813a1d9.png}
\end{center}
\vspace{2mm}
\begin{center}
\begin{minipage}{0.92\textwidth}
\small
Last week, thanks to supporters like you, I had the privilege of running a 90min workshop at SXSW EDU called “Teaching with Visual Thinking” in Austin, TX!
\end{minipage}
\end{center}
\vspace{5mm}
\begin{center}
\includegraphics[width=\textwidth, height=0.13\textheight, keepaspectratio]{C:/GitHub/VisualIdeas/data/images/988d7c27e022d76e882a95e27c12797d.jpg}
\end{center}
\vspace{2mm}
\begin{center}
\begin{minipage}{0.92\textwidth}
\small
Presenting at SXSW EDU
\end{minipage}
\end{center}
\vspace{5mm}
\begin{center}
\includegraphics[width=\textwidth, height=0.13\textheight, keepaspectratio]{C:/GitHub/VisualIdeas/data/images/74159f0066d521522b5f328f105fa361.jpg}
\end{center}
\vspace{2mm}
\begin{center}
\begin{minipage}{0.92\textwidth}
\small
The energy in the room was electric!

Speaking to a group of educators was super meaningful. Attendees came up throughout the rest of the conference to say things like this:
\end{minipage}
\end{center}
\vspace{5mm}
\begin{center}
\includegraphics[width=\textwidth, height=0.13\textheight, keepaspectratio]{C:/GitHub/VisualIdeas/data/images/3f3e060b39839d8ad5411ea5c19fd9c0.jpg}
\end{center}
\vspace{2mm}
\begin{center}
\begin{minipage}{0.92\textwidth}
\small
But here's the thing. I almost didn't go for it.

When I first decided to speak at SXSW, I had a choice:

Option 1: \textbf{20min presentation} (done it before, better odds of getting in)

Option 2: \textbf{90min workshop} (4X the prep! Harder to get in)

The cohort course I teach is 3 weeks long. Choosing option 2 would mean distilling 9-hours of lessons into 90 minutes! I wasn't ready for that.

Which was exactly why I knew I had to do it.

Truth is I've spent a large part of my life waiting. Thinking if I prepped enough, I'd finally feel "ready."

It wasn't until my 40s that I finally realized the problem with thinking your way to "ready" is you don't know what you don't know. No amount of prep changes that.

Every success I've had over the last 4 years of my creator journey has come from flipping the script: \textbf{Do first, ready last.}

So learn as you go. Just in time. It's scary, sloppy, and uncomfortable. And you have to be willing to look foolish while you figure it out. But it’s far better than the alternative:

Spending your whole life in the waiting room.
\end{minipage}
\end{center}
\vspace*{\fill}
\newpage
\renewcommand{\theposttitle}{The myth of "being ready"}
\renewcommand{\thepostdate}{2026-03-18}
\vspace*{\fill}
\begin{center}
{\Huge\bfseries ENCHANTING}\\[6mm]
\end{center}
\begin{center}
\includegraphics[width=\textwidth, height=0.55\textheight, keepaspectratio]{C:/GitHub/VisualIdeas/data/images/9707c6525c1bd249a9cd6a918038f1b0.png}
\end{center}
\vspace{2mm}
\begin{center}
\begin{minipage}{0.92\textwidth}
\small
Three words that change everything:\\
\textbf{"I am learning."}

It's not toxic positivity.\\
It's not pretending things are easy.\\
It's just...grace.

The same grace you'd give a friend who's trying something new.
\end{minipage}
\end{center}
\vspace*{\fill}
\newpage
\renewcommand{\theposttitle}{The myth of "being ready"}
\renewcommand{\thepostdate}{2026-03-18}
\vspace*{\fill}
\begin{center}
{\Huge\bfseries ANALOGY}\\[6mm]
\end{center}
\begin{center}
\includegraphics[width=\textwidth, height=0.55\textheight, keepaspectratio]{C:/GitHub/VisualIdeas/data/images/0ef221e87d442ec14cc74f0c349e8624.png}
\end{center}
\vspace{2mm}
\begin{center}
\begin{minipage}{0.92\textwidth}
\small
\textbf{Pull the weeds.}\\
\textbf{Feed the roots.}\\
\textbf{Reap the fruits.}
\end{minipage}
\end{center}
\vspace*{\fill}
\newpage
\renewcommand{\theposttitle}{How emotions are like threads}
\renewcommand{\thepostdate}{2026-04-01}
\vspace*{\fill}
\begin{center}
{\Huge\bfseries INTERESTING}\\[6mm]
\end{center}
\begin{center}
\includegraphics[width=\textwidth, height=0.55\textheight, keepaspectratio]{C:/GitHub/VisualIdeas/data/images/a239411ade075a042d60ee64b2840ff9.png}
\end{center}
\vspace*{\fill}
\newpage
\renewcommand{\theposttitle}{How emotions are like threads}
\renewcommand{\thepostdate}{2026-04-01}
\vspace*{\fill}
\begin{center}
{\Huge\bfseries DESIGN}\\[6mm]
\end{center}
\begin{center}
\includegraphics[width=\textwidth, height=0.24\textheight, keepaspectratio]{C:/GitHub/VisualIdeas/data/images/e94376b7c6316ef234743ec7dbc8988a.png}
\end{center}
\vspace{2mm}
\begin{center}
\begin{minipage}{0.92\textwidth}
\small
A reader asked the other week, \textbf{“How do you show adaptation?”}

One strategy is a simple line graph called \textbf{damped oscillation}. After a shock, the line swings back and forth, each arc smaller than the last, until it eventually settles in the middle.

It’s a great way to show resilience in the face of adversity (and one you can whip up on a whiteboard quickly whenever you need). For example:
\end{minipage}
\end{center}
\vspace{5mm}
\begin{center}
\includegraphics[width=\textwidth, height=0.24\textheight, keepaspectratio]{C:/GitHub/VisualIdeas/data/images/8d16d9ccd4c67d5996d4001453f22ed9.png}
\end{center}
\vspace{2mm}
\begin{center}
\begin{minipage}{0.92\textwidth}
\small
Example of Damped Oscillation at work

The key idea is wild swings at the beginning aren't a sign that something is wrong. It’s just part of a process.

By the way if you like stuff like this, Cohort 16 of Thinking in Visual Metaphors starts in late June. Sign up for the waitlist here.
\end{minipage}
\end{center}
\vspace*{\fill}
\newpage
\renewcommand{\theposttitle}{How emotions are like threads}
\renewcommand{\thepostdate}{2026-04-01}
\vspace*{\fill}
\begin{center}
{\Huge\bfseries ENCHANTING}\\[6mm]
\end{center}
\begin{center}
\includegraphics[width=\textwidth, height=0.55\textheight, keepaspectratio]{C:/GitHub/VisualIdeas/data/images/8c77158eca38571c8eb5b15b665d058d.png}
\end{center}
\vspace*{\fill}
\newpage
\renewcommand{\theposttitle}{How emotions are like threads}
\renewcommand{\thepostdate}{2026-04-01}
\vspace*{\fill}
\begin{center}
{\Huge\bfseries ANALOGY}\\[6mm]
\end{center}
\begin{center}
\includegraphics[width=\textwidth, height=0.55\textheight, keepaspectratio]{C:/GitHub/VisualIdeas/data/images/941d2ce77ca08c422e674d03443733e4.png}
\end{center}
\vspace{2mm}
\begin{center}
\begin{minipage}{0.92\textwidth}
\small
We don't "move on" from grief, \\
we move forward with it.
\end{minipage}
\end{center}
\vspace*{\fill}
\newpage
\renewcommand{\theposttitle}{The power of mistakes and discipline}
\renewcommand{\thepostdate}{2026-04-15}
\vspace*{\fill}
\begin{center}
{\Huge\bfseries INTERESTING}\\[6mm]
\end{center}
\begin{center}
\includegraphics[width=\textwidth, height=0.55\textheight, keepaspectratio]{C:/GitHub/VisualIdeas/data/images/8d719fbe84441e7d9bbc5cb038c7c191.png}
\end{center}
\vspace*{\fill}
\newpage
\renewcommand{\theposttitle}{The power of mistakes and discipline}
\renewcommand{\thepostdate}{2026-04-15}
\vspace*{\fill}
\begin{center}
{\Huge\bfseries DESIGN}\\[6mm]
\end{center}
\begin{center}
\begin{minipage}{0.92\textwidth}
\small
The gold standard of showing the power of compounding: \\
\textbf{the exponential curve.}

It’s a line that starts almost flat, growing so slowly it feels like nothing’s happening at all (this is where most people quit). But then, the sweet, sweet fruit of all that effort pays off in an “everything, everywhere, all at once” moment that pushes your results to take off!
\end{minipage}
\end{center}
\vspace{4mm}
\begin{center}
\includegraphics[width=\textwidth, height=0.55\textheight, keepaspectratio]{C:/GitHub/VisualIdeas/data/images/e8afd1bcc58679befd1678ab97648ea7.png}
\end{center}
\vspace{2mm}
\begin{center}
\begin{minipage}{0.92\textwidth}
\small
One of the most famous examples of this curve comes from James Clear’s \textbf{“Atomic Habits”.}
\end{minipage}
\end{center}
\vspace*{\fill}
\newpage
\renewcommand{\theposttitle}{The power of mistakes and discipline}
\renewcommand{\thepostdate}{2026-04-15}
\vspace*{\fill}
\begin{center}
{\Huge\bfseries ENCHANTING}\\[6mm]
\end{center}
\begin{center}
\includegraphics[width=\textwidth, height=0.55\textheight, keepaspectratio]{C:/GitHub/VisualIdeas/data/images/56f17e095666fc73782ca7c7751b06f3.png}
\end{center}
\vspace*{\fill}
\newpage
\renewcommand{\theposttitle}{The power of mistakes and discipline}
\renewcommand{\thepostdate}{2026-04-15}
\vspace*{\fill}
\begin{center}
{\Huge\bfseries ANALOGY}\\[6mm]
\end{center}
\begin{center}
\includegraphics[width=\textwidth, height=0.55\textheight, keepaspectratio]{C:/GitHub/VisualIdeas/data/images/366138628c6e1fe03d16f46e8462031c.png}
\end{center}
\vspace{2mm}
\begin{center}
\begin{minipage}{0.92\textwidth}
\small
I've never regretted pausing. \\
I've regretted plenty of reactions.
\end{minipage}
\end{center}
\vspace*{\fill}
\newpage
\renewcommand{\theposttitle}{The 2 types of talent}
\renewcommand{\thepostdate}{2026-04-29}
\vspace*{\fill}
\begin{center}
{\Huge\bfseries INTERESTING}\\[6mm]
\end{center}
\begin{center}
\includegraphics[width=\textwidth, height=0.55\textheight, keepaspectratio]{C:/GitHub/VisualIdeas/data/images/7107960c2e8198a17e4cc17dbd1fa19b.png}
\end{center}
\vspace{2mm}
\begin{center}
\begin{minipage}{0.92\textwidth}
\small
Natural talent and \textbf{Earned talent}.

Of course when you're born with traits that make you "look like a natural", you have a leg up. But if you stop there, you'll never meet your full potential.

When I was a kid, I never thought of myself as someone who could draw. I just enjoyed doing it and thought it was fun.

And you should see the visual metaphors I first posted 4 years ago! Each one took me an embarrassing amount of time and the result wasn't all that great.

But I kept at it, and after 900+ visuals…I got better! That's earned ability and something anyone can do.

It's time to earn your talent.

\textbf{* Side note:} When I posted this visual, there was meaningful dialogue around whether there is such a thing as “natural talent” in the first place. I have my own thoughts, but curious to hear what you think.
\end{minipage}
\end{center}
\vspace*{\fill}
\newpage
\renewcommand{\theposttitle}{The 2 types of talent}
\renewcommand{\thepostdate}{2026-04-29}
\vspace*{\fill}
\begin{center}
{\Huge\bfseries DESIGN}\\[6mm]
\end{center}
\begin{center}
\begin{minipage}{0.92\textwidth}
\small
The \textbf{J-Curve} (aka “the dip”) visualizes the experience of trying to do something to improve the current situation, BUT with a warning embedded in the curve: \textbf{it gets worse before it gets better. }
\end{minipage}
\end{center}
\vspace{4mm}
\begin{center}
\includegraphics[width=\textwidth, height=0.24\textheight, keepaspectratio]{C:/GitHub/VisualIdeas/data/images/3c882e734f13cf94ac84b6d6ae8320b4.png}
\end{center}
\vspace{2mm}
\begin{center}
\begin{minipage}{0.92\textwidth}
\small
Notice there’s a before state, a moment of change, and an after state.

It’s important the "after" level be much HIGHER than the "before" level, instead of a return to status quo. The dip acts like a rubber band pulled down first so it can fly higher later!

The following is an example inspired by a clip from "The Curiosity Shop" with Adam Grant and Brené Brown that introduces the idea of "adaptive guilt."
\end{minipage}
\end{center}
\vspace{5mm}
\begin{center}
\includegraphics[width=\textwidth, height=0.24\textheight, keepaspectratio]{C:/GitHub/VisualIdeas/data/images/53d7694b6e6f14e48b2fe7d216df1d13.png}
\end{center}
\vspace{2mm}
\begin{center}
\begin{minipage}{0.92\textwidth}
\small
In this example, the J-curve shows after feeling guilty for doing a bad thing, we can self-reflect and grow into a version of ourselves that meets a higher standard of behavior.

By the way if you like stuff like this, Cohort 16 of Thinking in Visual Metaphors starts in late June. Sign up for the waitlist here.
\end{minipage}
\end{center}
\vspace*{\fill}
\newpage
\renewcommand{\theposttitle}{The 2 types of talent}
\renewcommand{\thepostdate}{2026-04-29}
\vspace*{\fill}
\begin{center}
{\Huge\bfseries ENCHANTING}\\[6mm]
\end{center}
\begin{center}
\includegraphics[width=\textwidth, height=0.55\textheight, keepaspectratio]{C:/GitHub/VisualIdeas/data/images/1a7ecad72022ba45b85ae486636a484a.png}
\end{center}
\vspace{2mm}
\begin{center}
\begin{minipage}{0.92\textwidth}
\small
\textbf{Read the full post here. }
\end{minipage}
\end{center}
\vspace*{\fill}
\newpage
\renewcommand{\theposttitle}{The 2 types of talent}
\renewcommand{\thepostdate}{2026-04-29}
\vspace*{\fill}
\begin{center}
{\Huge\bfseries ANALOGY}\\[6mm]
\end{center}
\begin{center}
\includegraphics[width=\textwidth, height=0.55\textheight, keepaspectratio]{C:/GitHub/VisualIdeas/data/images/55e95c38ee4df9b4346b15c280a6efae.png}
\end{center}
\vspace*{\fill}
\end{document}
